% concatenated from main.tex
\documentclass[a4paper,11pt]{amsart}


\newfont{\cyr}{wncyr10 scaled 1100}

\setcounter{tocdepth}{2}

\usepackage[left=2.7cm,right=2.7cm,top=3.5cm,bottom=3cm]{geometry}
\usepackage{amsthm,amssymb,amsmath,amsfonts,mathrsfs,amscd,yfonts}
\usepackage{verbatim}
\usepackage{turnstile}
\usepackage[latin1]{inputenc}
%\usepackage[utf8]{inputenc}
\usepackage[all]{xy}
\usepackage{latexsym}
\usepackage{stmaryrd}
\usepackage{dsfont}
\usepackage{longtable}
\usepackage{multirow}
\usepackage[usenames,dvipsnames]{color}
\usepackage{listings}
\usepackage{mathtools}
\usepackage{tikz}
\usepackage{tikz-cd}
\usetikzlibrary{matrix,arrows,decorations.pathmorphing}
\usepackage{bm}
%\usepackage[notref,notcite]{showkeys}
\usepackage[shortlabels]{enumitem}
\usepackage[pagebackref]{hyperref}

\theoremstyle{plain}
\newtheorem*{thmA}{Theorem~A}
\newtheorem*{thmB}{Theorem~B}
\newtheorem*{thmC}{Theorem~C}
\newtheorem*{thmD}{Theorem~D}
\newtheorem*{hyp}{Hypothesis}
\newtheorem{theorem}{Theorem}[section]
\newtheorem{corollary}[theorem]{Corollary}
\newtheorem{lemma}[theorem]{Lemma}
\newtheorem{proposition}[theorem]{Proposition}
\newtheorem{propo}[theorem]{Proposition}
\newtheorem{coro}[theorem]{Corollary}
\newtheorem{conj}[theorem]{Conjecture}

\theoremstyle{definition}
\newtheorem{definition}[theorem]{Definition}
\newtheorem{defi}[theorem]{Definition}
\newtheorem{question}[theorem]{Question}
\newtheorem{examplewr}[theorem]{Example}
\newtheorem{ass}[theorem]{Assumption}

\theoremstyle{remark}
\newtheorem{obswr}[theorem]{Observation}
\newtheorem{remarkwr}[theorem]{Remark}
\newtheorem{hypotheses}[theorem]{Hypotheses}

\newenvironment{obs}{\begin{obswr}\begin{upshape}}{\end{upshape}\end{obswr}}
\newenvironment{remark}{\begin{remarkwr}\begin{upshape}}{\end{upshape}\end{remarkwr}}
\newenvironment{example}{\begin{examplewr}\begin{upshape}}{\end{upshape}\end{examplewr}}

\newcommand{\Fmar}[1]{\marginpar{\raggedright\tiny \textcolor{green}{#1}}}

\newcommand{\sk}{\vspace{0.1in}}
\newcommand{\any}{?}

%\newcommand{\hV}{{\mathbb{V}}}
\newcommand{\hW}{{\mathbb{W}}}
\newcommand{\hD}{{\mathbb{D}}}
\newcommand{\hU}{{\mathbb{U}}}

\newcommand{\bb}{\mathbb}
\newcommand{\frk}{\mathfrak}
\newcommand{\cl}{\mathcal}
\newcommand{\bd}{\mathscr{O}}
\newcommand{\univ}{\kappa}
\newcommand{\CM}{{\boldsymbol{\theta}_{\psi_0}}}

\newcommand{\kap}{{\kappa_{\psi,g,h}}}
\newcommand{\kapone}{{\kappa_{\psi,g,h,1}}}
\newcommand{\kapinfty}{{\kappa_{\psi,g,h,\infty}}}
\newcommand{\kapinftyone}{{\kappa_{\psi,g,h,1,\infty}}}
\newcommand{\kapinftyn}{{\kappa_{\psi,g,h,n,\infty}}}
\newcommand{\kapinftynq}{{\kappa_{\psi,g,h,nq,\infty}}}
\newcommand{\kapn}{{\kappa_{\psi,g,h,n}}}
\newcommand{\kapnq}{{\kappa_{\psi,g,h,nq}}}
\newcommand{\tkapn}{{\tilde{\kappa}_{\psi,g,h,n}}}
\newcommand{\tkapnq}{{\tilde{\kappa}_{\psi,g,h,nq}}}
\newcommand{\bkap}{{\boldsymbol{\kappa}_{\psi,g,h}}}
\newcommand{\tam}{\mathfrak{c}^{\rm tam}}

\newcommand{\kaps}{{\kappa_{\psi,\mathrm{ad}^0(g)}}}
\newcommand{\kapones}{{\kappa_{\psi,\mathrm{ad}^0(g),1}}}
\newcommand{\kapinftys}{{\kappa_{\psi,\mathrm{ad}^0(g),\infty}}}
\newcommand{\kapinftyones}{{\kappa_{\psi,\mathrm{ad}^0(g),1,\infty}}}
\newcommand{\kapinftyns}{{\kappa_{\psi,\mathrm{ad}^0(g),m,\infty}}}
\newcommand{\kapinftynqs}{{\kappa_{\psi,\mathrm{ad}^0(g),mq,\infty}}}
\newcommand{\kapns}{{\kappa_{\psi,\mathrm{ad}^0(g),n}}}
\newcommand{\tkapns}{{\tilde{\kappa}_{\psi,\mathrm{ad}^0(g),m}}}
\newcommand{\tkapnqs}{{\tilde{\kappa}_{\psi,\mathrm{ad}^0(g),mq}}}
\newcommand{\bkaps}{{\boldsymbol{\kappa}_{\psi,\mathrm{ad}^0(g)}}}

\newcommand{\Vrep}{V^\psi_{g,h}}
\newcommand{\Trep}{T^\psi_{g,h}}
\newcommand{\Arep}{A^\psi_{g,h}}

\newcommand{\Vsrep}{V^\psi_{\mathrm{ad}^0(g)}}
\newcommand{\Tsrep}{T^\psi_{\mathrm{ad}^0(g)}}
\newcommand{\Asrep}{A^\psi_{\mathrm{ad}^0(g)}}

\newcommand{\bT}{\mathbf{T}_{\mathrm{ad}(g)}^\psi}
\newcommand{\bA}{\mathbf{A}_{\mathrm{ad}(g)}^\psi}


\DeclareMathOperator{\dR}{\mathrm{dR}}
\DeclareMathOperator{\sign}{\mathrm{sign}}
\DeclareMathOperator{\sym}{Sym}
\DeclareMathOperator{\alt}{Alt}
\DeclareMathOperator{\wt}{wt}
\DeclareMathOperator{\rk}{rk}
\DeclareMathOperator{\PR}{PR}
\DeclareMathOperator{\Sym}{Sym}
\DeclareMathOperator{\BK}{BK}
\DeclareMathOperator{\BF}{BF}
\DeclareMathOperator{\sgn}{sgn}
\DeclareMathOperator{\cla}{cl}
\DeclareMathOperator{\cris}{cris}
\DeclareMathOperator{\la}{la}
\DeclareMathOperator{\et}{et}
\DeclareMathOperator{\id}{Id}
\DeclareMathOperator{\reg}{reg}
\DeclareMathOperator{\pr}{pr}
\DeclareMathOperator{\cyc}{cyc}
\DeclareMathOperator{\fin}{f}
\DeclareMathOperator{\sing}{sing}
\DeclareMathOperator{\nr}{nr}
\DeclareMathOperator{\nh}{nh}
\DeclareMathOperator{\hol}{hol}
\DeclareMathOperator{\sd}{sd}
\DeclareMathOperator{\spec}{sp}
\DeclareMathOperator{\prim}{prim}
\DeclareMathOperator{\Ext}{Ext}
\DeclareMathOperator{\Kum}{Kum}
\DeclareMathOperator{\can}{can}
\DeclareMathOperator{\rec}{rec}
\DeclareMathOperator{\cusp}{cusp}
\DeclareMathOperator{\Spf}{Spf}
\DeclareMathOperator{\alg}{alg}
\DeclareMathOperator{\grp}{grp}
\DeclareMathOperator{\red}{red}
\DeclareMathOperator{\Reg}{Reg}
\DeclareMathOperator{\Ind}{Ind}
\DeclareMathOperator{\ur}{ur}
\DeclareMathOperator{\Add}{Ad}
\DeclareMathOperator{\lcm}{lcm}
\DeclareMathOperator{\bal}{bal}
\DeclareMathOperator{\Eis}{Eis}
\DeclareMathOperator{\Ad}{Ad}
\DeclareMathOperator{\mot}{mot}
\DeclareMathOperator{\Log}{Log}
\DeclareMathOperator{\oc}{oc}
\DeclareMathOperator{\Res}{Res}
\DeclareMathOperator{\an}{an}
\DeclareMathOperator{\ind}{ind}
\DeclareMathOperator{\triv}{triv}
\DeclareMathOperator{\Image}{Image}
\DeclareMathOperator{\contt}{cont}
\DeclareMathOperator{\ad}{ad}
\DeclareMathOperator{\loc}{loc}
\DeclareMathOperator{\length}{length}
\DeclareMathOperator{\coker}{coker}
\DeclareMathOperator{\Gr}{bal}
\DeclareMathOperator{\str}{str}
\DeclareMathOperator{\rel}{rel}
\DeclareMathOperator{\Iw}{Iw}
\DeclareMathOperator{\im}{im}
\DeclareMathOperator{\Char}{Char}
\DeclareMathOperator{\tors}{tors}
\DeclareMathOperator{\Frac}{Frac}
\DeclareMathOperator{\ac}{ac}
\DeclareMathOperator{\ab}{ab}
\DeclareMathOperator{\cor}{cor}
\DeclareMathOperator{\Idd}{Id}
\DeclareMathOperator{\Tsym}{Tsym}
\DeclareMathOperator{\Symm}{Symm}
\DeclareMathOperator{\Det}{Det}
\DeclareMathOperator{\HS}{HS}
\DeclareMathOperator{\st}{st}
\newcommand{\eReg}{{\widetilde{\Reg}}}


\newcommand{\ukappa}{{\underline{\kappa}}}
\newcommand{\XX}{\mathbb{X}}
\newcommand{\fa}{{\mathfrak{a}}}
\newcommand{\fb}{{\mathfrak{b}}}
\newcommand{\fc}{{\mathfrak{c}}}


\newcommand{\half}{{{\ndtstile{2}{1}}}}
\newcommand{\mhalf}{{{\ndtstile{2}{-1}}}}

\newcommand{\hV}{{\mathbb{V}}}
\newcommand{\Deltafp}{{^\flat\!\Delta'}}
\newcommand{\Deltaf}{{^\flat\!\Delta}}
\newcommand{\Deltap}{{\Delta'}}
\newcommand{\hphi}{\boldsymbol{\phi}}
\newcommand{\hkappa}{\boldsymbol{\kappa}}
\newcommand{\hpartial}{\boldsymbol{\partial}}
\newcommand{\hlambda}{\boldsymbol{\lambda}}
\newcommand{\TT}{{\mathbb T}}
\newcommand{\uchi}{\mbox{\boldmath$\chi$}}

\newcommand{\testf}{\breve{f}}
\newcommand{\testg}{\breve{g}}
\newcommand{\testh}{\breve{h}}
\newcommand{\testhf}{\breve{\hf}}
\newcommand{\testhg}{\breve{\hg}}
\newcommand{\testhh}{\breve{\hh}}
\newcommand{\testG}{\breve{G}}
\newcommand{\testH}{\breve{H}}
\newcommand{\htf}{\breve{\mathbf{f}}}
\newcommand{\htg}{\breve{\mathbf{g}}}
\newcommand{\hth}{\breve{\mathbf{h}}}
\newcommand{\Lp}{{\mathscr{L}_p}}



\newcommand{\uvarepsilon}{\underline{\varvarepsilon}_{\cyc}}
\newcommand{\groupringr}{{\Lambda^{^{{\!\!\!\!\circ}}}_r}}
\newcommand{\groupring}{{\Lambda^{^{\!\!\!\!\circ}}}}
\newcommand{\grtripler}{{{\Lambda^{\!\!\!\circ}_r^{\otimes 3}}}}
\newcommand{\grtriple}{{{\Lambda^{\!\!\!\circ \ \ \ \otimes 3}}}}

\newcommand{\cF}{\mathcal F}
\newcommand{\cH}{\mathcal H}
\newcommand{\co}{o}
\newcommand{\cC}{\mathcal C}
\newcommand{\V}{\mathcal V}
\newcommand{\cE}{\mathcal E}
\newcommand{\cU}{\mathcal U}
\newcommand{\cW}{\mathcal W}
\newcommand{\cA}{\mathcal A}
\newcommand{\g}{\gamma}
\newcommand{\G}{\Gamma}
\newcommand{\Gap}{\Gamma_0^D(p M)}
\newcommand{\Ga}{\Gamma_0^D(M)}
\newcommand{\hGa}{\hat{\Gamma }_0(M)}
\newcommand{\Q}{\mathbb{Q}}
\newcommand{\Z}{\mathbb{Z}}
\newcommand{\F}{\mathbb{F}}
\newcommand{\K}{\mathbb{K}}
\newcommand{\C}{\mathbb{C}}
\newcommand{\I}{\mathbb{I}}
\newcommand{\Ql}{\Q_{\ell}}
\newcommand{\GQ}{G_{\Q}}
\newcommand{\PP}{\mathbb{P}}
\newcommand{\Sel}{\mathrm{Sel}}
\newcommand{\grSel}{\mathfrak{Sel}}
\newcommand{\Sha}{\mathrm{Sha}}
\newcommand{\Gal}{\mathrm{Gal\,}}
\newcommand{\GL}{\mathrm{GL}}
\newcommand{\Cl}{\mathrm{Cl}}
\newcommand{\Tr}{\mathrm{Tr}}
\newcommand{\Div}{\mathrm{Div}}
\newcommand{\Fil}{\mathrm{Fil}}
\newcommand{\rank}{\mathrm{rank}}
\newcommand{\cond}{\mathrm{cond}}
\newcommand{\Frob}{\mathrm{Fr}}
\newcommand{\End}{\mathrm{End}}
\newcommand{\Jac}{\mathrm{Jac}}
\newcommand{\disc}{\mathrm{disc}}
\newcommand{\Id}{\mathrm{Id}}
\newcommand{\MS}{\mathrm{MS}}
\newcommand{\Aut}{\mathrm{Aut}}
\newcommand{\Fr}{\mathrm{Fr}}
\newcommand{\diag}{\mathrm{diag}}
\newcommand{\Up}{\mathrm{U}_p}
\newcommand{\ord}{{\mathrm{ord}}}
\newfont{\gotip}{eufb10 at 12pt}
\newcommand{\gn}{\mbox{\gotip n}}
\newcommand{\gm}{\mbox{\gotip m}}
\newcommand{\ga}{\mbox{\gotip a}}
\newcommand{\gb}{\mbox{\gotip b}}
\newcommand{\gf}{\mbox{\gotip f}}
\newcommand{\go}{{\mathcal O}_{A_f}}
\newcommand{\OAf}{{\mathcal O}_{A_f}}
\newcommand{\gp}{\mbox{\gotip p}}
\newcommand{\gL}{\mbox{\gotip L}}
\newcommand{\gd}{\mbox{\gotip d}}
\newcommand{\cO}{{\mathcal O}}
\newcommand{\cM}{{\mathcal M}}
\newcommand{\cN}{{\mathcal N}}
\newcommand{\cP}{{\mathcal P}}
\newcommand{\cL}{{\mathcal L}}
\newcommand{\cS}{{\mathcal S}}
\newcommand{\cG}{{\mathcal G}}
\newcommand{\GG}{{\mathcal G}}
\newcommand{\cD}{{\mathcal D}}
\newcommand{\om}{{\omega}}
\newcommand{\ra}{\rightarrow}
\newcommand{\lra}{\longrightarrow}
\newcommand{\hra}{\hookrightarrow}
\newcommand{\qbar}{\bar\Q}
\newcommand{\Norm}{\mathrm{N}}
\newcommand{\NS}{\mathrm{NS}}
\newcommand{\SL}{{\mathrm {SL}}}
\newcommand{\SO}{{\mathrm {SO}}}
\newcommand{\n}{{\mathrm{n}}}
\newcommand{\Pic}{{\mathrm{Pic}}}
\newcommand{\CH}{{\mathrm{CH}}}
\newcommand{\AJ}{{\mathrm{AJ}}}
\newcommand{\tr}{{\mathrm{tr}}}
\newcommand{\si}{\sigma}
\newcommand{\eps}{\varepsilon}
\newcommand{\ta}{{\tau}}
\newcommand{\R}{{\mathbb R}}
\newcommand{\M}{{\mathrm{M}}}
\newcommand{\HC}{\mathcal F_{\rm har}}
\newcommand{\PGL}{{\mathrm{PGL}}}
\newcommand{\Supp}{\mathrm{Supp}}
\newcommand{\ep}{\varvarepsilon}
\newcommand{\cX}{\mathcal X}
\newcommand{\hE}{{\mathbf{E}}}
\newcommand{\hf}{{\mathbf{f}}}
\newcommand{\hg}{{\mathbf{g}}}
\newcommand{\hh}{{\mathbf h}}
\newcommand{\hB}{{\mathbf B}}
\newcommand{\VV}{{\mathbf V}}
\newcommand{\DD}{{\mathbf D}}
\newcommand{\homega}{{\mathbf \omega}}
\newcommand{\heta}{{\mathbf \eta}}
\newcommand{\scQ}{{\mathscr{Q}}}
\newcommand{\scN}{{\mathscr{N}}}
\DeclareMathOperator{\cts}{cts}


\newcommand{\SP}{\vspace{0.5cm}}
\newcommand{\T}{\mathbb T}
\newcommand{\new}{\mathrm{new}}
\DeclareMathOperator{\Hom}{Hom} \DeclareMathOperator{\Meas}{Meas}
\newcommand{\X}{\mathbb X}
\newcommand{\Y}{\mathbb Y}
\newcommand{\U}{\mathbb U}
\newcommand{\Rat}{{\rm Rat}\big(\PP^1(\Q_p),\C_p\big)}
\newcommand{\Ker}{{\mathrm Ker}}
\newcommand{\res}{\mathrm{res}}

\newcommand{\Hida}{{\mathcal H}}


\newcommand{\rom}{\mathrm}
\newcommand{\fr}{\mathfrak}
\newcommand{\fp}{{\mathfrak p}}
%\newcommand{\cl }{\mathcal}

\newcommand{\longmono}{\mbox{$\lhook\joinrel\longrightarrow$}}
\newcommand{\longleftmono}{\mbox{$\longleftarrow\joinrel\rhook$}}
\newcommand{\longepi}{\mbox{$\relbar\joinrel\twoheadrightarrow$}}
\newcommand{\mat}[4]{\left(\begin{array}{cc}#1&#2\\#3&#4\end{array}\right)}
\newcommand{\smallmat}[4]{\big(\begin{smallmatrix}#1&#2\\#3&#4\end{smallmatrix}\big)}
\newcommand{\dirlim}{\mathop{\varinjlim}\limits}
\newcommand{\invlim}{\mathop{\varprojlim}\limits}
\newcommand{\dBr}[1]{\llbracket{#1}\rrbracket}


\newcommand{\D}{\mathbb D}

\numberwithin{equation}{section}


\include{thebibliography}

\begin{document}

\title[Anticyclotomic diagonal cycles and Beilinson--Flach]{Anticyclotomic diagonal classes and Beilinson--Flach elements}

\author{Ra\'ul Alonso, Lois Omil-Pazos and \'Oscar Rivero}

\begin{abstract}
We present a comparison between the anticyclotomic Euler system of diagonal cycles associated with the convolution of two modular forms and the cyclotomic Beilinson--Flach Euler system. This extends the seminal work of Bertolini, Darmon, and Venerucci, who established a link between (anticyclotomic) Heegner points and the Beilinson--Kato system. Our approach hinges on a detailed analysis of $p$-adic $L$-functions and Perrin-Riou maps and exploits the Eisenstein degeneration of diagonal cycles along Hida families, working with a CM family which specializes to an irregular Eisenstein series in weight one. We use these results to derive some arithmetic applications.
\end{abstract}


%\date{\today\;{\color{blue} Draft still in progress}}
\date{\today}

\address{R. A.: School of Mathematics and Statistics, University College Dublin, Belfield, Dublin 4, Ireland}
\email{raul.alonsorodriguez@ucd.ie}

\address{L. O.-P.: CITMAga, Santiago de Compostela, Spain}
\email{lois.omil.pazos@usc.es}

\address{O. R.: CITMAga and Departamento de Matem\'aticas, Universidade de Santiago de Compostela, Santiago de Compostela, Spain}
\email{oscar.rivero@usc.es}


\subjclass[2010]{11R23; 11F85}


\maketitle


\setcounter{tocdepth}{1}
\tableofcontents


\section{Introduction}

The recent work of Bertolini, Darmon, and Venerucci \cite{BDV} establishes a conjecture relating Heegner cycles to Beilinson--Kato elements, linking both objects to $p$-adic families of Beilinson--Flach elements in the higher Chow groups of products of two modular curves. The comparison between Heegner points and Beilinson--Kato elements had previously been investigated by B\"uy\"ukkboduk \cite{buyukboduk16} and Venerucci \cite{Ven}, within the framework of the exceptional zero conjectures. In this work, we extend the results of \cite{BDV} by relating Beilinson--Flach classes to the diagonal cycles studied in \cite{BSV} and \cite{DR3}, which in turn can be interpreted in terms of the anticyclotomic classes constructed in \cite{ACR}, \cite{ACR2}, and \cite{castella-do}.

This work may also be interpreted in the context of the recent developments of Loeffler and the third author \cite{LR1}, where they study the Eisenstein degeneration of Euler systems, i.e. Euler systems obtained from critical Eisenstein series. Our setting, however, differs in a significant way, since it involves as input a (non-cohomological) weight-one Eisenstein series. We will return to this point throughout the text.

Following \cite{BeDiPo}, the central feature in the setting of irregular weight-one modular forms attached to Dirichlet characters $(\chi_1,\chi_2)$ with $\chi_1(p)=\chi_2(p)$ is the existence of three (ordinary) families passing through it: the families of Eisenstein series $\hE(\chi_1,\chi_2)$ and $\hE(\chi_2,\chi_1)$, as well as a third Hida family that is generically cuspidal. The key idea, both in this note and in earlier related works, is the introduction of an auxiliary modular form that can be suitably deformed along a Hida or Coleman family to produce the required Galois representation. In this sense, the present work may be viewed as an instance of CM degeneration, while the work of Loeffler and the third author provides an example of Eisenstein degeneration.

\subsection{The set-up}

Let $p$ be an odd prime and fix once and for all an embedding $\iota_p:\overline{\bb{Q}}\xrightarrow{} \overline{\bb{Q}}_p$ and an embedding $\iota: \overline{\bb{Q}}\hookrightarrow \bb{C}$. The latter determines a choice of complex conjugation in $G_{\bb{Q}}$ which will be denoted by $c$.

Let $K/\Q$ be an imaginary quadratic field in which $p$ splits and let $\varepsilon_K$ be the quadratic character attached to it. Let $\hf$ be the CM Hida family introduced in \S\ref{subsec:wt1}, which specializes in weight one to the irregular weight-one Eisenstein series $f = \Eis_1(\varepsilon_K)$. Let $(\hg,\hh)$ be a pair of Hida families of coprime tame levels $(N_g,N_h)$ and characters $(\chi_\hg,\chi_\hh)$ such that $\chi_\hg \chi_\hh =\varepsilon_K \omega^{2r}$ for some $r\in \bb{Z}$. We also assume that $\hg$ and $\hh$ are residually irreducible and $p$-distinguished. Let $g=\hg_{y_0}^\circ$ and $h=\hh_{z_0}^\circ$ be good crystalline specializations of $\hg$ and $\hh$ of weights $l_0\geq 2$ and $m_0\geq 1$, respectively. We assume that $p\nmid \mathrm{cond}(h)$. Set $c_0=(l_0+m_0-1)/2$.

%Write $h_K$ for the class number of $K$. Let $g\in S_l(N_g,\chi_g)$ and $h\in S_m(N_h,\chi_h)$ be newforms of weights $l\geq m\geq 2$ of different parity and nebentypus $\chi_g$ and $\chi_h$, with $\chi_g \chi_h = \varepsilon_K$. Let $k>0$ be an odd integer, and let $\psi$ be a Hecke character of $K$ of infinity type $(1-k,0)$, conductor $\frk{f}$, and central character $\varepsilon_K$. Observe that while in our previous work \cite{ACR} $k$ was even and $l$ and $m$ had the same parity, everything works verbatim just assuming that $k+l+m$ is even. Let $N=\lcm(N_g,N_h)$

In \S\ref{subsec:families} we introduce our notations regarding Galois representations attached to families of modular forms. In particular, attached to $\hg$ (and similarly for $\hf$ and $\hh$) there is a locally-free rank-two module $\mathbb V_{\hg}$, defined over a finite flat extension of $\Lambda = \mathbb Z_p[[1+p\mathbb Z_p]]$ and equipped with a continuous action of the absolute Galois group $G_{\mathbb Q}$.

\subsection{Euler systems and $p$-adic $L$-functions}

We now present the two cohomology class that can naturally be attached to the triple $(f,\hg,\hh)$. Firstly, the diagonal cycle class, which may be understood as an anticyclotomic Euler system as discussed in \cite{ACR}; secondly, the cyclotomic system of Beilinson--Flach, as developed e.g. in \cite{KLZ}.

\begin{enumerate}
    \item[(A)] The work of Bertolini--Seveso--Venerucci \cite{BSV} and Darmon--Rotger \cite{DR3} provides us with a diagonal cycle class \[ \kappa(\hf,\hg,\hh) \in H^1_{\bal}(\mathbb Q,  \bb{V}_{\hf} \hat\otimes \mathbb V_{\hg} \hat \otimes \mathbb V_{\hh}(2-{\bf t})), \] where $2{\bf t} = {\bf k} + {\bf l} + {\bf m}$. This class is constructed via the $p$-adic variation of diagonal cycles, so it may be understood as a geometric object. By specilizing the first variable, it yields a class $\kappa(f,\hg,\hh)\in H^1_{\bal}(\bb{Q},V_f\otimes\bb{V}_{\hg}\hat{\otimes}\bb{V}_{\hh}(2-\mathbf{t}_1))$, where $2\mathbf{t}_1=1+\mathbf{l}+\mathbf{m}$.

    \item[(B)] The work of Kings--Loeffler--Zerbes \cite{KLZ} provides us with a (cyclotomic) Beilinson--Flach class
    \[
    \kappa_{\hg,\hh}\in H^1_{\bal}(\mathbb Q, \mathbb V_{\hg} \hat \otimes \mathbb V_{\hh} (2-\mathbf{t}_1))
    \]
    attached to the pair $(\hg,\hh)$, with the conventions that we later recall. Similarly, we may consider a Beilinson--Flach class $\kappa_{\hg,\hh\otimes\varepsilon_K}$ attached to the families $(\hg, \hh \otimes \varepsilon_K)$, where $\hh\otimes \varepsilon_K$ denotes the twist of $\hh$ by the quadratic character $\varepsilon_K$.
\end{enumerate}


Note, however, that the construction of the Beilinson--Flach classes proceeds in an ostensibly different way, since it involves working with modular units, which are absent in the theory of diagonal cycles. Roughly speaking, the construction of diagonal cycles proceeds in purely geometric terms, so the construction is amenable to be generalized to other settings where there are no modular units (e.g. Shimura curves) and Beilinson--Flach classes are not available.

The main result of this note is a result connecting both classes. However, to ensure that they live in the same space, one first defines a Beilinson--Flach class \[ \BF(f,\hg,\hh) \in H^1(\mathbb Q, V_f \otimes \mathbb V_{\hg} \hat{\otimes} \mathbb V_{\hh}(2-{\mathbf t}_1)), \] obtained as a suitable weighted combination of $\kappa_{\hg,\hh}$ and $\kappa_{\hg,\hh \otimes \varepsilon_K}$. The intuition for that comes from the observation that $V_f \otimes \mathbb V_{\hg} \hat \otimes \mathbb V_{\hh}(2-{\mathbf t}_1)$ decomposes as the direct sum of the Galois representations $\mathbb V_{\hg} \hat \otimes \mathbb V_{\hh}(2-{\mathbf t}_1)$ and $\mathbb V_{\hg} \hat \otimes \mathbb V_{\hh}(2-{\mathbf t}_1)(\varepsilon_K)$, and hence one can naturally construct a cohomology class by gluing the Euler systems for each of the two pieces.

The key point for developing a comparison of Euler systems comes through the theory of $p$-adic $L$-functions. More precisely, the arithmetic of the Beilinson--Flach system is dictated by the Hida--Rankin $p$-adic $L$-functions, which interpolate special values of the complex $L$-function $L(g \otimes h,s)$. Similarly, diagonal cycles are related, via explicit reciprocity laws, with the triple product $p$-adic $L$-function constructed by Hsieh \cite{Hs}, and later extended by Andreatta--Iovita \cite{AI} to the non-ordinary case.

The main result of this note may be interpreted as a comparison between a cyclotomic and an anticyclotomic cohomology class attached to the representation $V_f \otimes \mathbb V_{\hg} \hat \otimes \mathbb V_{\hh}(2-{\mathbf t}_1)$. As a piece of notation, we put $\bb{V}_{f\hg\hh}^\dagger=V_f\otimes\bb{V}_{\hg}\hat{\otimes} \bb{V}_\hh (2-\mathbf{t}_1)$ and $\bb{V}_{\hg\hh}^\dagger=\bb{V}_{\hg}\hat{\otimes}\bb{V}_\hh(2-\mathbf{t}_1)$. Before stating the theorem, we introduce the following objects; their precise definitions are recalled in the main body of the article.

\begin{itemize}
    \item[-] The Selmer group $H^1_{\cl{G}\cup +}(\bb{Q}, \bb{V}_{f\hg\hh}^\dagger)$, corresponding to a certain local condition at $p$. This is defined in  \S\ref{subsec:Selmergroups}, and it is a module over $\Lambda_{\hg\hh}=\Lambda_{\hg}\hat{\otimes}\Lambda_{\hh}$, where $\Lambda_{\hg}$ (resp. $\Lambda_{\hh}$) is the Iwasawa algebra over which the Hida family $\hg$ (resp. $\hh$) is defined.
    \item[-] The Selmer group $H^1_{\bal}(\bb{Q}, \bb{V}_{f\hg\hh}^\dagger)$, corresponding to the balanced local condition at $p$.
    \item[-] The Atkin--Lehner pseudo-eigenvalue $\lambda_{N_g}(\hg)$.
    \item[-] The $p$-adic period $\Omega_{f,\gamma}$, depending on the choice of the isomorphism $\gamma$ of \eqref{eq:iso}.
    \item[-] The triple product $p$-adic $L$-function $\Lp^g(\hf,\hg,\hh)$, interpolating central values of the corresponding triple product complex $L$-functions along the $\hg$-unbalanced region. In this case, $\Lp^g(f,\hg,\hh)$ corresponds to specializing $\hf$ at $f$.
\end{itemize}

In the course of developing the theory and proving the results of this work, we impose a non-vanishing assumption on $\Lp^g(f,\hg,\hh)$; see Assumption \ref{ass:nonvanishing}.

\begin{theorem}
Assume that $H^1_{\cl{G}\cup +}(\bb{Q}, \bb{V}_{f\hg\hh}^\dagger)$ is a torsion-free $\Lambda_{\hg\hh}$-module of rank $1$ and that the $p$-adic $L$-function $\Lp^g(f,\hg,\hh)$ is not identically zero. Then $\BF(f,\hg,\hh)$ belongs to $H^1_{\bal}(\bb{Q}, \bb{V}_{f\hg\hh}^\dagger)$ and \[ \lambda_{N_g}(\hg)\cdot\BF(f,\hg,\hh)=\Omega_{f,\gamma}\cdot \Lp^g(f,\hg,\hh) \cdot \kappa(f,\hg,\hh). \]
\end{theorem}

% (Esto ya se menciona despu\'es) More generally, this equality can be considered for fixed specializations $(g,h)$ of the families $(\hg,\hh)$. 
Furthermore, the result may be also understood in the framework of the comparison between different instances of Euler systems, developed e.g. in \cite{LR1}, but where one considers instead CM families passing through a weight-one Eisenstein series, while in \emph{loc.\,cit.} the key input was the use of families passing through the critical $p$-stabilization of an Eisenstein series of weight at least two.

\begin{remark}
Contrary to the situation in \cite{BDV}, where Beilinson--Kato classes and Beilinson--Flach elements were compared in Iwasawa cohomology, we are not working at that level here. In our case, such a comparison is not possible, since the main result should be viewed as relating a cyclotomic Euler system to an anticyclotomic one. Therefore, we restrict our analysis to the bottom layers.
\end{remark}

\begin{remark}
The slogan of this result may be summarized under the sentence {\it Heegner points are to Kato classes what anticyclotomic diagonal cycles are to Beilinson--Flach classes}. Unfortunately, this setting remains much more mysterious in many instances, like the lack of a proof of the Iwasawa main conjecture (only one divisibility is known).
\end{remark}

\begin{comment}
\begin{theorem}\label{thm:main}
Assume that the complex $L$-function $L(g \otimes h, s)$ vanishes at $s=\frac{l+m-1}{2}$. Then, the following equality holds in $\mathbb Q_p$ up to multiplication by a non-zero algebraic number: \[ \log_p(\res_p(\kappa_{g,h})) = \log_p(\res_p(\kappa_{\psi,g,h})), \] where the logarithm is the usual Perrin-Riou regulator of a suitable projection of the local class followed by the pairing with the canonical differentials.
\end{theorem}

The key idea for the proof of Theorem \ref{thm:main} is introducing the diagonal cycle class $\kappa_{\hf,g,h}$ attached to $(\hf,g,h)$, where $\hf$ is a CM family such that its weight one specialization is the irregular Eisenstein series $f = E_1(1,\varepsilon_K)$. 
\end{comment}

The proof proceeds in three main steps.

\begin{enumerate}
\item[(a)] Compare the triple product $p$-adic $L$-function $\Lp^g(f,\hg,\hh)$ with the Hida--Rankin $p$-adic $L$-functions attached to the pairs $(\hg,\hh)$ and $(\hg, \hh \otimes \varepsilon_K)$.
\item[(b)] Define the weighted Beilinson--Flach class and show that it satisfies the appropriate local condition. This requires a careful analysis of the structure of the corresponding Selmer groups.
\item[(c)] Relate the two cohomology classes via the explicit reciprocity laws.
\end{enumerate}

As a consequence of this result in families, we also obtain a comparison upon specializing $\hg$ and $\hh$; see \S\ref{subsec:nonexceptional} for details. Another corollary concerns a factorization of a certain big logarithm, which can be applied to the Beilinson--Flach class. To state it, we introduce the following objects:
\begin{itemize}
    \item[-] The triple product $p$-adic $L$-function of Hsieh \cite{Hs} and \cite{AI}, denoted $\Lp^f(\hf,\hg,\hh)$, attached to the triple $(\hf,\hg,\hh)$, where $\hf$ is the CM family considered above. Its specialization at weight one in the first variable is denoted by $\Lp^f(f,\hg,\hh)$.
    \item[-] The big-logarithm map $\Log_{\omega_{\hg} \otimes \omega_{\hh}}$, introduced in Def. \ref{def:big-log}.
\end{itemize}


\begin{corollary}
Assume that $H^1_{\cl{G}\cup +}(\bb{Q}, \bb{V}_{f\hg\hh}^\dagger)$ is a torsion-free $\Lambda_{\hg\hh}$-module of rank $1$ and that the $p$-adic $L$-function $\Lp^g(f,\hg,\hh)$ is not identically zero. Then we have the following factorization of the image of the Beilinson--Flach class under the Perrin-Riou map $\Log_{\bm{\omega}_{\hg} \otimes \bm{\omega}_{\hh}}$: \[ \Omega_{f,\gamma} \cdot \Lp^g(f,\hg,\hh) \cdot \Lp^f(f, \hg, \hh) =  \lambda_{N_g}(\hg) \cdot \Log_{\omega_{\hg} \otimes \omega_{\hh}}(\BF(f, \hg, \hh)). \] 
\end{corollary}


\subsection{Exceptional zeros}

There is a case which is especially interesting from the point of view of arithmetic applications, which corresponds to the following assumption.

\begin{ass}
With the previous notations, it holds that \[ \alpha_g \beta_h = p^{\frac{l_0+m_0-3}{2}}. \] 
\end{ass}

Since $g$ and $h$ are $p$-ordinary and $p\nmid \mathrm{cond}(h)$, it follows from this assumption together with the Ramanujan--Petersson conjecture that $(l_0,m_0)=(2,1)$ and that $g$ has conductor $N_gp$. Note that this situation includes in particular the case where $g$ is the modular form attached to an elliptic curve over $\bb{Q}$ with multiplicative reduction at $p$.

Hence, we encounter an arithmetic situation that is particularly intriguing, as {\em two} of the Euler factors associated with the $p$-adic $L$-function vanish. In this setting, improved $p$-adic $L$-functions are available, both in the Hida--Rankin case and for the triple product of modular forms. However, the vanishing of these two Euler factors occurs for different reasons: one is introduced by the big logarithm map, while the other arises from the interpolation in families of Beilinson--Flach classes and diagonal cycles. To complete the picture, we therefore need to (a) construct improved big logarithm maps, and (b) construct improved cohomology classes. At this point, we rely on a standard conjecture in the theory of exceptional zeros to guarantee the existence of the improved Beilinson--Flach class in this setting, while the improved diagonal cycle class has already been constructed in \cite{BSV}.

As a piece of notation for the next proposition, we use the following terminology.
\begin{itemize}
    \item[-] The line $\mathcal C$ corresponds to the points in the pair of families with weights of the form $(l,l-1)$; see \S\ref{sec:expectional} for the precise definition.
    \item[-] The triple-product improved $p$-adic $L$-function, introduced in Theorem \ref{thm:improvedtripleproductpadicLfunction}, is denoted by $\widehat{\Lp}^g(f,\hg,\hh)$.
    \item[-] The improved diagonal cycle of \cite{BSV} is denoted by $\widehat{\kappa}(f,\hg,\hh)$.
    \item[-] The (still conjectural) improved Beilinson--Flach class is denoted by $\widehat{\mathrm{BF}}(f, \hg, \hh)$.
\end{itemize}

\begin{theorem}
Assume that $H^1(\bb{Q},\bb{V}_{f\hg\hh}^\dagger\vert_{\cl{C}})$ is a torsion-free $\cl{O}_{\cl{C}}$-module, $H^1_{\cl{G}\cup +}(\bb{Q},\bb{V}_{f\hg\hh}^\dagger)$ is a torsion-free $\Lambda_{\hg\hh}$-module of rank $1$ and $\Lp^g(f,\hg,\hh)$ is not identically zero. Then $\widehat{\mathrm{BF}}(f, \hg, \hh)$ belongs to the Selmer group $H^1_{\bal}(\bb{Q},\bb{V}_{f\hg\hh}^\dagger\vert_{\cl{C}})$ and
\[
\Omega_{f,\gamma} \cdot \widehat{\Lp}^g(f,\hg,\hh) \cdot \widehat{\kappa}(f,\hg,\hh) = \lambda_{N_g}(\hg)\cdot\widehat{\mathrm{BF}}(f, \hg, \hh).
\]
\end{theorem}

The proof follows the same ideas as in the non-exceptional case, but requires certain modifications to account for the improved setting. In particular, one needs to work with improved big logarithms, since in this situation the Euler factor in the numerator of the Perrin-Riou map vanishes at the point corresponding to the pair $(g,h)$.

\begin{comment}
Following the analogy with the cyclotomic case, it is natural to formulate a conjecture about the behaviour of the class $\kappa_{\psi,g,h}$. This may be seen as an extension of the conjecture of Perrin-Riou \cite{PR93}.

\begin{conj}
The class $\kappa_{\psi,g,h}$ has infinite order if and only if the complex Hida--Rankin $L$-function $L(g \otimes h,s)$ has a simple zero at $s=\frac{l+m-1}{2}$.
\end{conj}

That is, since the anticyclotomic class $\kappa_{\psi,g,h}$ plays an analogous role to Heegner points in the setting of two modular forms, we expect that the vanishing of the Hida--Rankin $L$-function must dictate whether or not it has infinite order.
\end{comment}



\subsection{Related works}

This project, together with \cite{BDV}, provides an example of how two distinct types of Euler systems can be related through a {\it CM degeneration technique}. We now highlight its connections with other analogous phenomena:

\begin{itemize}
    \item[(1)] {\bf Critical Eisenstein series.} In \cite{LR1} and \cite{PR25}, the authors study Euler systems in families where one of the modular forms passes through a point corresponding to the critical-slope $p$-stabilization of an Eisenstein series. In these cases, the specialization of the associated Galois representation admits a projection onto a one-dimensional quotient, allowing for a direct comparison with a {\it smaller} Euler system.

    \item[(2)] {\bf Irregular Euler systems.} The degeneration technique used in this work, inspired by \cite{BSV}, considers a CM family intersecting an irregular weight-one Eisenstein series. This can be viewed as a {\it CM degeneration}, in which the Euler system in families decomposes into a sum of two distinct Euler systems, each associated with a different component of the CM representation. In both situations, the key idea is the same: to exploit a cuspidal Hida or Coleman family which, at a certain point, specializes to an Eisenstein series.
    
    \item[(3)] {\bf (Eisenstein) congruences among modular forms.} Instead of working inside a family, one may also consider an Euler system attached to a cuspidal modular form that satisfies a congruence relation with an Eisenstein series. In this setting, similar connections to those described in (1) can be established.  

\end{itemize}

\subsection{Acknowledgements}
We are deeply grateful to Francesc Castella for stimulating discussions related to the subject of this work. We also thank K\^az\i m B\"uy\"ukboduk for many insightful comments and suggestions. The third-named author further thanks David Loeffler and Victor Rotger for numerous conversations that inspired several of the ideas developed in this note.

During the preparation of this article, the authors were supported by PID2023-148398NA-I00, funded by MCIU/AEI/10.13039/501100011033/FEDER, UE. The first author was also supported by Taighde \'Eireann -- Research Ireland under Grant number IRCLA/2023/849 (HighCritical).


\section{Preliminaries}

Along this section we recall the main properties about the Galois representations attached to modular forms and families which are needed along this work. Further, we state the assumptions which are needed to guarantee the \'etaleness of the Coleman--Mazur--Buzzard eigencurve around a classical point, following the seminal works of \cite{bellaiche-dimitrov} and \cite{BeDiPo}. This has a special relevance in our study, since we are going to consider a family passing through an irregular weight-one Eisenstein series.

Let $p>2$ be a prime and fix once and for all an embedding $\iota_p:\overline{\bb{Q}}\xrightarrow{} \overline{\bb{Q}}_p$ and an embedding $\iota: \overline{\bb{Q}}\hookrightarrow \bb{C}$. The latter determines a choice of complex conjugation in $G_{\bb{Q}}$ which will be denoted by $c$.

\subsection{Deligne representations}

Let $\xi=\sum_{n=1}^\infty a_n(\xi)q^n\in S_k(N_\xi,\chi_\xi)$ be a normalized newform of weight $k\geq 2$, level $N_\xi$, and nebentypus $\chi_\xi$. %Let $p>2$ be a prime and let $E=L_\mathfrak{P}$ be a finite extension of $\mathbb Q_p$ with ring of integers $\mathcal{O}$ arising as the completion of the Hecke field $L$ of $f$ at a prime $\mathfrak{P}$ of $L$ above $p$. \textcolor{red}{En otras secciones $L$ es una extensi\'on de $\bb{Q}_p$. Yo quiz\'a adoptar\'ia la misma convenci\'on aqu\'i.}
Let $L$ be a finite extension of $\bb{Q}_p$ containing the Fourier coefficients of $\xi$ (under the embedding $\iota_p$ fixed above) and let $\cl{O}$ be the ring of integers of $L$. By work of Eichler--Shimura and Deligne, there is a two-dimensional representation
$$
\rho_\xi\,:\,G_\mathbb Q\longrightarrow \GL_L(V_\xi)\simeq\GL_2(L)
$$
unramified outside $pN_\xi$ and characterized by the property
$$
{\rm trace}\,\rho_\xi(\Fr_{q})=a_{q}(\xi)
$$
for all primes $q\nmid pN_\xi$, where $\Fr_{q}$ denotes an arithmetic Frobenius element at $q$.
Let $Y_1(N_\xi)$ be the open modular curve over $\mathbb Q$ parameterizing pairs $(A,P)$ consisting of an elliptic curve $A$ and a point $P\in A$ of order $N_\xi$. Let $\mathscr{L}_{k-2}$ be the $p$-adic sheaf over $Y_1(N_\xi)$ introduced in \cite[\S2.3]{BSV}. We shall work with the geometric realization of $V_\xi$ arising as the maximal quotient of
\[
H^1_{\et}(Y_{1}(N_\xi)_{\overline{\mathbb Q}},\mathscr{L}_{k-2}(1))\otimes_{\Z_p}L
\]
on which the dual Hecke operators $T_q'$ and $\langle d\rangle'$ act as multiplication by $a_q(\xi)$ and $\chi_\xi(d)$ for all primes $q\nmid N_\xi$ and all $d\in(\Z/N_\xi\Z)^\times$. %Let $T_\xi\subset V_\xi$ be the $\mathcal{O}$-lattice defined by the natural image of
%\begin{align*}
%H^1_{\et}(Y_{1}(N_\xi)_{\overline{\mathbb Q}}, \mathscr{L}_{k-2}(1))\otimes_{\Z_p}\cl{O}
%\end{align*}
%under the quotient map $H^1_{\et}(Y_{1}(N_\xi)_{\overline{\mathbb Q}},\mathscr{L}_{k-2}(1))\otimes_{\Z_p}L\twoheadrightarrow V_\xi$.

%\textcolor{blue}{No estoy seguro de que necesitemos la realizaci\'on geom\'etrica ni el ret\'iculo $T$.}

%In the case where the modular forms is of finite slope, and not necessarily ordinary, there is also the same notion of Galois representation, but some of its properties are subtler and need to use the theory of $(\varphi,\Gamma)$-modules, as developed by Kedlaya--Pottharst--Xiao in \cite{KPX}.

If $\xi$ is $p$-ordinary, then we have a filtration
\[
0\rightarrow V_\xi^+\rightarrow V_\xi \rightarrow V_\xi^-\rightarrow 0
\]
of $G_{\bb{Q}_p}$-representations, where $V_\xi^{\pm}$ are $1$-dimensional and $V_\xi^-$ is unramified with $\Fr_p$ acting as multiplication by the unit root of the $p$-th Hecke polynomial of $\xi$.

%We write $D(V_f)$ for the Dieudonn\'e module corresponding to the modular form $f$; this will arise later in the study of the Perrin-Riou maps.

\subsection{Hida families and big Galois representations}\label{subsec:families}

In this subsection, we recall the fundamental concepts and main notations related to Hida families that will appear in this work.

Let $\Lambda=\bb{Z}_p[[1+p\bb{Z}_p]]$. Given an integer $r$ and a finite-order character $\epsilon:1+p\bb{Z}_p\rightarrow \overline{\bb{Q}}_p^\times$, we define the character $\nu_{r,\epsilon}:1+p\bb{Z}_p\rightarrow \overline{\bb{Q}}_p^\times$ by $\nu_{r,\epsilon}(z)=z^r\epsilon(z)$. This character extends to a ring homomorphism $P_{r,\epsilon}:\Lambda \rightarrow \overline{\bb{Q}}_p$, and we will use the same notation $P_{r,\epsilon}$ to denote the corresponding point in $\mathrm{Spec}(\Lambda)(\overline{\bb{Q}}_p)$. Points in $\mathrm{Spec}(\Lambda)(\overline{\bb{Q}}_p)$ of this form will be called \textit{arithmetic points}, and the set of such points will be denoted by $\cl{W}$. Given a finite flat extension $\cl{R}$ of $\Lambda$, we say that a point $x\in \mathrm{Spec}(\cl{R})(\overline{\bb{Q}}_p)$ is \emph{arithmetic} if it lies above an arithmetic point $P_{r,\epsilon}\in \mathrm{Spec}(\Lambda)(\overline{\bb{Q}}_p)$. In this case, we refer to $k_x=r+2$ as the \emph{weight} of $x$.

Let $N$ be a positive integer coprime to $p$ and let $\chi:(\bb{Z}/Np\bb{Z})^\times \rightarrow \overline{\bb{Q}}^\times$ be a Dirichlet character. A \emph{Hida family} of tame level $N$ and character $\chi$ is a formal $q$-expansion
\[
\bm{\xi}=\sum_{n\geq 1} a_n(\bm{\xi}) q^n \in \Lambda_{\bm{\xi}}[[q]]
\]
with coefficients in a normal domain $\Lambda_{\bm{\xi}}$ finite flat over $\Lambda$ such that for every arithmetic point $x\in \mathrm{Spec}(\Lambda_{\bm{\xi}})(\overline{\bb{Q}}_p)$ lying above some $P_{r,\epsilon}$ with $r\geq 0$ the corresponding specialization $\bm{\xi}_x$ is a $p$-ordinary eigenform in $S_{r+2}(Np^s,\chi \epsilon \omega^{-r})$, where $s=\max\lbrace 1,\ord_p(\mathrm{cond}(\epsilon)\rbrace$. We define $\cl{O}_{\bm{\xi}}=\Lambda_{\bm{\xi}}[1/p]$. We set $U_{\bm{\xi}}=\mathrm{Spec}(\Lambda_{\bm{\xi}})$ and $\cl{U}_{\bm{\xi}}=\mathrm{Spec}(\cl{O}_{\bm{\xi}})$ and we denote by $\kappa_{\bm{\xi}}:U_{\bm{\xi}}\rightarrow \mathrm{Spec}(\Lambda)$ the homomorphism corresponding to the inclusion $\Lambda \xhookrightarrow{} \Lambda_{\bm{\xi}}$. We denote by $\cl{W}_{\bm{\xi}}$ the set of arithmetic points in $U_{\bm{\xi}}(\overline{\bb{Q}}_p)$. We say that an arithmetic point $x\in \cl{W}_{\bm{\xi}}$ of weight $k_x\geq 1$ is \emph{classical} if the corresponding specialization $\bm{\xi}_x$ is a classical modular form (which is always the case if $k_x\geq 2$). In this case $\bm{\xi}_x$ is a $p$-stabilized newform. If $\bm{\xi}_x$ is the $p$-stabilization of a newform of level $N$, we denote the latter by $\bm{\xi}_x^{\circ}$. Otherwise, we write $\bm{\xi}_x^\circ=\bm{\xi}_x$. We note that, if $k_x>2$, then $\bm{\xi}_x$ is always $p$-old.

%Let $\bm{\xi}$ be a Hida family of tame level $N_{\hf}$ and tame character $\chi_{\hf}$ parametrized by a connected affinoid disc $U_{\hf}$ centered at an integer $k_0 \geq 1$ in the weight space $\mathcal W_L = \mathcal W \times_{\mathbb Q_p} L$ over a finite extension $L$ of $\mathbb Q_p$. We denote by $\mathcal{O}_\hf$ the ring of analytic functions on $U_\hf$. We have that $\hf = \sum_{n \geq 1} a_n(\hf)q^n$ is a formal $q$-expansion with coefficients in $\mathcal O_\hf$, the ring of analytic functions on $U_{\hf}$.

%For each integer $k$ in a cofinite subset $U_{\hf}^{\cla}$ of $U_{\hf} \cap \mathbb Z^{\geq k_0}$, the specialization $\hf_k$ is a modular form of weight $k$, level $\Gamma_1(N_{\hf}) \cap \Gamma_0(p)$ and character $\chi_{\hf}$. The $p$-adic valuation of $a_p(\hf_k)$ is independent of $k$. If $\hf_k$ is the $p$-stabilization of a modular form of level $\Gamma_1(N_{\hf})$, we denote the latter by $\hf_k^{\circ}$
%(in particular, this is true if $k>2$ \textcolor{blue}{esto es solo para familias de Hida, para familias de Coleman la condici\'on $k>2$ creo que se ha de cambiar por una condici\'on que depende de la pendiente}).
%Otherwise, we write $\hf_k^\circ=\hf_k$.

% (Creo que al final no usamos esto. Lo dejo comentado) For each classical point $k\geq 2$, we have a Galois representation $V(\hf_k^\circ)$ defined as in the preceding section. We say that a classical point $k \geq 2$ is {\it good} if $p$ does not divide the conductor of $\hf_k^\circ$, the $p$-th Hecke polynomial $X^2- a_p(\hf_k^\circ) \cdot X + \chi_{\hf}(p)p^{k-1}$ has distinct roots and $\hf_k^\circ$ is not theta-critical. Assume that $k$ is good and let $\xi=\hf_k$. Then there is a short exact sequence \[ 0 \longrightarrow D(\xi)_{\alpha}^+ \longrightarrow D(\xi) \longrightarrow D(\xi)_{\alpha}^- \longrightarrow 0 \] of $(\varphi,\Gamma)$-modules over the Robba ring $\mathcal R_L = \mathcal R \otimes_{\mathbb Q_p} L$, where $D(\xi)_{\alpha}^{\pm}$ is isomorphic to a concrete $(\varphi,\Gamma)$-module associated with a certain character $\delta_{\xi,\alpha}^{\pm}$.

%Let $U_{\hf} \hookrightarrow \mathcal W_L$ be a connected open disc centered at $k_0$, 

Attached to a Hida family $\bm{\xi}$ % of tame level $N_{\bm{\xi}}$ and character $\chi_{\bm{\xi}}$
there is a locally-free rank-two $\Lambda_{\bm{\xi}}$-module $\bb{V}_{\bm{\xi}}$ equipped with a continuous action of $G_{\bb{Q}}$ such that for any arithmetic point $x\in \cl{W}_{\bm{\xi}}$ of weight $k_x\geq 2$ the specialization $\bb{V}_{\bm{\xi}}\otimes_{\Lambda_{\bm{\xi}},x}\overline{\bb{Q}}_p$ recovers the $G_{\bb{Q}}$-representation $V_{\bm{\xi}_x^\circ}$ attached to $\bm{\xi}_{x}^\circ$. We call $\bb{V}_{\bm{\xi}}$ the \emph{big Galois representation} attached to $\bm{\xi}$. As a $G_{\bb{Q}_p}$-representation, $\bb{V}_{\bm{\xi}}$ admits a filtration
\[
0\rightarrow \bb{V}_{\bm{\xi}}^+\rightarrow \bb{V}_{\bm{\xi}}\rightarrow \bb{V}_{\bm{\xi}}^-\rightarrow 0,
\]
where $\bb{V}_{\bm{\xi}}^{+}$ is a free rank-one $\Lambda_{\bm{\xi}}$-module, $\bb{V}_{\bm{\xi}}^-$ is a locally-free rank-one $\Lambda_{\bm{\xi}}$-module, and the $G_{\bb{Q}_p}$-action on $\bb{V}_{\bm{\xi}}^-$ is unramified with $\Fr_p$ acting as multiplication by $a_p(\bm{\xi})$. For any arithmetic point $x\in \cl{W}_{\bm{\xi}}$ of weight $k_x\geq 2$, the corresponding specialization of the exact sequence above recovers the filtration
\[
0\rightarrow V_{\bm{\xi}_x^\circ}^+\rightarrow V_{\bm{\xi}_x^\circ}\rightarrow V_{\bm{\xi}_x^\circ}^-\rightarrow 0.
\]
Moreover, if $\bm{\xi}$ is residually irreducible and $p$-distinguished, then the $\Lambda_{\bm{\xi}}$-modules $\bb{V}_{\bm{\xi}}$ and $\bb{V}_{\bm{\xi}}^{-}$ are free.

%Let $\mathbb V(\hf)$ be the big Galois representation attached to $\hf$. It is a free $\mathcal O_{\hf}$-module of rank two such that for each good classical point $k$ in $U_{\hf}^{\cla}$ there are canonical specialization isomorphisms \[ \rho_k \colon V(\hf) \otimes_k L \cong \mathbb V(\hf_k). \] The same applies verbatim to the dual Galois representation $\mathbb V^*(\hf)$, and there exists a perfect $G_{\mathbb Q}$-equivariant pairing \[ \langle \cdot, \cdot \rangle_{\hf} \colon \mathbb V(\hf) \otimes_{\mathcal O_{\hf}} \mathbb V^*(\hf) \longrightarrow \mathcal O_{\hf}, \] compatible with the usual dualities at each specialization.

\begin{comment}
\subsection{Cyclotomic and anticyclotomic twists}\label{sec:twists}

\textcolor{blue}{La notaci\'on introducida aqu\'i no parece que se use despu\'es.}

In this section we introduce the notations regarding the cyclotomic and antityclotomic twists we will be using in forthcoming sections.

Write $\langle a; b\rangle$, with $a\in(\mathbb Z/N\mathbb Z)^\times$ and $b\in(\mathbb Z/p^r\mathbb Z)^\times$, for the diamond operator $\langle d\rangle$, where $d\in(\mathbb Z/Np^r)^\times$ is congruent to $a$ modulo $N$ and to $b$ modulo $p^r$. We also write $\epsilon_N \colon G_\mathbb Q\rightarrow (\mathbb Z/N\mathbb Z)^\times$ for the mod $N$ cyclotomic character.

For any positive integer $r$, let
$$
G_r=1+p(\mathbb Z/p^r\mathbb Z), \quad \tilde{G}_r=(\mathbb Z/p^r\mathbb Z)^\times,
$$
and define
$$
\Lambda_r=\mathbb Z_p[G_r],\quad \tilde{\Lambda}_r=\mathbb Z_p[\tilde{G}_r], \quad \Lambda=\varprojlim_r\Lambda_r=\mathbb Z_p[\![1+p\mathbb Z_p]\!],\quad \tilde{\Lambda}=\varprojlim_r \tilde{\Lambda}_r=\mathbb Z_p[\![\mathbb Z_p^\times]\!].
$$
We have natural factorizations $(\mathbb Z/p^r\mathbb Z)^\times=\mu_{p-1}\times(1+p\mathbb Z/p^r\mathbb Z)$ and $\mathbb Z_p^\times=\mu_{p-1}\times(1+p\mathbb Z_p)$ which give natural embeddings $\Lambda_r\xhookrightarrow{} \tilde{\Lambda}_r$ and $\Lambda\xhookrightarrow{} \tilde{\Lambda}$. For any integer $i$ modulo $p-1$, let $\kappa_i:\mathbb Z_p^\times\rightarrow \Lambda^\times$ be the character defined by $z\mapsto \omega^i(z)[\langle z\rangle]$ and let $\bm{\kappa}_i=\kappa_i\circ\epsilon_{\text{cyc}}:G_\mathbb Q\rightarrow \Lambda^\times$.

There is a character $\kappa_{r_1}:\mathbb Z_p^\times \rightarrow \Lambda^\times$ given by $u\mapsto \omega^{r_1}(u)[\langle u\rangle]$. Fix a square root $\kappa_{r_1}^{1/2}$ of this character given by $u\mapsto \omega^{r_1/2}(u)[\langle u\rangle^{1/2}]$. If $\Lambda_{\hf}$ is the finite extension of $\Lambda$ over which a certain Hida family $\hf$ is defined, we let $\kappa_\hf$ and $\kappa_\hf^{1/2}$ be the composition of $\kappa_{r_1}$ and $\kappa_{r_1}^{1/2}$, respectively, with the embedding $\Lambda^\times\xhookrightarrow{} \Lambda_\hf^\times$.
\end{comment}

\subsection{The Coleman--Mazur eigencurve}

Let $\bm{\xi}$ be a Hida family of tame level $N_{\bm{\xi}}$ and character $\chi_{\bm{\xi}}$. Let $x\in \cl{W}_{\bm{\xi}}$ be a classical point of weight $k_x\geq 1$. We say that $x$ is a \emph{crystalline} point if $\bm{\xi}_x$ has nebentypus of conductor coprime to $p$. We say that $x$ is a \emph{good} point if in addition any of the following conditions holds: %  We say that $\xi$ is a \emph{crystalline} specialization of $\bm{\xi}$ if $\xi$ is has nebentypus of conductor coprime to $p$. Moreover, we say that $\xi$ is \emph{good} if any of the following conditions holds:
\begin{enumerate}
\item[(a)] $k_x \geq 2$;% and $\bm{\xi}_{u_0}$ is a non-critical $p$-regular eigenform.
\item[(b)] $k_x=1$ and $\bm{\xi}_{x}$ is a $p$-stabilisation of a classical $p$-regular cuspidal weight-one newform of level $N_{\bm{\xi}}$ without real multiplication by a real  quadratic field $F$ where the prime $p$ splits;
\item[(c)] $k_x=1$ and $\bm{\xi}_{x}$ is the $p$-stabilisation of a $p$-irregular weight-one Eisenstein series of conductor $N_{\bm{\xi}}$.
\end{enumerate}

Note that these conditions imply in particular that $\bm{\xi}_{x}$ is an \'etale point of the cuspidal Coleman--Mazur--Buzzard $p$-adic eigencurve of tame level $N_{\bm{\xi}}$. When (a) holds, this follows from the work of Hida and Coleman (see, e.g. \cite{bellaiche11a}); when (b) holds, this is proved in \cite{bellaiche-dimitrov}; when (c) holds, it is proved in \cite{BeDiPo}.

If $x\in \cl{W}_{\bm{\xi}}$ is a good point of weight $1$, then it follows from \cite[Prop.~2.2]{BDV} that $\bb{V}_{\bm{\xi}}\otimes_{\Lambda_{\bm{\xi}},x}\overline{\bb{Q}}_p$ recovers the Deligne--Serre Artin representation $V_{\bm{\xi}_x^\circ}$ attached to $\bm{\xi}_{x}^\circ$. In particular, this allows us to obtain a filtration
\[
0\rightarrow V_{\bm{\xi}_x^\circ}^+ \rightarrow V_{\bm{\xi}_x^\circ} \rightarrow V_{\bm{\xi}_x^\circ}^- \rightarrow 0
\]
for $V_{\bm{\xi}_x^\circ}$ as a $G_{\bb{Q}_p}$-representation by specializing the filtration
\[
0\rightarrow \bb{V}_{\bm{\xi}}^+ \rightarrow \bb{V}_{\bm{\xi}} \rightarrow \bb{V}_{\bm{\xi}}^- \rightarrow 0
\]
at the point $x$.
%Let $f \in S_k(N_f, \chi_f)$ be a newform of weight $k \geq 1$. Let $\hf$ be a Coleman family passing through a $p$-stabilization of $f$; we put $\hf_x^{\circ} = f$. Along this work, we are going to impose the following assumption.

\begin{comment}

\begin{ass}\label{ass:etale}
The eigenform $\bm{\xi}_{u_0}$ is an \'etale point of the cuspidal Coleman--Mazur--Buzzard $p$-adic eigencurve of tame level $N_{\bm{\xi}}$.
\end{ass}

This assumption is satisfied in many cases. In particular, it is satisfied if any of the following conditions hold:
\begin{enumerate}
\item[(a)] $k \geq 2$ and $\bm{\xi}_{u_0}$ is a non-critical $p$-regular eigenform. This follows from the work of Hida and Coleman; see e.g. \cite{bellaiche11a}.
\item[(b)] $k=1$ and $\bm{\xi}_{u_0}$ is a $p$-stabilisation of a classical $p$-regular cuspidal weight-one newform of level $N_{\bm{\xi}}$ without real multiplication by a real quadratic field $F$ where the prime $p$ splits. This is the case considered in \cite{bellaiche-dimitrov}.
\item[(c)] $k=1$ and $\bm{\xi}_{u_0}$ is the $p$-stabilisation of a $p$-irregular weight-one Eisenstein series of conductor $N_{\bm{\xi}}$. This is the case considered in \cite{BeDiPo}.
\end{enumerate}

In particular  we can consider specializations of families in weight one points, and they inherit the properties regarding the existence of filtrations. See e.g. \cite[\S2]{DR2b}. \textcolor{blue}{Para poder usar la proposici\'on 2.2 de BDV hay que asumir alguna de las hip\'otesis de arriba, y no solo la Assumption 2.1. Es cierto que DR2b parece afirmar algo parecido, pero no s\'e si me fiar\'ia mucho.}  \textcolor{red}{Cierto, mejor no meterse a usar ese artículo que no demuestran mucho. En cualquier caso aqui no nos hace falta.} \textcolor{blue}{Bueno, hace falta para considerar especializaciones en peso 1. As\'i que quiz\'a ser\'ia mejor asumir que se cumple una de las tres condiciones de arriba.} \textcolor{red}{Si si, me referia a que tipicamente vamos a trabajar en el Eisenstein irregular case, que no es problema precisamente suponer esto.} Note however that the assumption is not always true. Recent works of Betina and Dimitrov, and Betina, Maksoud and Pozzi show how this condition can fail for irregular weight-one points, either in the CM case \cite{betina-dimitrov} or in a much more general set-up.

\end{comment}


\subsection{Irregular weight-one Eisenstein series}\label{subsec:wt1}

Let $K$ be an imaginary quadratic field of discriminant $-d_K$ in which $p$ splits. Write $p \mathcal O_K = \mathfrak p \bar{\mathfrak p}$, with $\mathfrak p$ the prime of $\mathcal O_K$ corresponding to the fixed embedding $\iota_p \colon \overline{\mathbb Q} \hookrightarrow \overline{\mathbb Q}_p$. Let $\varepsilon_K$ stand for the quadratic character attached to the imaginary quadratic field $K$. In particular, for any prime $q\nmid d_K$, we have $\varepsilon_K(q) = +1$ if $q$ splits in $K$ and $\varepsilon_K(q) = -1$ if $q$ is inert in $K$.

\begin{defi}
The weight-one Eisenstein series of character $\varepsilon_K$ is the weight-one modular form given by the Fourier expansion \[ \Eis_1(\varepsilon_K) = \frac{1}{2} L(\varepsilon_K,0) + \sum_{n \geq 1} q^n \sum_{d|n} \varepsilon_K(d) \in M_1(-d_K, \varepsilon_K). \] It has level $\Gamma_1(-d_K)$ and character $\varepsilon_K$. 
\end{defi} 

%Cambio la notaci\'on para que sea consistente con otras secciones
We will write $f=\Eis_1(\varepsilon_K)$. Note that the eigenform $f$ is a $p$-irregular modular form, i.e., its $p$-th Hecke polynomial, given by $X^2-2X+1$, has a double root. Define \[ f_\alpha = f(q) - f(q^p) \in M_1(-pd_K, \varepsilon_K) \] as the unique $p$-stabilisation of $f$.

A result of Betina, Dimitrov and Pozzi \cite{pozzi-thesis}, \cite{BeDiPo} extending the prior work of Bella\"iche and Dimitrov \cite{bellaiche-dimitrov}, establishes the following.

\begin{propo}
The eigenform $f_\alpha$ is an \'etale point of the cuspidal Coleman--Mazur eigencurve. In particular, there exists a unique (up to conjugation) Hida family $\hf$ of tame level $-d_K$ and tame character $\chi_f = \varepsilon_K$ passing through $f_\alpha$. Moreover, the Hida family $\hf$ has complex multiplication by $K$.
\end{propo}
\begin{proof}
This is \cite[Thm. A]{BeDiPo}.
\end{proof}

\begin{remark}
Similar results exist if we consider instead the Eisenstein series $\Eis_1(\chi_1,\chi_2 \varepsilon_K)$, with $\chi_1$ and $\chi_2$ Dirichlet characters of prime-to-$p$ conductors $N_1$ and $N_2$ with $\chi_1(p)=\chi_2(p)$. In that case, one would need to define $f_\alpha$ as \[ f_\alpha = \Eis_1(\varepsilon_K)(q) - \chi_1(p) \Eis_1(\varepsilon_K)(q^p) \in M_1(-pd_KN_1N_2, \varepsilon_K \chi_1 \chi_2). \]
\end{remark}


%\textcolor{blue}{Yo modificar\'a la descripci\'on de $\hf$, que est\'a ahora en el lenguaje de BSV, y usar\'ia el lenguaje de BSTW, que es el que hace falta para hacer teor\'ia de Iwasawa anticiclot\'omica. Puesto que lo que necesitamos saber sobre $\hf$ es m\'inimo, quiz\'a me restringir\'ia a lo esencial.}

Let $K^{\frk{p}}_{\infty}$ be the $\bb{Z}_p$-extension of $K$ unramified away from $\frk{p}$. Let $\Gamma_{\frk{p}}=\Gal(K^{\frk{p}}_{\infty}/K)$ and let $\Lambda_\hf=\cl{O}[[\Gamma_{\frk{p}}]]$. Let $\rec_K:\bb{A}_K^\times\rightarrow G_K^{\ab}$ denote the (arithmetically normalized) global reciprocity map and let $\Theta_\frk{p}$ denote the composition of $\rec_K$ with the canonical projection $G_K^{\ab} \twoheadrightarrow \Gamma_{\frk{p}}$. Let $\rec_{K,\frk{p}}:K_{\frk{p}}^\times \rightarrow G_{K_{\frk{p}}}^{\ab}$ denote the (arithmetically normalized) local reciprocity map. We identify $G_{K_\frk{p}}$ with the decomposition group above $\frk{p}$ determined by our fixed embedding $\overline{\bb{Q}}\hookrightarrow \overline{\bb{Q}}_p$, and thus we regard $G_{K_\frk{p}}$ as a subgroup of $G_K$. Thus we obtain a homomorphism $G_{K_\frk{p}}^{\ab} \rightarrow G_K^{\ab} \twoheadrightarrow \Gamma_{\frk{p}}$ (which is surjective if $p\nmid h_K$) and we define $\theta_\frk{p}:K_\frk{p}^\times\rightarrow \Gamma_{\frk{p}}$ as the composition of $\rec_{K,\frk{p}}$ with this homomorphism. The character $1+p\bb{Z}_p\rightarrow \Lambda_{\hf}^\times$ defined by $z\mapsto z^{-1}\theta_{\frk{p}}(z)$ extends to an embedding $\Lambda\hookrightarrow \Lambda_{\hf}$, whereby $\Lambda_{\hf}$ becomes a finite flat extension of $\Lambda$. For each non-zero fractional ideal $\frk{a}$ of $K$, let $x_\frk{a}\in \bb{A}_{K,f}^\times$ be a finite id\`ele with $\ord_w(x_{\frk{a},w})=\ord_w(\frk{a})$ at each finite place $w$ with $\ord_w(\frk{a})\neq 0$ and $x_{\frk{a},w}=1$ for all other places $w$. Then, the Hida family $\hf$ is given by
\[
\hf=\sum_{\frk{a}\subseteq \cl{O}_K,\,\frk{p}\nmid\frk{a}}\Theta_\frk{p}(x_\frk{a}) q^{N_{K/\bb{Q}}(\frk{a})} \in \Lambda_\hf[[q]].
\]
The specialization of $\hf$ at the point $x_0\in \cl{W}_{\hf}$ corresponding to the trivial character of $\Gamma_{\frk{p}}$ recovers the modular form $f_\alpha$.


\begin{comment}
We still need to recall some notations to explain the properties and the behavior of the family $\hf$. Let $\mathbb A_K^{\times}$ be the group of id\`eles of $K$ and set $U_{\mathfrak p} = K^{\times} \cdot \mathbb C^{\times} \cdot \prod_{\mathfrak q \neq \mathfrak p} \mathcal O_{\mathfrak q}^{\times} \cdot \mu_{\mathfrak p}$, where $\mu_{\mathfrak p} = \mu_{p-1}$ is the torsion subgroup of $\mathcal{O}_{\frk{p}}^\times$. Let $G_{\mathfrak p} = \mathbb A_K^{\times}/U_{\mathfrak p}$ and fix an extension \[ \varphi_{\mathfrak p} \colon \mathbb A_K^{\times}/K^{\times} \twoheadrightarrow G_{\mathfrak p} \rightarrow \overline{\mathbb Q}_p^{\times} \] of the character $1+p\mathcal{O}_{\mathfrak{p}}\rightarrow \overline{\mathbb{Q}}_p^\times$ defined by $u\mapsto u^{-1}$. Then $\varphi_{\mathfrak p}$ is an algebraic $p$-adic Hecke character of weights $(1,0)$, conductor $\mathfrak p$ and central character the Teichm\"uller lift $\omega$. Let $\psi_{\mathfrak p} \colon \mathbb A_K^{\times}/K^{\times} \rightarrow \mathbb C^{\times}$ be the associated algebraic Hecke character of infinity type $(1,0)$ and conductor $\mathfrak p$. We also denote by $\psi_{\frk{p}}$ the corresponding character $I_K(\frk{p})\rightarrow \bb{C}^\times$.

Let $\langle \psi_\frk{p}\rangle$ denote the composition of $\psi_{\frk{p}}$ with the projection onto the group of principal units of $\cl{O}$, and let $\bm{\psi}:I_K(\frk{p})\rightarrow \cl{O}_{\hf}^\times$ be the unique character such that $\bm{\psi}(\frk{a})(l)=\langle\psi_\frk{p}\rangle (\frk{p})^{l-1}$ for all $\frk{a}$ in $I_K(\frk{p})$ and for all $l$ in $U_\hf\cap\bb{Z}_{\geq 1}$. Then the Hida family $\hf$ is given by
\[ \hf = \sum_{(\frk{a},\frk{p})=1} \bm{\psi}(\mathfrak a) \cdot q^{\mathbb N \mathfrak a}, \] where the sum is over the non-zero ideals $\frk{a}$ of $\mathcal O_K$ coprime to $\mathfrak p$.

%\begin{propo}
%Consider the character $\bm{\psi} \colon I_K(\mathfrak p) \rightarrow \mathcal O_{\hf}^{\times}$. Then, the family $\hf$ may be expressed as \[ \hf = \sum \bm{\psi}(\mathfrak a) \cdot q^{\mathbb N \mathfrak a}, \] where the sum is over the non-zero ideal of $\mathcal O_K$ coprime to $\mathfrak p$.
%\end{propo}

%\begin{proof}
%This is immediate from the definition of the character $\psi$.
%\end{proof}

For each integer multiple $m$ of $p-1$ we can extend the $m$-th power of $\psi_{\mathfrak p}$ to a Hecke character $\psi_m \colon I_K \rightarrow \overline{\mathbb Q}^{\times}$ of infinity type $(m,0)$ and trivial conductor.  
Then, for each integer $l$ in $U_{\hf}\cap \bb{Z}_{> 1}$ congruent to $1$ modulo $q_E-1$ (where $q_E$ is the cardinality of the residue field of $E$) the specialization of $\hf$ at $l$ is the ordinary $p$-stabilization of the theta series
\[
\theta(\psi_{l-1}) = \sum_{\mathfrak a} \psi_{l-1}(\mathfrak a) \cdot q^{\mathbb N \mathfrak a} \in S_{l}(-d_K,\varepsilon_K),
\]
where the sum is over non-zero ideals $\frk{a}$ of $\cl{O}_K$.

%\begin{propo}
%The specializations of at $\psi_m$ are given by \[ \theta(\psi_m) = \sum_{\mathfrak a \subseteq \mathcal O_K} \psi_m(\mathfrak a) \cdot q^{\mathbb N \mathfrak a} \in S_{m+1}(-d_K,\varepsilon_K), \] where the sum is over all the non-zero ideals of $\mathcal O_K$ and $\mathbb N \mathfrak a$ is the norm of the ideal $\mathfrak a$.
%\end{propo}

%\begin{proof}
%This follows from \cite[Thm. 4.8.3]{miyake},
%\end{proof}

For each integer multiple $m$ of $p-1$, let $\varphi_m:I_K\rightarrow \overline{\bb{Q}}_p^\times$ be the $p$-adic Hecke character associated with $\psi_m$. Then there exists a unique character
\[
\varphi \colon G_K \rightarrow \mathcal O_{\hf}^{\times}
\]
specializing to $\varphi_{l}$ at each integer $l$ in $U_\hf$ congruent to 1 modulo $q_E-1$.

%The $\mathcal O_{\hf}$-adic representations $V(\hf)$ and $\Ind_K^{\mathbb Q} \varphi$ are irreducible and unramified outside $d_Kp$. Further, for $\ell$ prime not dividing $d_Kp$, the arithmetic Frobenius at $\ell$ acts on them with trace $a_{\ell}(\hf)$. We also write $\varphi_m \colon G_K \rightarrow \overline{\mathbb Q}_p^{\times}$ for the $p$-adic Galois character corresponding to $\psi_m$ by global class field theory. Further, there exists a unique character \[ \varphi \colon G_K \rightarrow \mathcal O_{\hf}^{\times} \] specializing to $\varphi_{l-1}$ at each integer $l$ congruent to 1 modulo $q_L-1$, where $q_L$ is the cardinality of the residue field of $L$. We recall the following result.

\end{comment}

Let $\bm{\varphi}:G_K\twoheadrightarrow \Gamma_{\frk{p}}$ denote the canonical projection. Then, we have the following result.

\begin{proposition}\label{prop:BSTW}
The $\Lambda_{\hf}[G_{\mathbb Q}]$-modules $\bb{V}_{\hf}$ and $\Ind_K^{\mathbb Q} \Lambda_\hf(\bm{\varphi})$ are isomorphic.
\end{proposition}
\begin{proof}
This is \cite[Thm.~1.21]{BSTW}. %\textcolor{red}{Creo que el resultado de BDV no es suficiente si queremos por ejemplo interpretar $\kappa(\hf,\hg,\hh)$ como una clase sobre la extensi\'on anticiclot\'omica. Este es un problema que ya ten\'iamos en ACR, que en BSV se restring\'ian a un disco demasiado pequeno. Yo citar\'ia en su lugar \cite{BSTW}. No s\'e c\'omo lo ves?} \textcolor{blue}{Si si, toda la razon en esto.}
\end{proof}

Fix once and for all an isomorphism of $\Lambda_{\hf}[G_{\mathbb Q}]$-modules 
\begin{equation}\label{eq:iso}
\gamma \, : \, \bb{V}_{\hf} \cong \Ind_K^{\mathbb Q} \bm{\varphi}.
\end{equation}
This isomorphism is not canonical, so the next decompositions will depend on this choice.
Since $p$ splits in $K$, the restrictions of $\Ind_K^{\mathbb Q} \bm{\varphi}$ to $G_K$ and $G_{\mathbb Q_p}$ decompose as the direct sum of $\bm{\varphi}$ and its complex conjugate $\bm{\varphi}^c$. Note that the character $\bm{\varphi}^c|_{G_{\mathbb Q_p}}$ is unramified and maps the arithmetic Frobenius $\Fr_p$ to the $p$-th Fourier coefficient $a_p(\hf) = \Theta(x_{\bar{\mathfrak p}})$ of $\hf$. Therefore, the restriction of $\bb{V}_{\hf}$ to $G_{\bb{Q}_p}$ decomposes
\[
\bb{V}_{\hf} = \bb{V}_{\hf}^+ \oplus \bb{V}_{\hf}^-, \quad \text{ with } \quad \gamma(\bb{V}_{\hf}^+) = \Lambda_{\hf}(\bm{\varphi}|_{G_{\mathbb Q_p}}) \text{ and } \gamma(\bb{V}_{\hf}^-) = \Lambda_\hf(\bm{\varphi}^c|_{G_{\mathbb Q_p}}).
\]
Recall that the specialization $\bb{V}_{\hf}\otimes_{\Lambda_{\hf},x_0} L$ recovers the Deligne--Serre Artin representation $V_f$. Setting, as before, $V_f^{\pm}=\bb{V}_{\hf}^{\pm}\otimes_{\Lambda_{\hf},x_0} L$, we have a decomposition $V_f=V_f^+\oplus V_f^-$ of $G_{\bb{Q}_p}$-representations. Also, specializing the isomorphism $\gamma$ at $x_0$, we obtain an isomorphism
\[
\gamma:V_f\cong L(\bm{1})\oplus L(\varepsilon_K) 
\]
of $L[G_{\bb{Q}}]$-modules.

Let $\bm{v}^+$ and $\bm{v}^-$ be the canonical $\Lambda_{\hf}$-bases of the $G_K$-submodules $\Lambda_{\hf}(\bm{\varphi})$ and $\Lambda_{\hf}(\bm{\varphi}^c)$ of $\Ind_K^{\mathbb Q} \Lambda_\hf(\bm{\varphi})$, respectively. The maps $\bm{v}^{\pm} \colon G_{\mathbb Q} \rightarrow \Lambda_{\hf}$ are determined by $(\bm{v}^+(1),\bm{v}^+(c))=(1,0)$ and $(\bm{v}^-(1),\bm{v}^-(c)) = (0,1)$, where $c$ denotes our fixed choice of complex conjugation. Set $\bm{v}_\hf^{\pm} = \gamma^{-1}(\bm{v}^{\pm})$ in $\bb{V}_{\hf}^{\pm}$, let $v_f^{\pm}$ in $V_f^{\pm}$ be their specializations at $x_0$ and define \[ v_{f,1} = v_f^+ + v_f^-, \quad v_{f,\varepsilon_K} = v_f^+ - v_f^-. \] Note that $c$ exchanges the vectors $\bm{v}^+$ and $\bm{v}^-$, and therefore the elements $\gamma(v_{f,1})$ and $\gamma(v_{f,\varepsilon_K})$ give bases of the $G_{\mathbb Q}$-representations $L(\bm{1})$ and $L(\varepsilon_K)$, respectively. These choices of vectors will be used later on.

\section{$p$-adic $L$-functions}

Let $(\hg,\hh)$ be a pair of cuspidal Hida families and let $\hf$ be the Hida family introduced in \S\ref{subsec:wt1}. This work relies on the comparison between the Hida--Rankin $p$-adic $L$-function $L_p(\hg,\hh)$ and the triple product $p$-adic $L$-function $L_p(\hf,\hg,\hh)$. In this section we introduce these objects.
  

%This work revolves around the comparison between the Hida--Rankin $p$-adic $L$-function and the triple product $p$-adic $L$-function in the particular case where at least one of the forms has CM and passes through an Eisenstein point. We introduce the two $p$-adic $L$-functions we will be comparing.

%Let $\hf$, $\hg$ and $\hh$ be three families of tame levels $N_{\hf}$, $N_{\hg}$ and $N_{\hh}$, and tame characters $\chi_{\hf}$, $\chi_{\hg}$ and $\chi_{\hh}$, parameterised by connected affinoid discs $U_{\hf}$, $U_{\hg}$ and $U_{\hh}$ centered at integers $(k_0,l_0,m_0)$ in the weight space $\mathcal W_L = \mathcal W \times_{\mathbb Q_p} L$ over a finite extension $L$ of $\mathbb Q_p$. 
%Most of the time, we will be considering a situation where two of the families are indeed ordinary \textcolor{blue}{(para tener las f\'ormulas de interpolaci\'on precisas necesitamos que sean ordinarias las tres, no?) si si, claro, para un resultado a la Hsieh hace falta.. dejo esto comentado}; this is an assumption which can be easily removed, but which simplifies the notation in a few points.

\subsection{Hida--Rankin $p$-adic $L$-function}\label{subsec:Hida-Rankin}

%In the setting of two Hida families, one can define a three-variable $p$-adic $L$-function, depending on the two weights and the cyclotomic twists. We recall in this section the main properties of the Hida--Rankin $p$-adic $L$-function in the ordinary case, since this suffices for our purposes in this work. For this part, we mainly follow the exposition in \cite{KLZ}.%, where both the properties of the $p$-adic $L$-function and the reciprocity law are developed in the non-ordinary setting.

Here we introduce the three-variable Hida--Rankin $p$-adic $L$-function. We follow mainly the exposition in \cite{KLZ}.

Let $(\hg,\hh)$ be a pair of cuspidal Hida families of tame levels $(N_g,N_h)$ and characters $(\chi_{\hg},\chi_{\hh})$. We make the following assumption.

\begin{ass}
    $\gcd(N_g,N_h)=1$.
\end{ass}

Put $\cl{W}_{\hg \hh \mathbf{s}}\coloneqq \cl{W}_{\hg}\times \cl{W}_{\hh}\times \cl{W}$.
For $\phi\in\lbrace \hg,\hh\rbrace$, let $\cl{W}_{\phi}^\circ$ denote the subset of crystalline points in $\cl{W}_{\phi}$. Also, fix an integer $s_0$ and define
\[
\cl{W}^\circ=\lbrace P_{s,\bm{1}}\in \cl{W}\;\vert\; s\equiv s_0 \pmod{p-1}\rbrace.
\]
In an abuse of notation, a point $P_{s,\bm{1}}\in \cl{W}^\circ$ might be denoted simply by $s$.
%denote the subset of points in $\cl{W}$ parameterising characters of $\mathbb{Z}_p^\times$ of the form $z\mapsto z^s$ with $s\in \mathbb{Z}$. %Put $\cl{W}_{\hg \hh\mathbf{s}}^{\circ}=\cl{W}_{\hg}^{\circ} \times \cl{W}_{\hh}^{\circ} \times \cl{W}^\circ\subseteq \cl{W}_{\hg \hh \mathbf{s}}$.

Let $(y_0,z_0)\in \cl{W}_{\hg}^\circ\times \cl{W}_{\hh}^\circ$ be a good crystalline point with corresponding weights $(l_0,m_0)$ satisfying $l_0\geq 2$, $m_0\geq 1$. We define
\begin{align*}
\cl{W}_{\hg}^{\mathrm{cl}}&\coloneqq \lbrace y\in \cl{W}_{\hg}^\circ \;\vert\; k_y\geq 2 \rbrace; \\
\cl{W}_{\hh}^{\mathrm{cl}}&\coloneqq \lbrace z\in \cl{W}_{\hh}^\circ \;\vert\; k_z\geq 2\rbrace \cup \lbrace z_0\rbrace.
\end{align*}
Note that, for $\phi\in\lbrace \hg,\hh\rbrace$, the specialization of $\phi$ at a point in $\cl{W}_{\phi}^{\mathrm{cl}}$ is a classical $p$-stabilized newform. For each $x\in \cl{W}_{\phi}^{\mathrm{cl}}$, we define $\alpha_{\phi_x}\coloneqq a_p(\phi_x)$ and $\beta_{\phi_x}\coloneqq\chi_{\phi_x}(p)p^{k_x-1}\alpha_{\phi_x}^{-1}$. Note that, if $k_x\geq 2$, then $\alpha_{\phi_x}$ is the unit root of the $p$-th Hecke polynomial of $\phi_x^{\circ}$.

Let $\cl{W}_{\hg \hh \mathbf{s}}^{\mathrm{cl}}=\cl{W}_{\hg}^{\mathrm{cl}}\times \cl{W}_{\hh}^{\mathrm{cl}}\times \cl{W}^\circ$ and let $\cl{W}_{\hg \hh \mathbf{s}}^{g}$ denote the subset of points $(y,z,s)\in \cl{W}_{\hg \hh \mathbf{s}}^{\mathrm{cl}}$ of weights $(k_y,k_z)=(l,m)$ satisfying $m\leq s<l$. This is the range of interpolation for the three-variable Rankin $p$-adic $L$-function $L_p(\hg,\hh,\mathbf{s})$ introduced below. Similarly, we define $\cl{W}_{\hg\hh\mathbf{s}}^h$ as the subset of points $(y,z,s)\in \cl{W}_{\hg \hh \mathbf{s}}^{\mathrm{cl}}$ of weights $(k_y,k_z)=(l,m)$ satisfying $l\leq s<m$. Note that both $\cl{W}_{\hg\hh\mathbf{s}}^g$ and $\cl{W}_{\hg\hh\mathbf{s}}^h$ exclude points with $l=m$. Set $\Lambda_{\hg \hh \mathbf{s}}=\Lambda_{\hg}\hat{\otimes}_{\cl{O}} \Lambda_{\hh}\hat{\otimes}_{\bb{Z}_p} \Lambda$ and $\cl{O}_{\hg \hh \mathbf{s}}=\Lambda_{\mathbf{hgs}}[1/p]$. Let $I_\hg\subseteq \Lambda_\hg$ be the congruence ideal of $\hg$ and let $I_{\hh}\subseteq \Lambda_{\hh}$ be the congruence ideal of $\hh$.

%The interpolation range of the three-variable $p$-adic $L$-function is given by the following set (which, in particular, does not include any point where $m=l$).

%\begin{defi}
%The {\it critical} range is the set of points $(y,z,\sigma) \in \mathcal W_{\hg \hh}^\circ$ of weights $(l,m,s)$ such that $l,m \geq 2$ and $m \leq s < l$.
%\end{defi}

%Set $\cl{O}_{\hg \hh}:=\cl{O}_{\hg} \hat \otimes_{\cl{O}} \cl{O}_{\hh} \hat \otimes_{\mathbb Z_p} \Lambda$. %and $\mathcal{W}_{\hg \hh}:= \Spf(\Lambda_{\hg \hh})$.  Let $\mathcal Q_{\hg \hh}$ denote the fraction field of $\Lambda_{\hg \hh}$, and set $\mathcal W_{\hg \hh}^\circ:= \mathcal W_{\hg}^\circ \times \mathcal W_{\hh}^\circ \times \mathcal W^{\text{cl}}$.

 Hida constructed in \cite{hida1} and \cite{hida2} a three-variable $p$-adic Rankin $L$-function $L_p(\hg, \hh, \mathbf{s})$ in $I_{\hg}^{-1} \cl{O}_{\mathbf{ghs}}$ interpolating the algebraic parts of the critical values $L(\hg^\circ_y,\hh^\circ_z,s)$ for every triple of critical points $(y,z,s)$ in $\mathcal W_{\hg \hh \mathbf{s}}^{g}$. The following formulation is taken from \cite[Thm.~7.7.2]{KLZ}. % More precisely, Theorem 5.1d of Hida (citar) asserts the following.

\begin{theorem}(Hida)\label{thm:Hida-3var}
There exists a $p$-adic $L$-function
\[
L_p(\hg,\hh,\mathbf{s}) \in I_{\hg}^{-1}\cl{O}_{\hg \hh \mathbf{s}}
\]
such that for every $(y,z,s)\in\mathcal W_{\hg \hh}^g$ of weights $(k_y,k_z)=(l,m)$ with $p\nmid \mathrm{cond}(\hg_y)\cdot \mathrm{cond}(\hh_z)$ we have
\[
L_p(\hg,\hh,\mathbf{s})(y,z,s) = \frac{\mathcal{E}(y,z,s)}{\mathcal{E}_0(y)\mathcal{E}_1(y)} \frac{\Gamma(s)\Gamma(s-m+1)}{\pi^{2s-m+1}(-i)^{l-m}2^{2s+l-m} \langle \hg_y^\circ,\hg_y^\circ \rangle}  \times L(\hg_y^\circ,\hh_z^\circ,s)
\]
%where $C$ is a non-zero algebraic number in the finite extension $\mathbb Q(\hg_y^{\circ},\hh_z^{\circ})$ generated by the Fourier coefficients of $\hg_y^{\circ}$ and $\hh_z^{\circ}$,
where $\langle \hg_y^{\circ}, \hg_y^{\circ} \rangle$ is the Petersson norm as normalized in \cite{KLZ} and the Euler factors are defined by
\begin{eqnarray}
      \nonumber \mathcal E_0(y) &:=& 1-\chi_g^{-1}(p) \beta_{\hg_y}^2p^{1-l}, \\
      \nonumber \mathcal E_1(y) &:=& 1-\chi_g(p)\alpha_{\hg_y}^{-2}p^{l-2}, \\
      \mathcal E(y,z,s) &:=& \Big(1-\frac{p^{s-1}}{\alpha_{\hg_y} \alpha_{\hh_z}} \Big) \Big(1-\frac{p^{s-1}}{\alpha_{\hg_y} \beta_{\hh_z}} \Big) \Big(1-\frac{\beta_{\hg_y} \alpha_{\hh_z}}{p^{s}}\Big) \Big(1-\frac{\beta_{\hg_y} \beta_{\hh_z}}{p^{s}}\Big).
\end{eqnarray}
\end{theorem}

\begin{remark}
Similarly, there is a $p$-adic $L$-function $L_p(\hh,\hg,\mathbf{s}) \in I_{\hh}^{-1}\cl{O}_{\hg \hh \mathbf{s}}$ interpolating the Rankin--Selberg $L$-values in the region $\cl{W}_{\hg\hh\mathbf{s}}^h$. Note that the $p$-adic $L$-functions $L_p(\hg,\hh,\mathbf{s})$ and $L_p(\hh,\hg,\mathbf{s})$ are different. %Further, observe that if $(g,h)$ are two modular forms without $p$ in the level, the $p$-adic $L$-function $L_p(\hg,\hh)$ depends on the choice of a $p$-stabilization for $g$, but not for $h$.
\end{remark}

This construction was later generalized by Urban to the case where the modular forms are not necessarily ordinary \cite{Urban-nearly-overconvergent}, \cite[App.]{AI}, using nearly overconvergent modular forms of finite order and their spectral theory. In this case, the interpolation property does not completely determine the $p$-adic $L$-function, and one needs to impose further conditions.

\subsection{Triple product $p$-adic $L$-function}\label{subsec:tripleproduct}

Let $(\hf,\hg,\hh)$ be a triple of Hida families of tame levels $(N_f,N_g,N_h)$ and characters $(\chi_{\hf},\chi_{\hg},\chi_{\hh})$. In this section, we introduce the three-variable $p$-adic $L$-function interpolating the central values of the triple product $L$-functions associated with classical specializations of $(\hf,\hg,\hh)$ in the $g$-dominant region. %In the ordinary case (i.e., when $(\hf,\hg,\hh)$ is a triple of Hida families) the existence of such a $p$-adic $L$-function with a precise interpolation formula was established by \cite{Hs}. The finite slope case is treated by Andreatta--Iovita \cite{AI} and also in the work of \cite{huang}, although in that case the interpolation property depends on some (undefined) constant; see e.g \cite[Cor. 5.13]{AI} and \cite[Prop. 2.18]{huang}. Therefore, in the following we restrict ourselves to the case when $(\hf,\hg,\hh)$ is a triple of Hida families.

Put $\cl{W}_{\hf \hg \hh}\coloneqq \cl{W}_{\hf}\times \cl{W}_{\hg}\times \cl{W}_{\hh}$. For $\phi\in\lbrace \hf,\hg,\hh\rbrace$, let $\cl{W}_{\phi}^\circ$ denote the subset of crystalline points in $\cl{W}_{\phi}$ and put $\cl{W}_{\hf \hg \hh}^{\circ}=\cl{W}_{\hf}^{\circ} \times \cl{W}_{\hg}^{\circ} \times \cl{W}_{\hh}^{\circ}\subseteq \cl{W}_{\hf \hg \hh}$. Let $w_0=(x_0,y_0,z_0)\in \cl{W}_{\hf \hg \hh}^\circ$ be a good crystalline point with corresponding weights $(k_0,l_0,m_0)$ satisfying $k_{0}+l_{0}+m_{0}\equiv 0\pmod{2}$, $k_0,m_0\geq 1$, $l_0\geq 2$. Put $(f,g,h)\coloneqq (\hf_{x_0}^\circ,\hg_{y_0}^\circ,\hh_{z_0}^\circ)$. Note that $(f,g,h)$ is a triple of $p$-stabilized newforms with characters $(\chi_f,\chi_g,\chi_h)=(\chi_{\hf}^{(p)},\chi_{\hg}^{(p)},\chi_{\hh}^{(p)})$, where $\chi^{(p)}$ stands for the prime-to-$p$ part of the character $\chi$. %Thus $(f,g,h)$ is a triple of newforms of levels $(N_f,N_g,N_h)$, weights $(k_{x_0},k_{y_0},k_{z_0})$ and characters $(\chi_f,\chi_g,\chi_h)=(\chi_{\hf}^{(p)},\chi_{\hg}^{(p)},\chi_{\hh}^{(p)})$, where $\chi^{(p)}$ stands for the prime-to-$p$ part of the character $\chi$.

We make the following assumption.

\begin{ass} 
\hfill
    \begin{enumerate}
        \item $\gcd(N_f,N_g,N_h)$ is square-free;
        \item $\chi_f\chi_g\chi_h=1$;
        \item the residual representation $\bar{\rho}_{g}$ is absolutely irreducible and $p$-distinguished.
    \end{enumerate}
\end{ass}


Put $V_{fgh}\coloneqq V_f\otimes_L V_g\otimes_L V_h$. Then $V_{fgh}$ is a $G_\bb{Q}$-representation of dimension $8$ over $L$. Moreover, putting $c_0\coloneqq(k_{0}+l_{0}+m_0-2)/2$, the twisted representation $V_{fgh}^\dagger\coloneqq V_{fgh}(1-c_0)$ is Kummer self-dual, i.e., $(V_{fgh}^\dagger)^\vee\simeq V_{fgh}^\dagger(1)$.

%\textcolor{red}{Parece que Andreatta--Iovita no dan una f\'ormula de interpolaci\'on precisa, as\'i que igual es mejor restringirse al caso de familias de Hida.} {\color{blue} Si, es cierto que hay una ambiguedad en el factor de interpolacion, tanto en Andreatta--Iovita como luego en Huang cuando hace la ley de reciprocidad... para ciertas cosas entiendo que valdria igual, pero si queremos ser precisos es lo que dices.}

We now add the following assumption.

\begin{ass}
    For all prime $\ell\mid N_fN_gN_h$, $\epsilon_\ell(V_{fgh}^\dagger)=+1$.
\end{ass}


%Put $\bb{V}_{\hf \hg \hh}\coloneqq \bb{V}_{\hf}\otimes_{\cl{O}} \b{V}_{\hg}\otimes_{\cl{O}} \bb{V}_{\hh}$. This module is free of rank $8$ over $\cl{R}\coloneqq\cl{O}_{\hf}\otimes_{\cl{O}} \cl{O}_{\hg}\cl{O}\otimes\cl{O}_{\hh}$. The second condition in the previous assumption implies that $\det(\bb{V}_{\hf \hg \hh})=\cl{X}^2\epsilon_{\cyc}$ for some character $\cl{X}:G_{\bb{Q}}\rightarrow \cl{R}^\times$ unramified at all finite primes away from $p$. Fix a choice of such character $\cl{X}$ and put $\bb{V}_{\hf \hg \hh}^\dagger\coloneqq \bb{V}_{\hf \hg \hh}(\cl{X}^{-1})$. This representation is Kummer self-dual, in the sense that $\bb{V}_{\hf \hg \hh}^\dagger(1)\cong \bb{V}_{\hf \hg \hh}^\dagger$.








%\textcolor{red}{Yo quiz\'a obviar\'ia completamente la cuesti\'on de los test vectors y enunciar\'ia directamente el teorema diciendo que hay una funci\'on L p-\'adica que satisface la f\'ormula de interpolaci\'on que queremos. Sin embargo, s\'i creo que ser\'ia importante recordar las hip\'otesis en Hsieh respecto a local signs y quiz\'a otras cosas, porque el teorema no es cierto sin estas condiciones.} Recall from \cite[Section 1.4]{DR3} the notion of test vector. As proved in Section 3.5 of \emph{loc.\,cit.} following \cite{Hs}, there is a canonical choice of test vectors for which there exists a {\it square-root} $p$-adic $L$-function \[ \Lp^f(\hf, \hg, \hh): \mathcal W_{{\bf fgh}} \rightarrow \mathbb C_p, \] characterized by an interpolation property relating its values at classical points $(x,y,z) \in \mathcal W_{{\bf fgh}}^f$ to the square root of the central critical value of Garrett's triple-product complex $L$-function $L(\hf_x,\hg_y,\hh_z,s)$ associated to the triple of classical eigenforms $(\hf_x,\hg_y,\hh_z)$. For the following proposition, let $\alpha_{\hf_x}$ and $\beta_{\hf_x}$ be the roots of the $p$-th Hecke polynomial of $\hf_x$, ordered in such a way that $\ord_p(\alpha_{\hf_x}) \leq \ord_p(\beta_{\hf_x})$. The following result is \cite[Proposition 5.1]{DR3}.


We define
\begin{align*}
\tilde{\cl{W}}_{\hf}^{\mathrm{cl}}&\coloneqq \lbrace x\in \cl{W}_{\hf}^\circ \;\vert\; k_x\geq 2 \text{ and } k_x\equiv k_0\;(\mathrm{mod}\,2(p-1))\rbrace \cup \lbrace x_0\rbrace; \\ 
\tilde{\cl{W}}_{\hg}^{\mathrm{cl}}&\coloneqq \lbrace y\in \cl{W}_{\hg}^\circ \;\vert\; k_y\geq 2 \text{ and } k_y\equiv l_0 \;(\mathrm{mod}\,2(p-1))\rbrace; \\
\tilde{\cl{W}}_{\hh}^{\mathrm{cl}}&\coloneqq \lbrace z\in \cl{W}_{\hh}^\circ \;\vert\; k_z\geq 2 \text{ and } k_z\equiv m_0 \;(\mathrm{mod}\,2(p-1))\rbrace \cup \lbrace z_0\rbrace.
\end{align*}
Note that, for $\phi\in\lbrace \hf,\hg,\hh\rbrace$, the specialization of $\phi$ at a point in $\tilde{\cl{W}}_{\phi}^{\mathrm{cl}}$ is a classical $p$-stabilized newform. For each $x\in \tilde{\cl{W}}_{\phi}^{\mathrm{cl}}$, we define $\alpha_{\phi_x}\coloneqq a_p(\phi_x)$ and $\beta_{\phi_x}\coloneqq\chi_{\phi_x}(p)p^{k_x-1}\alpha_{\phi_x}^{-1}$. Note that, if $k_x\geq 2$, then $\alpha_{\phi_x}$ is the unit root of the $p$-th Hecke polynomial of $\phi_x^{\circ}$.



Put $\cl{W}_{\hf \hg \hh}^{\mathrm{cl}}=\tilde{\cl{W}_{\hf}}^{\mathrm{cl}}\times \tilde{\cl{W}}_{\hg}^{\mathrm{cl}}\times \tilde{\cl{W}}_{\hh}^{\mathrm{cl}}$. This set admits the natural partition
\[
\mathcal W_{{\bf fgh}}^{\mathrm{cl}} = \mathcal W_{{\bf fgh}}^{f} \sqcup \mathcal W_{{\bf fgh}}^{g} \sqcup \mathcal W_{{\bf fgh}}^{h} \sqcup \mathcal W_{{\bf fgh}}^{\bal},
\]
where
\begin{itemize}
\item $\mathcal W_{{\bf fgh}}^g$ denotes the set of points $(x,y,z) \in \mathcal W_{{\bf fgh}}^{\mathrm{cl}}$ of weights $(k,l,m)$ such that $l \geq k+m$;
\item $\mathcal W_{{\bf fgh}}^f$ (resp. $\mathcal W_{{\bf fgh}}^h$) is defined similarly, replacing the role of $\hg$ by $\hf$ (resp. $\hh$);
\item $\mathcal W_{{\bf fgh}}^{\bal}$ is the set of balanced triples, consisting of points $(x,y,z)\in \mathcal W_{{\bf fgh}}^{\mathrm{cl}}$ of weights $(k,l,m)$ such that each of the weights is strictly smaller than the sum of the other two.
\end{itemize}
Set $\Lambda_{\hf \hg \hh}=\Lambda_{\hf}\hat{\otimes}_{\cl{O}} \Lambda_{\hg}\hat{\otimes}_{\cl{O}} \Lambda_{\hh}$ and $\cl{O}_{\hf \hg \hh}=\Lambda_{\hf \hg \hh}[1/p]$. Let $I_\hg\subseteq \Lambda_\hg$ be the congruence ideal of $\hg$.


%Set $\Lambda_{{\bf fgh}}=\Lambda_{\hf} \hat \otimes \Lambda_{\hg} \hat \otimes \Lambda_{\hh}$ and let $\mathcal W_{{\bf fgh}}^{\circ} := \mathcal W_{\hf}^{\circ} \times \mathcal W_{\hg}^{\circ} \times \mathcal W_{\hh}^{\circ} \subset \mathcal W_{{\bf fgh}} = \Spf(\Lambda_{{\bf fgh}})$ be the set of triples of crystalline classical points, at which the three  families specialize to modular forms with trivial nebentype at $p$. 

In \cite{DR1}, Darmon and Rotger constructed a triple product $p$-adic $L$-function interpolating the square roots of the central values of the triple product $L$-functions associated with classical specializations of $(\hf,\hg,\hh)$ in $\mathcal W_{{\bf fgh}}^g$. The precise interpolation formula given below is due to Hsieh \cite{Hs}

\begin{theorem}(Darmon--Rotger, Hsieh)\label{thm:hsieh}
%Fix test vectors $(\tilde{\hf},\tilde{\hg},\tilde{\hh})$ as in \cite[Section 3]{Hs}. Then $\Lp^f(\tilde{\hf},\tilde{\hg},\tilde{\hh})$ lies in $\Lambda_{{\bf fgh}}$ and for every $(x,y,z) \in \mathcal W_{{\bf fgh}}^f$ of weights $(k,l,m)$ we have \[ \Lp^f(\tilde{\hf},\tilde{\hg},\tilde{\hh})^2(x,y,z) = \frac{\mathfrak a(k,l,m)}{\langle \hf_x^{\circ},\hf_x^{\circ} \rangle^2} \cdot \mathfrak e^2(x,y,z) \times L(\hf_x^{\circ},\hg_y^{\circ},\hh_z^{\circ},c), \] where
There exists a $p$-adic $L$-function
\[
\Lp^g(\hf,\hg,\hh) \in I_\hg^{-1}\cl{O}_{\hf \hg \hh}
\]
such that for every $(x,y,z) \in \mathcal W_{{\bf fgh}}^g$ of weights $(k,l,m)$ with $p\nmid \mathrm{cond}(\hf_x)\cdot\mathrm{cond}(\hg_y)\cdot\mathrm{cond}(\hh_z)$ we have
\[
\Lp^g(\hf,\hg,\hh)^2(x,y,z) = \frac{\mathfrak a(k,l,m)}{\langle \hg_y^{\circ},\hg_y^{\circ} \rangle^2} \cdot \mathfrak e^2(x,y,z) \cdot L(\hf_x^{\circ},\hg_y^{\circ},\hh_z^{\circ},c),
\]
where
\begin{enumerate}
\item $c = \frac{k+l+m-2}{2}$,
\item $\mathfrak a(k,l,m) = (2\pi i)^{-2(l-2)} \cdot \Big( \frac{k+l+m-4}{2} \Big) ! \cdot \Big( \frac{-k+l+m-2}{2} \Big) ! \cdot \Big( \frac{k+l-m-2}{2} \Big) ! \cdot \Big( \frac{-k+l-m}{2} \Big) !$,
\item $\mathfrak e(x,y,z) = \mathcal E(x,y,z)/\mathcal E_0(x) \mathcal E_1(x)$ with
    \begin{eqnarray}
      \nonumber \mathcal E_0(y) &:=& 1-\chi_g^{-1}(p) \beta_{\hg_y}^2p^{1-l}, \\
      \nonumber \mathcal E_1(y) &:=& 1-\chi_g(p)\alpha_{\hg_y}^{-2}p^{l-2}, \\
      \nonumber \mathcal E(x,y,z) &:=& \Big(1-\chi_g(p) \alpha_{\hf_x} \alpha_{\hg_y}^{-1}\alpha_{\hh_z}p^{\frac{-k+l-m}{2}} \Big) \times \Big(1-\chi_g(p) \alpha_{\hf_x} \alpha_{\hg_y}^{-1}\beta_{\hh_z}p^{\frac{-k+l-m}{2}} \Big)\\ \nonumber
       & & \times \Big(1-\chi_g(p) \beta_{\hf_x} \alpha_{\hg_y}^{-1}\alpha_{\hh_z}p^{\frac{-k+l-m}{2}} \Big) \times \Big(1-\chi_g(p) \beta_{\hf_x} \alpha_{\hg_y}^{-1}\beta_{\hh_z}p^{\frac{-k+l-m}{2}} \Big).
    \end{eqnarray}
\end{enumerate}
\end{theorem}

\begin{remark}
    To be precise, each choice of test vectors $(\breve{\hf},\breve{\hg},\breve{\hh})$ for $(\hf,\hg,\hh)$ determines a $p$-adic $L$-function $\Lp^g(\breve{\hf},\breve{\hg},\breve{\hh})$, and Hsieh shows that there exists an optimal choice of test vectors $(\breve{\hf}^\ast,\breve{\hg}^\ast,\breve{\hh}^\ast)$ for which the $p$-adic $L$-function $\Lp^g(\breve{\hf}^\ast,\breve{\hg}^\ast,\breve{\hh}^\ast)$, which we denote simply by $\Lp^g(\hf,\hg,\hh)$, satisfies the precise interpolation formula given above. Using these same test vectors, one can also define $p$-adic $L$-functions $\Lp^f(\hf,\hg,\hh)$ and $\Lp^h(\hf,\hg,\hh)$ interpolating square roots of classical $L$-values in the regions $\mathcal W_{{\bf fgh}}^f$ and $\mathcal W_{{\bf fgh}}^h$, respectively. Note that our choice of test vectors might not be optimal for these regions, so $\Lp^f(\hf,\hg,\hh)$ and $\Lp^h(\hf,\hg,\hh)$ might not precisely satisfy an interpolation formula analogous to the one above. However, if $\gcd(N_g,N_h)=1$, $\hg$ and $\hh$ are residually irreducible and $p$-distinguished, and $\hf$ is the CM Hida family introduced in \S\ref{subsec:wt1}, which is the situation that we will consider, then it follows from \cite{Hs} that one can choose test vectors which are optimal for both $\mathcal W_{{\bf fgh}}^g$ and $\mathcal W_{{\bf fgh}}^h$.
\end{remark}

%Similarly, there exist $p$-adic $L$-functions $\Lp^f$ and $\Lp^h$ interpolating square roots of classical $L$-values in the regions $\mathcal W_{{\bf fgh}}^f$ and $\mathcal W_{{\bf fgh}}^h$, respectively.

%\begin{remark}
%Observe that this is not exactly the interpolation property given by \cite[Thm. 5.7]{AI}, who consider instead a sum of two $p$-adic $L$-functions, one corresponding to each of the $p$-stabilizations of the dominant form.
%\end{remark}

%Since some of the results we develop depend on a suitable choice of test vectors, we fix them once and for all.

\section{Families of cohomology classes}

In this section we introduce the two kinds of cohomology classes that are involved in our comparison, together with the corresponding reciprocity laws connecting them with the $p$-adic $L$-functions introduced in the previous section.

The cohomology classes introduced in this section are the following:
\begin{enumerate}
\item[(a)] the Beilinson--Flach classes attached to a pair of Hida families $(\hg,\hh)$ constructed in \cite{KLZ};% The construction was extended to the case of Coleman families in \cite{LZ-coleman}.
\item[(b)] the diagonal cycles attached to a triple of Hida families $(\hf,\hg,\hh)$ constructed in \cite{BSV} and \cite{DR3}. %The construction was extended to the case of Coleman families in \cite{huang}. We will be interested in the case when $\hf$ is the Hida family introduced in \S\ref{subsec:wt1} passing through the irregular weight-one Eisenstein series $\Eis_1(\varepsilon_K)$.
\end{enumerate}

%We further discuss the behavior of the last Euler system when one of the families is a CM form (passing in this case through an irregular weight one Eisenstein series).

\subsection{Beilinson--Flach classes}\label{subsec:BF}

We use the notations and assumptions introduced in \S\ref{subsec:Hida-Rankin}. In particular, $(\hg,\hh)$ is a pair of Hida families of tame levels $(N_g,N_h)$ and characters $(\chi_\hg,\chi_\hh)$. 
We also make the following assumption.

\begin{ass}
    The Hida families $\hg$ and $\hh$ are residually irreducible and $p$-distinguished.
\end{ass}

Let $\mathbf{s}:\bb{Z}_p^\times\rightarrow \Lambda^\times$ be the character defined by $z\mapsto \omega^{s_0}(z)[\langle z\rangle]$, where $\langle z\rangle =\omega^{-1}(z)z$ and $[\langle z\rangle]$ denotes the corresponding group-like element. Note that the specialization of $\mathbf{s}$ at a point $P_{s,\bm{1}}\in \cl{W}^\circ$ is the character $z\mapsto z^s$. We denote by $\mathbf{l}:\bb{Z}_p^\times\rightarrow \Lambda_\hg^\times$ the weight character of $\hg$, defined by $z\mapsto \chi_{\hg,p}(z) z^2[\langle z\rangle]_{\Lambda_{\hg}}$. Note that the specialization of $\mathbf{l}$ at a crystalline point $y\in \cl{W}_{\hg}^\circ$ of weight $l$ is the character $z\mapsto z^l$. Similarly, we denote by $\mathbf{m}:\bb{Z}_p^\times \rightarrow \Lambda_\hh^\times$ the weight character of $\hh$, defined by $z\mapsto \chi_{\hh,p}(z)z^2[\langle z\rangle]_{\Lambda_{\hh}}$.

%Let $N$ be a positive integer such that $N_g,N_h\mid N$.
Given specializations $\hg_y$ and $\hh_z$, we denote by $\eta_{\hg_y}^\alpha$ and $\omega_{\hh_z}$ the associated differentials as defined in \cite[\S3]{KLZ}. As shown in \cite[\S10]{KLZ}, %in the case of Hida families, and more generally in \cite[\S6]{LZ-coleman} for Coleman families, 
one can define objects $\bm{\eta}_{\hg}$ and $\bm{\omega}_{\hh}$ interpolating these differentials.

  %and let $\underline{\varepsilon}_{\cyc}:G_{\bb{Q}}\rightarrow \Lambda^\times$ denote the $\Lambda$-adic cyclotomic character, i.e., the character defined by $\sigma\mapsto [\varepsilon_{\cyc}(\sigma)]$.  

We first introduce a Perrin-Riou logarithm that will be used for the formulation of an explicit reciprocity law. 

%We refer the reader to \cite[\S10]{KLZ} for the definition of the canonical differentials $\eta_g$ and $\omega_h$ involved in the reciprocity law. In \emph{loc.\,cit.} the authors also establish that it is possible to interpolate the canonical differentials in families, so we have objects $\eta_{\hg}$ and $\omega_{\hh}$.

\begin{propo}\label{perrin1}
There exists a
homomorphism of $\Lambda_{\hg \hh \mathbf{s}}$-modules
\[
\mathcal L_{\hg \hh \mathbf{s}}^{-+}: H^1(\mathbb Q_p, \mathbb V_{\hg}^- \hat \otimes \mathbb V_{\hh}^+ \hat \otimes \Lambda(1-\mathbf{s})) \rightarrow  I_{\hg}^{-1} I^{-1}\cl{O}_{\hg \hh \mathbf{s}},
\]
where $I$ is the augmentation ideal of $\Lambda$, such that for every point $(y,z,s) \in \mathcal W_{{\bf ghs}}^{\mathrm{cl}}$ of weights $(k_y,k_z)=(l,m)$ with $\alpha_{\hg_y}\beta_{\hh_z}\neq p^s$ the specialization of $\mathcal L_{\hg \hh \mathbf{s}}^{-+}$ at $(y,z,s)$ is the homomorphism
\[
\mathcal L_{{\bf ghs}}^{-+}(y,z,s): H^1(\mathbb Q_p, V_{\hg_y}^- \otimes V_{\hh_z}^+(1-s)) \rightarrow \mathbb C_p
\]
given by
\[
\mathcal L_{{\bf ghs}}^{-+}(y,z,s) = \frac{1-p^{s-1} \alpha_{\hg_y}^{-1} \beta_{\hh_z}^{-1}}{1-p^{-s} \alpha_{\hg_y} \beta_{\hh_z}} \times \begin{cases} \frac{(-1)^{m-s-1}}{(m-s-1)!} \times \langle \log_{\BK}(\cdot), \eta_{\hg_y}^\alpha \otimes \omega_{\hh_z} \rangle & \text{ if } s < m \\ (s-m)! \times \langle \exp_{\BK}^*(\cdot), \eta_{\hg_y}^\alpha \otimes \omega_{\hh_z} \rangle & \text{ if } s \geq m, \end{cases}
\]
where $\log_{\BK}$ is the Bloch--Kato logarithm and $\exp_{\BK}^*$ is the dual exponential map.
\end{propo}
\begin{proof}
This follows from \cite[Theorem~8.2.8, Proposition~10.1.1]{KLZ}.
\end{proof}

%\textcolor{blue}{Para Coleman families $\mathbb{V}_{\hg}^-$ no est\'a definido. Igual es mejor restringirse a Hida families para no complicarse. De hecho para Coleman families los ciclos diagonales no se pueden ver como clases variando en la direcci\'on anticiclot\'omica.}

%\textcolor{red}{O: Mmm bueno, $\mathbb V_{\hg}^-$ siempre se puede definir con la teoria de $(\varphi,\Gamma)$-modulos, como en el articulo de David y Sarah del Coleman proceedings (del 16). Pero es cierto que igual es mas comodo restringirse a familias ordinarias si prefieres.}

%\textcolor{blue}{R: Como veas, pero si usamos Coleman families habr\'ia que introducir m\'as notaci\'on, y quiz\'a habr\'ia que usar analytic Iwasawa cohomology para enunciar la proposici\'on anterior.}

%\textcolor{red}{O: si si, tienes razon, que al final poco ganamos. Inicialmente habiamos pensado en hacerlo asi motivados por BSV, que tambien esta para Coleman families. Lo podemos dejar todo para Hida y hacer un comentario sobre que los resultados se extienden de forma natural al caso Coleman.}

%\begin{remark}
%From the general theory of Perrin-Riou maps, the image of $\mathcal L_{\hg \hh}^{-+}$ lands indeed in $I^{-1} \Lambda_{\hg \hh}$, where $I$ is a certain ideal related with the existence of exceptional zeros.
%\end{remark}

\begin{remark}
    We recall that the map in the previous theorem is obtained as the composition
    \[
    \mathcal L_{\hg \hh \mathbf{s}}^{-+} = \langle \log_{\hg \hh \mathbf{s}}^{-+}(\cdot), \bm{\eta}_{\hg} \otimes \bm{\omega}_{\hh} \rangle,
    \]
    where $\log_{\hg \hh \mathbf{s}}^{-+}(\cdot)$ is the Perrin--Riou big logarithm introduced in \cite[Theorem~8.2.8]{KLZ}. %Given a specialization $(g,h)$ of $(\hg,\hh)$, we write $\mathcal L_{gh}^{-+}$ and $\log_{gh}^{-+}$ for the corresponding maps.
\end{remark}

\begin{remark}
    Interchanging the roles of $\hg$ and $\hh$, we also have a map
    \[
    \mathcal L_{\hg \hh\mathbf{s}}^{+-}: H^1(\mathbb Q_p, \mathbb V_{\hg}^+ \hat \otimes \mathbb V_{\hh}^- \hat \otimes \Lambda(1-\mathbf{s})) \rightarrow  I_{\hh}^{-1} I^{-1}\cl{O}_{\hg \hh \mathbf{s}},
    \]
    with analogous interpolation properties.
\end{remark}

We also recall the following result.

\begin{propo}\label{prop:local-property}
    The inclusion $\bb{V}_{\hh}^+\xhookrightarrow{} \bb{V}_{\hh}$ induces an injection
    \[
    H^1(\bb{Q}_p,\bb{V}_{\hg}^-\hat{\otimes}\bb{V}_{\hh}^+\hat{\otimes}\Lambda(1-\mathbf{s}))\xhookrightarrow{}H^1(\bb{Q}_p,\bb{V}_{\hg}^-\hat{\otimes}\bb{V}_{\hh}\hat{\otimes}\Lambda(1-\mathbf{s}))
    \]
\end{propo}
\begin{proof}
    This is part of \cite[Prop.~8.1.7]{KLZ}.
\end{proof}

\begin{remark}
    By the previous proposition, we can regard the module $H^1(\bb{Q}_p,\bb{V}_{\hg}^-\hat{\otimes}\bb{V}_{\hh}^+\hat{\otimes}\Lambda(1-\mathbf{s}))$ as a submodule of $H^1(\bb{Q}_p,\bb{V}_{\hg}^-\hat{\otimes}\bb{V}_{\hh}\hat{\otimes}\Lambda(1-\mathbf{s}))$. Similarly, we can also regard the module $H^1(\bb{Q}_p,\bb{V}_{\hg}^+\hat{\otimes}\bb{V}_{\hh}^-\hat{\otimes}\Lambda(1-\mathbf{s}))$ as a submodule of $H^1(\bb{Q}_p,\bb{V}_{\hg}\hat{\otimes}\bb{V}_{\hh}^-\hat{\otimes}\Lambda(1-\mathbf{s}))$.
\end{remark}

%As a piece of notation, we put $\mathcal L_{\hg \hh}^{-+} = \langle \log_{\hg \hh}^{-+}, \eta_{\hg} \otimes \omega_{\hh} \rangle$, that is, the map we have denoted by $\mathcal L_{\hg \hh}^{-+}$ is the composition of the map interpolating the Bloch--Kato logarithm or the dual exponential followed by the pairing with Ohta's canonical differentials. If we consider a fixed specialization $(g,h)$ of $(\hg,\hh)$, we write $\mathcal L_{gh}^{-+}$ and $\log_{gh}^{-+}$ for the corresponding maps.

To shorten notation, put $\bb{V}_{\hg\hh\mathbf{s}}=\bb{V}_{\hg}\hat{\otimes}\bb{V}_{\hh}\hat{\otimes}\Lambda(1-\mathbf{s})$. Let $\mathscr{F}^2\mathbb V_{\hg\hh\mathbf{s}}\subset \mathbb V_{\hg\hh\mathbf{s}}$ be the $G_{\bb{Q}_p}$-stable $\Lambda_{\hg \hh \mathbf{s}}$-submodule of rank $3$ defined by
\[
\mathscr{F}^2\mathbb V_{\hg\hh\mathbf{s}}=(\bb{V}_{\hg}^+\hat{\otimes}\bb{V}_{\hh} +\bb{V}_{\hg}\hat{\otimes}\bb{V}_{\hh}^+)\hat{\otimes}\Lambda(1-\mathbf{s}).
\]
Fix a finite set $\Sigma$ of places of $\mathbb Q$ containing $\infty$ and the primes dividing $N_gN_hp$ and let $\mathbb Q^\Sigma$ be the maximal extension of $\mathbb Q$ unramified outside $\Sigma$.




\begin{definition}\label{def:Sel1}
The \textit{balanced Selmer group of $\bb{V}_{\hg \hh \mathbf{s}}$} is defined by
\[
H^1_{\bal}(\bb{Q},\mathbb V_{\hg \hh \mathbf{s}}) = \ker \bigg(
H^1(\mathbb Q^\Sigma/\mathbb Q, \mathbb V_{\hg \hh \mathbf{s}}) \longrightarrow
\frac{H^1(\mathbb Q_p,\mathbb V_{\hg \hh \mathbf{s}})}{H^1_{\bal}(\bb{Q}_p,\mathbb V_{\hg \hh \mathbf{s}})}\bigg),
\]
where
\[
H^1_{\bal}(\bb{Q}_p,\mathbb V_{\hg \hh \mathbf{s}})=\ker \big(H^1(\mathbb Q_p, \mathbb V_{\hg \hh \mathbf{s}}) \longrightarrow H^1(\mathbb Q_p, \mathbb V_{\hg \hh \mathbf{s}}/\mathscr{F}^2\mathbb V_{\hg \hh \mathbf{s}}) \big).
\]
\end{definition}

Let $\mathrm{pr}^{-+}:H^1_{\bal}(\bb{Q}_p,\bb{V}_{\hg\hh \mathbf{s}})\rightarrow H^1(\bb{Q}_p,\bb{V}_{\hg}^-\hat{\otimes}\bb{V}_\hh^+\hat{\otimes}\Lambda(1-\mathbf{s}))$ be the map induced by the natural projection $\mathscr{F}^2\bb{V}_{\hg\hh\mathbf{s}}\rightarrow \bb{V}_\hg^-\hat{\otimes}\bb{V}_\hh^+\hat{\otimes}\Lambda(1-\mathbf{s})$ and let
\[
\mathcal{L}_{\hg\hh\mathbf{s}}^{\hg}:H^1_{\bal}(\bb{Q},\bb{V}_{\hg\hh\mathbf{s}})\longrightarrow I_{\hg}^{-1}I^{-1}\cl{O}_{\hg\hh\mathbf{s}}
\]
be the map defined by $\mathcal{L}_{\hg\hh\mathbf{s}}^\hg=\mathcal{L}_{\hg\hh\mathbf{s}}^{-+}\circ \mathrm{pr}^{-+}\circ \res_p$. Similarly we define $\mathcal{L}_{\hg\hh\mathbf{s}}^\hh$.

The construction of Beilinson--Flach classes was first carried out by Lei--Loeffler--Zerbes for fixed modular forms $(g,h)$ of weight two \cite{LLZ0}, and was later extended to Hida families $(\hg,\hh)$ in \cite{KLZ0}. Later, in \cite{LZ-coleman}, the construction was extended to the Coleman case; as discussed e.g. in \cite{LR1} there are situations where one must be more cautious and there may be some poles, for instance, at the critical $p$-stabilization of an Eisenstein series. However, in the setting considered in this note we will not deal with this kind of issues. 

Recall that, given a newform $\xi$ of level $N_{\xi}$, its image under the Atkin--Lehner operator $W_{N_\xi}$ is a scalar multiple of the conjugate eigenform $\xi^*$. Then, we define the {\it Atkin--Lehner pseudo-eigenvalue} $\lambda_{N_\xi}(\xi)$ of $\xi$ by $W_{N_\xi}(\xi) = \lambda_N(\xi)\xi^*$. As explained in \cite[\S10]{KLZ}, given a Hida family $\bm{\xi}$ of tame level $N_\xi$, there exists an element $\lambda_{N_\xi}(\bm{\xi})\in\cl{O}_{\bm{\xi}}^\times$ interpolating the Atkin--Lehner pseudo-eigenvalues of the crystalline specializations of $\bm{\xi}$.

%In this scenario, we formulate the main result of \cite{KLZ}, which asserts that there exists a family of cohomology classes indexed by points of $\mathcal W_{\hg \hh }$ and whose image under the previous Perrin-Riou map agrees with the Hida--Rankin $p$-adic $L$-function.

We now state the main result of \cite{KLZ}.

\begin{theorem}\label{reclaw}
Fix an integer $c>1$ relatively prime to $6pN_gN_h$. Then, there exists a global cohomology class \[ _c\kappa(\hg,\hh,\mathbf{s}) \in H^1_{\bal}(\mathbb Q, \mathbb V_{\hg} \hat \otimes \mathbb V_{\hh} \hat \otimes \Lambda(1-\mathbf{s})) \] such that
\[
\cl{L}_{\hg\hh\mathbf{s}}^{\hg}(_c\kappa(\hg,\hh,\mathbf{s}))=\frac{(-1)^{\mathbf{s}}}{\lambda_{N_g}(\hg)} \cdot (c^2-c^{2\mathbf{s}-\mathbf{l}-\mathbf{m}+2}\chi_{\hg}^{(p)}(c)^{-1}\chi_{\hh}^{(p)}(c)^{-1}) \times L_p(\hg,\hh,\mathbf{s})
\]
and
\[
\cl{L}_{\hg\hh\mathbf{s}}^{\hh}(_c\kappa(\hg,\hh,\mathbf{s}))=\frac{(-1)^{\mathbf{s}}}{\lambda_{N_h}(\hh)} \cdot (c^2-c^{2\mathbf{s}-\mathbf{l}-\mathbf{m}+2}\chi_{\hg}^{(p)}(c)^{-1}\chi_{\hh}^{(p)}(c)^{-1}) \times L_p(\hh,\hg,\mathbf{s}).
\]
%\begin{enumerate}
%\item The image of the local class $\res_p(_c \kappa(\hg,\hh,\mathbf{s}))$ in $H^1(\mathbb Q_p, \mathbb V_{\hg}^- \hat \otimes \mathbb V_{\hh}\hat{\otimes}\Lambda(1-\mathbf{s}))$ lands in \[ H^1(\mathbb Q_p, \mathbb V_{\hg}^- \hat \otimes \mathbb V_{\hh}^+\hat{\otimes}\Lambda(1-\mathbf{s})). \]
%\item Letting $_c \kappa_p^{-+}(\hg,\hh)$ denote the local cohomology class in the above space, we have \[ \mathcal L_{{\bf gh}}^{-+}(_c \kappa_p^{-+}(\hg,\hh,\mathbf{s})) = \frac{(-1)^{\mathbf{s}}}{\lambda_{N_g}(\hg)} \cdot (c^2-c^{2\mathbf{s}-\mathbf{l}-\mathbf{m}+2}\chi_{\hg}^{-1}(c)\chi_{\hh}^{-1}(c)) \times L_p(\hg,\hh,\mathbf{s}). \] %where $\lambda_{\hg}$ denotes the pseudo-eigenvalue of $\hg$, an Iwasawa function interpolating the pseudo-eigenvalue at $N$ of the crystalline classical specializations of $\hg$.
%\end{enumerate}
\end{theorem}
\begin{proof}
The global cohomology class $_c \kappa(\hg,\hh,\mathbf{s})$ is introduced in \cite[Definition 8.1.1]{KLZ}. The result follows from \cite[Proposition 8.1.7]{KLZ} and \cite[Theorem~B]{KLZ}.
\end{proof}

%We have used the variable $s$ instead of the most common $j$; the relation is just $s = j+1$. Towards working with Euler systems, since $j$ may be taken as the {\it cyclotomic-twist} variable, it is convenient to keep this relation at hand.

\begin{remark}\label{rk:dependenceofBFonc}
After tensoring with $\operatorname{Frac} (\mathcal{O}_{\mathbf{ghs}})$, the class
\[
\kappa(\hg, \hh,\mathbf{s}) \coloneqq C_c^{-1} \otimes {}_c \kappa(\hg, \hh,\mathbf{s})
\]
is independent of $c$, where
\begin{equation}
 \label{eq:defCd}
  C_c(\hg,\hh,{\bf s}) \coloneqq c^2 - c^{(2{\bf s}-{\bf l}-{\bf m}+2)}\chi_\hg^{(p)}(c)^{-1} \chi_\hh^{(p)}(c)^{-1}.
\end{equation}
%However, since $c$ is fixed throughout, we may sometimes drop it from the notation and just write $\kappa_{\hg,\hh}$ for the class $_c \kappa(\hg,\hh,\mathbf{s})$. The constant does make an appearance in fudge factors accounting for the interpolation properties satisfied by the Euler system, but in the case we are interested in these fudge factors do not vanish. We typically refer to this class as the {\it Beilinson--Flach class} or the {\it Beilinson--Flach element attached to $\hg$ and $\hh$}.  
\end{remark}

%\begin{remark}
 %   Theorem~\ref{reclaw}(1) is equivalent to the statement that the local class $_c\kappa(\hg,\hh,\mathbf{s})$ maps to zero under the map induced in cohomology by the projection $V_\hg \otimes V_\hh \rightarrow V_\hg^- \otimes V_\hh^-$.
%\end{remark}



\subsection{Diagonal cycles}\label{subsec:diagonalcycles}

In this section, we recall the main results regarding the existence of $p$-adic families of diagonal cycles obtained in \cite{DR3} and \cite{BSV2} building on the previous geometric constructions of \cite{DR2}.

We use the notations and assumptions introduced in \S\ref{subsec:tripleproduct}. In particular, $(\hf,\hg,\hh)$ is a triple of Hida families of tame levels $(N_\hf,N_\hg,N_\hg)$ and characters $(\chi_\hf,\chi_\hg,\chi_\hh)$ and $(f,g,h)=(\hf_{x_0}^\circ,\hg_{y_0}^\circ,\hh_{z_0}^\circ)$ is a triple of good crystalline specializations of characters $(\chi_f,\chi_g,\chi_h)$ and weights $(k_0,l_0,m_0)$ with $k_{0}+l_{0}+m_{0}\equiv 0\pmod{2}$. %, living over discs $V_1$, $V_2$, $V_3$ in weight space, with all classical specialisations non-critical-slope cusp forms, as in the previous section.
We denote by $\mathbf{k}:\bb{Z}_p^\times \rightarrow \Lambda_\hf^\times$, $\mathbf{l}:\bb{Z}_p^\times\rightarrow \Lambda_\hg^\times$ and $\mathbf{m}:\bb{Z}_p^\times \rightarrow \Lambda_\hh^\times$ the corresponding tautological weight characters.

The running assumptions imply that $\chi_{\hf} \chi_{\hg} \chi_{\hh} = \omega^{2r}$ for some $r\in \mathbb{Z}$, where $\omega$ denotes the Teichm\"uller character. Thus we can choose a character ${\bf t}: \mathbb Z_p^\times \to \Lambda_{\hf\hg\hh}^\times=(\Lambda_\hf \hat{\otimes} \Lambda_\hg \hat{\otimes}\Lambda_\hh)^\times$ satisfying $2{\bf t} = {\bf k} + {\bf l} + {\bf m}$. There are two choices for $\mathbf{t}$, and we choose the one determined by imposing that the specialization of $\mathbf{t}$ at the fixed crystalline point $(x_0,y_0,z_0)$ is the character $\varepsilon_{\cyc}^{(k_0+l_0+m_0)/2}$. % This imposes an additional condition on our specialisations: we shall say a point $P$ of $V_1 \times V_2 \times V_3$ is an ``integer point'' if $({\bf k}, {\bf l}, {\bf m})$ specialise at $P$ to integers $(k,l, m) \ge -1$, and, in addition, ${\bf t}$ specialises to $\frac{k + l + m}{2}$ (rather than its product with the quadratic character mod $p$), which amounts to requiring that $k + l + m$ lie in a particular congruence class modulo $2(p-1)$.
We also make the following assumption.

\begin{ass}
\hfill
    \begin{enumerate}[(i)]
        \item The Hida families $\hg$ and $\hh$ are residually irreducible and $p$-distinguished.
        \item $\bb{V}_{\hf}$ and $\bb{V}_{\hf}^-$ are free $\Lambda_\hf$-modules.
    \end{enumerate}
\end{ass}

Note that, by Proposition~\ref{prop:BSTW}, the second part of the assumption holds when $\hf$ is the CM Hida family introduced in \S\ref{subsec:wt1}, which is the case that we will consider in later sections.


%Following a construction of Andreatta and Iovita \cite{AI}, there is a $\hf$-dominant square root $p$-adic $L$-function $\Lp^{\hf}(\hf,\hg,\hh)$ over $V_1 \times V_2 \times V_3$.  This extends the construction of Hsieh \cite{Hs} in the ordinary case. The square of its value at an integer point $(k, l, m)$ with $k > l+m$ is $(\star) \cdot L(f_k, g_l, h_m, \frac{k+l+m}{2}+2)$, where $(\star)$ is the usual m\'elange of Euler factors, complex periods etc. (Note that with our conventions $f_k$ has weight $k$, etc, so $\frac{k+l+m}{2}-1$ is the center of the functional equation.  %\textcolor{blue}{Esta convenci\'on parece un poco problem\'atica. En otras secciones estamos usando $k$ para denotar el peso de una especializaci\'on de $\hf$.})

We first introduce a Perrin-Riou logarithm that will be used for the formulation of an explicit reciprocity law. 

\begin{propo}\label{perrin2}
There exists a
%injective (no veo claro que sea siempre inyectivo)
homomorphism of $\Lambda_{\hf \hg \hh}$-modules
\[
\mathcal L_{\hf \hg \hh}^{+-+}: H^1(\mathbb Q_p, \bb{V}_\hf^+ \hat{\otimes}\mathbb V_{\hg}^- \hat \otimes \mathbb V_{\hh}^+ (2-\mathbf{t})) \rightarrow I_\hg^{-1} \cl{O}_{\hf \hg \hh },
\]
such that for every point $(x,y,z) \in \mathcal W_{{\bf fgh}}^{\mathrm{cl}}$ of weights $(k,l,m)$ with $\beta_{\hf_x}\alpha_{\hg_y}\beta_{\hh_z}\neq p^{(k+l+m-2)/2}$ the specialization of $\mathcal L_{\hf \hg \hh}^{+-+}$ at $(x,y,z)$ is the homomorphism
\[
\mathcal L_{{\bf fgh}}^{+-+}(x,y,z): H^1(\mathbb Q_p, V_{\hf_x}^+\otimes V_{\hg_y}^- \otimes V_{\hh_z}^+(1-c)) \rightarrow \mathbb C_p
\]
given by
\[
\mathcal L_{{\bf fgh}}^{+-+}(x,y,z) = \frac{1-p^{-c}\alpha_{\hf_x} \beta_{\hg_y}\alpha_{\hh_z}}{1-p^{-c} \beta_{\hf_x} \alpha_{\hg_y} \beta_{\hh_z}} \times \begin{cases} \frac{(-1)^{c-l}}{(c-l)!} \times \langle \log_{\BK}(\cdot), \omega_{\hf_x}\otimes\eta_{\hg_y}^\alpha \otimes \omega_{\hh_z} \rangle & \text{ if } l < k+ m \\ (l-c-1)! \times \langle \exp_{\BK}^*(\cdot), \omega_{\hf_x} \otimes\eta_{\hg_y}^\alpha \otimes \omega_{\hh_z} \rangle & \text{ if } l \geq k+m, \end{cases}
\]
where $\log_{\BK}$ denotes the Bloch-Kato logarithm, $\exp_{\BK}^*$ denotes the dual exponential map and $c=(k+l+m-2)/2$.
\end{propo}
\begin{proof}
This follows from \cite[Proposition~7.3]{BSV}.
\end{proof}

%As a piece of notation, we put $\mathcal L_{\hf \hg \hh}^{-++} = \langle \log_{\hf \hg \hh}^{-++}, \eta_{\hf} \otimes \omega_{\hg} \otimes \omega_{\hh}\rangle$, that is, the map we have denoted by $\mathcal L_{\hf \hg \hh}^{-++}$ is the composition of the map interpolating the Bloch--Kato logarithm or the dual exponential followed by the pairing with Ohta's canonical differentials.

\begin{remark}
    We recall that the map in the previous theorem is obtained as the composition
    \[
    \mathcal L_{\hg \hh}^{+-+} = \langle \log_{\hf \hg \hh}^{+-+}(\cdot), \bm{\omega}_{\hf}\otimes\bm{\eta}_{\hg} \otimes \bm{\omega}_{\hh} \rangle,
    \]
    where $\log_{\hf \hg \hh}^{+-+}(\cdot)$ is the Perrin--Riou big logarithm introduced in \cite[Proposition~7.1]{BSV}. %Given a specialization $(f,g,h)$ of $(\hf,\hg,\hh)$, we write $\mathcal L_{fgh}^{+-+}$ and $\log_{fgh}^{+-+}$ for the corresponding maps.
\end{remark}

\begin{remark}
    Letting $\hf$ (resp. $\hh$) play the role of $\hg$ in Proposition~\ref{perrin2}, we obtain a similar map $\cl{L}_{\hf\hg\hh}^{-++}$ (resp. $\cl{L}_{\hf\hg\hh}^{++-}$) with analogous properties.
\end{remark}


To shorten notation, put $\bb{V}_{\hf\hg\hh}^\dagger=\bb{V}_\hf\hat{\otimes}\bb{V}_\hg\hat{\otimes}\bb{V}_\hh(2-\mathbf{t})$. Let $\mathscr{F}^2\mathbb V_{\hf\hg\hh}^\dag\subset \mathbb V_{\hf\hg\hh}^{\dag}$ be the $G_{\bb{Q}_p}$-stable $\Lambda_{\hf \hg \hh}$-submodule of rank $4$ defined by
\[
\mathscr{F}^2\mathbb V_{\hf\hg\hh}^\dag=(\bb{V}_\hf\hat{\otimes}\bb{V}_{\hg}^+\hat{\otimes}\bb{V}_{\hh}^+ +\bb{V}_\hf^+\hat{\otimes}\bb{V}_{\hg}\hat{\otimes}\bb{V}_{\hh}^+ +\bb{V}_\hf^+\hat{\otimes}\bb{V}_{\hg}^+\hat{\otimes}\bb{V}_{\hh})(2-\mathbf{t}).
\]
Fix a finite set $\Sigma$ of places of $\mathbb Q$ containing $\infty$ and the primes dividing $N_fN_gN_hp$ and let $\mathbb Q^\Sigma$ be the maximal extension of $\mathbb Q$ unramified outside $\Sigma$.

%There is a $G_{\mathbb Q_p}$-stable rank-four $\Lambda_{\hf\hg\hh}$-submodule $\mathcal F^2\mathbb V_{\hf\hg\hh}^\dag\subset \mathbb V_{\hf\hg\hh}^{\dag}$, so we may write
%\[
%\mathbb V_{\hf\hg\hh}^f:=\mathbb V_{\hf}^+\hat\otimes_{\mathcal{O}}\mathbb V_{\hg}\hat\otimes_{\mathcal O}\mathbb V_{\hh}(2-{\bf t}).
%\]
%As before, we let $\mathcal F^2\mathbb V_{\hf gh}$ and $\mathbb V_{\hf gh}^f$ denote the corresponding specializations.

%\begin{remark}
%When the forms are no longer ordinary, one must work instead with the filtrations coming from $(\varphi,\Gamma)$-modules, but the local behavior of the classes is completely analogous.
%\end{remark}



%We can now formulate the result which asserts the existence of {\it big diagonal cycles}, together with an explicit reciprocity law linking them with the $p$-adic $L$-function of Andreatta--Iovita.


%Me parece que solo nos hace falta el balanced.




\begin{definition}\label{def:Sel2}
The \textit{balanced Selmer group of $\bb{V}_{\hf \hg \hh}^\dagger$} is defined by
\[
H^1_{\bal}(\bb{Q},\mathbb V_{\hf \hg \hh}^{\dag}) = \ker \bigg(
H^1(\mathbb Q^\Sigma/\mathbb Q, \mathbb V_{\hf \hg \hh}^{\dag}) \longrightarrow
\frac{H^1(\mathbb Q_p,\mathbb V_{\hf \hg \hh}^\dag)}{H^1_{\bal}(\bb{Q}_p,\mathbb V_{\hf \hg \hh}^\dag)}\bigg),
\]
where
\[
H^1_{\bal}(\bb{Q}_p,\mathbb V_{\hf \hg \hh}^\dag)=\ker \big(H^1(\mathbb Q_p, \mathbb V_{\hf \hg \hh}^{\dag}) \longrightarrow H^1(\mathbb Q_p, \mathbb V_{\hf \hg \hh}^{\dag}/\mathscr{F}^2\mathbb V_{\hf \hg \hh}^{\dag}) \big).
\]
\end{definition}

Let $\mathrm{pr}^{+-+}:H^1_{\bal}(\bb{Q}_p,\bb{V}_{\hf\hg\hh}^\dagger)\rightarrow H^1(\bb{Q}_p,\bb{V}_{\hf}^+\hat{\otimes}\bb{V}_{\hg}^-\hat{\otimes}\bb{V}_\hh^+(2-\mathbf{t}))$ be the map induced by the natural projection $\mathscr{F}^2\bb{V}_{\hf\hg\hh}^\dagger\rightarrow \bb{V}_{\hf}^+\hat{\otimes}\bb{V}_\hg^-\hat{\otimes}\bb{V}_\hh^+(2-\mathbf{t})$ and let
\[
\mathcal{L}_{\hf\hg\hh}^{\hg}:H^1_{\bal}(\bb{Q},\bb{V}_{\hf\hg\hh}^\dagger)\longrightarrow I_{\hg}^{-1}\cl{O}_{\hf\hg\hh}
\]
be the map defined by $\mathcal{L}_{\hf\hg\hh}^\hg=\mathcal{L}_{\hf\hg\hh}^{+-+}\circ \mathrm{pr}^{+-+}\circ \res_p$. Similarly we define $\mathcal{L}_{\hf\hg\hh}^\hf$ and $\mathcal{L}_{\hf\hg\hh}^\hh$. 



\begin{comment}
\begin{definition}\label{def:Sel}
For $\mathcal{L}\in\{\bal,\mathcal{F}\}$ define the Selmer group $\Sel_{\mathcal{L}}(\mathbb V_{\hf gh}^{\dag})$ by
\[
\Sel_{\mathcal{L}}(\mathbb V_{\hf gh}^{\dag}) = \ker \bigg(
H^1(\mathbb Q^\Sigma/\mathbb Q, \mathbb V_{\hf gh}^{\dag}) \longrightarrow
\frac{H^1(\mathbb Q_p,\mathbb V_{\hf gh}^\dag)}{H^1_{\mathcal{L}}(\mathbb Q_p,\mathbb V_{\hf gh}^\dag)}\bigg),
\]
where
\[
H_{\mathcal L}^1(\mathbb Q_p, \mathbb V_{\hf gh}^{\dag}) =
\begin{cases}
\ker \big(H^1(\mathbb Q_p, \mathbb V_{\hf gh}^{\dag}) \longrightarrow H^1(\mathbb Q_p, \mathbb V_{\hf gh}^{\dag}/\mathscr{F}^2\mathbb V_{\hf gh}^{\dag}) \big) &\textrm{if $\mathcal{L}={\rm bal}$},\\[0.5em]
\ker \big(H^1(\mathbb Q_p, \mathbb V_{\hf gh}^{\dag}) \longrightarrow H^1(\mathbb Q_p^{}, \mathbb V_{\hf gh}^{\dag}/ \mathbb V_{\hf gh}^{f}) \big) &\textrm{if $\mathcal{L}=\mathcal{F}$}.
\end{cases}
\]
We call $\Sel_{\bal}(\mathbb V_{\hf gh}^{\dag})$ (resp. $\Sel_{\mathcal F}(\mathbb V_{\hf gh}^{\dag})$)  the \emph{balanced} (resp. \emph{$f$-unbalanced}) Selmer group.
\end{definition}
\end{comment}

We now state main result of \cite{DR3} and \cite{BSV}.

\begin{theorem}\label{thm:reclawdiagonalcycles}
There exists a global cohomology class
\[
\kappa(\hf,\hg,\hh) \in H^1_{\bal}(\mathbb Q, \mathbb V_{\hf} \hat \otimes \mathbb V_{\hg} \hat \otimes \mathbb V_{\hh}(2-{\bf t}))
\]
such that, for $\bm{\xi}\in \lbrace \hf,\hg,\hh\rbrace$,
we have that
\[
\mathcal L_{\hf\hg\hh}^{\bm{\xi}}(\kappa(\hf, \hg, \hh)) = \Lp^{\xi}(\hf, \hg, \hh).
\]
%\begin{enumerate}
%\item The projection of the local class $\res_p(\kappa(\hf,\hg,\hh))$ to $H^1(\mathbb Q_p, \mathbb V_{\hf}^- \hat \otimes \mathbb V_{\hg} \hat \otimes \mathbb V_{\hh}(1-{\bf c}))$ lands in $H^1(\mathbb Q_p, \mathbb V_{\hf}^- \hat \otimes \mathbb V_{\hg}^+ \hat \otimes \mathbb V_{\hh}^+(1-{\bf c}))$.
%\item Letting $\kappa_p^{-++}(\hf,\hg,\hh)$ denote the local cohomology class in the above space, \[ \mathcal L_{\hf \hg \hh}^{-++}(\kappa_p^{-++} (\hf, \hg, \hh)) = \Lp^f(\hf, \hg, \hh). \]
%\end{enumerate}

\end{theorem}

\begin{proof}
This is \cite[Theorem~A]{BSV} or \cite[Theorem~5.1]{DR3}.
%, specifically considered in the general case by \cite{huang}.

%No veo por qu\'e no saldr\'ia directamente de BSV. Comento lo de debajo.
%Alternatively, one may follow the approach of \cite[Prop. 2.3]{BDV} to establish that we can still use the techniques of \cite{BSV} to construct a family of cohomology classes when one of the (cuspidal) families passes through a irregular weight one Eisenstein point.
\end{proof}

%Although the result in \emph{loc.\,cit.} is stated just for Hida families, the argument remains valid in the present generality, using the $p$-adic $L$-function of Andreatta--Iovita and the general theory of $(\varphi,\Gamma)$-modules.


\subsection{Anticyclotomic diagonal cycles}\label{subsec:ac-cycles}

\begin{comment}
\textcolor{blue}{Esto no parece muy relevante para el art\'iculo, y de hecho ni siquiera se puede aplicar. No s\'e si conviene mantenerlo.}

\textcolor{blue}{O: Yo lo veía relevante para luego interpretarlo en cierto sentido como comparar un sistema ciclotomico (BF) con uno anticiclotomico (el que se obtiene a partir de los ciclos diagonales), aunque solo sea en el bottom layer.}

\textcolor{blue}{Pero es que en realidad en este caso los resultados de ACR no se pueden aplicar, as\'i que no est\'a claro que tengamos un sistema de Euler anticiclot\'omico. Y en todo caso el sistema de Euler ciclot\'omico de momento tampoco lo hemos introducido, solo hemos hablado del bottom layer tambi\'en en el caso BF. Yo creo que igual pod\'iamos restringirnos al bottom layer tambi\'en en el caso anticiclot\'omico, sin hablar de las tame norm relations}

\textcolor{blue}{O: vale, perfecto, si si, tienes razon, me parece razonable.}
\end{comment}

We keep the notations and assumptions in the previous subsection, as well as the notations and assumptions in \S\ref{subsec:wt1}, and specialize the discussion in the previous subsection to the case in which $\hf$ is the Hida family introduced in \S\ref{subsec:wt1} with $f=\hf_{x_0}^\circ=\Eis_1(\varepsilon_K)$.

In this case, we have a class
\[
\kappa(\hf,\hg,\hh)\in H^1(\mathbb Q, \Ind_K^{\mathbb Q} \Lambda_{\hf}(\bm{\varphi}) \hat \otimes \mathbb V_{\hg} \hat \otimes \mathbb V_{\hh}(2-{\bf t}))
\]
which under Shapiro's isomorphism can also be seen as a class in
\[
H^1(K,\Lambda_\hf(\bm{\varphi})\hat{\otimes}\mathbb V_{\hg} \hat \otimes \mathbb V_{\hh}(2-\bf{t})).
\]
Moreover, after specializing $\hf$ to $f$, we obtain a class
\begin{align*}
    \kappa(f,\hg,\hh)\in H^1(\bb{Q}, (1\oplus\varepsilon_K)\otimes \bb{V}_\hg \hat{\otimes} \bb{V}_\hh (2-\mathbf{t}_1))\cong H^1(K,\bb{V}_\hg\hat{\otimes} \bb{V}_\hh(2-\bf{t}_1))
\end{align*}
where $\bf{t}_1$ is the composition of $\bf{t}$ with the map $(\Lambda_{\hf}\hat{\otimes}\Lambda_{\hg}\hat{\otimes}\Lambda_\hh)^\times\rightarrow(\Lambda_{\hg}\hat{\otimes}\Lambda_\hh)^\times$ induced by the point $x_0:\Lambda_\hf\rightarrow \cl{O}$.


Assume now in addition that $p$ does not divide the class number of $K$. Then the classes $\kappa(\hf,\hg,\hh)$ yield classes
\[
\tilde{\kappa}(\hf,\hg,\hh)\in H^1(K,\Lambda_{\ac}(\bm{\kappa}_{\ac}^{-1})\hat{\otimes} \bb{V}_\hg \hat{\otimes} \bb{V}_\hh (2-\mathbf{t}_1)) \cong H^1_{\mathrm{Iw}}(K_\infty^{\ac},\bb{V}_\hg\hat{\otimes}\bb{V}_{\hh}(2-\mathbf{t}_1)),
\]
where $K_\infty^{\ac}$ denotes the anticyclotomic $\bb{Z}_p$-extension of $K$, $\Lambda_{\ac}=\bb{Z}_p[[\Gal(K_\infty^{\ac}/K)]]$ denotes the corresponding Iwasawa algebra and $\bm{\kappa}_{\ac}:G_K\rightarrow \Lambda_{\ac}^\times$ denotes the tautological character. When $\hf$ is residually non-Eisenstein, \cite{ACR} realises the class $\tilde{\kappa}(\hf,\hg,\hh)$ as the bottom layer of an anticyclotomic Euler system. However, this does not apply to the choice of $\hf$ in this article. We intend to come back to this in the future.


%Let $g \in S_l(N_g,\chi_g)$ and $h \in S_m(N_h,\chi_h)$ be newforms of weights $l \geq m \geq 2$ of the same parity. Let $K/\mathbb Q$ be an imaginary quadratic field of discriminant $-D<0$. Let $k>0$ be an odd integer, and let $\psi$ be a Hecke character of infinity type $(1-k,0)$, conductor $\mathfrak f$, and central character $\varepsilon_{\psi} = \bar \chi_g \bar \chi_h$. Let $p \nmid N_g N_h$ be an odd prime with $(\mathfrak f,p)=1$, and an mebedding $\iota_p \colon \overline{\mathbb Q} \hookrightarrow \overline{\mathbb Q}_p$, and let $E = L_{\mathfrak P}$ be a finite extension of $\mathbb Q_p$ containing the image under $\iota_p$ of the values of $\psi$ and the Fourier coefficients of both $g$ and $h$. Let \[ V_{g,h}^{\psi} := V_g \otimes V_h(\psi_{\mathfrak P}^{-1})(1-c), \] where $c = (k+l+m-2)/2$. We also assume that $p$ splits in $K$ and $p \nmid h_K$, where $h_K$ is the class number of $K$.

%In that case, proceeding as in \cite{ACR}, one can construct a cohomology class \[ \kappa_{\psi,g,h} \in H^1(K, T_{g,h}^{\psi}) \] corresponding to taking a CM form in the triple product case and applying Shapiro's lemma in a suitable way. In \cite{ACR}, when we work under the {\it non-Eisenstein assumption}, this is indeed the bottom layer of an anticyclotomic Euler system.

\begin{comment}

\begin{theorem}
Suppose that $p$ splits in $K$ and $p \nmid h_K$, and that both $g$ and $h$ are ordinary at $p$. Let $\mathcal S$ be the set of all squarefree products of primes split $q$ in $K$ and coprime to $DN_gN_h \mathfrak f$ and denote by $K[n]$ the maximal $p$-extensions of $K$ inside the ring class field of conductor $n$. If the character $\psi$ is non-Eisenstein (in the terminology of \cite{LLZ}), there exists a family of cohomology classes \[ \kappa_{\psi,g,h,np^r} \in H^1(K[np^r], T_{g,h}^{\psi}) \] for all $n \in \mathcal S$ and $r \geq 0$, where $T_{g,h}^{\psi}$ is a fixed $G_K$-stable lattice inside $V_{g,h}^{\psi}$, satisfying the Euler system compatibility relations detailed in \cite[\S4, 6]{ACR}.
\end{theorem}

\begin{proof}
This is \cite[Thm. A]{ACR}. Note that there we had assumed that $l$ and $m$ had the same parity, and we had then taken an even integer $k$. Note that exactly the same proof works when $g$ and $h$ are modular forms with weights of different parity and $k$ is odd, since the only requirement is that $k+l+m$ is even.
\end{proof}

\begin{defi}
Let $\kappa_{\psi,g,h} := \kappa_{\psi,g,h,1} \in H^1(K,T_{g,h}^{\psi})$ denote the bottom layer of the previous anticyclotomic Euler system.
\end{defi}

\end{comment}

%In the cases where $\psi$ is Eisenstein, with the current methods we cannot claim that such an Euler system exists. However, the bottom layer $\kappa_{\psi,g,h}$ is still available. Furthermore, the class $\kappa_{\psi,g,h}$ satisfies an explicit reciprocity law connecting it with the triple product $p$-adic $L$-function; see e.g. \cite[Thm. 7.4]{ACR}. We expect to be able to interpret $\kappa_{\psi,g,h}$ as the bottom layer of an anticyclotomic Euler system in forthcoming work.



\section{Comparison of classes}\label{sec:comparison}


We keep the assumptions and notations introduced in previous sections. In particular, $K$ is an imaginary quadratic field in which $p$ splits and $\hf$ is the CM Hida family introduced in \S\ref{subsec:wt1} with $f=\hf_{x_0}^\circ=\Eis_1(\varepsilon_K)$. Also, $(\hg,\hh)$ is a pair of Hida families of tame levels $(N_g,N_h)$ and characters $(\chi_\hg,\chi_\hh)$ such that $\chi_\hg \chi_\hh =\varepsilon_K \omega^{2r}$ for some $r\in \bb{Z}$, and $g=\hg_{y_0}^\circ$ and $h=\hh_{z_0}^\circ$ are good crystalline specializations of $\hg$ and $\hh$ of weights $l_0\geq 2$ and $m_0\geq 1$, respectively. They have characters $(\chi_g,\chi_h)=(\chi_{\hg}^{(p)},\chi_\hh^{(p)})$. We define $\alpha_g=a_p(\hg_{y_0})$, $\beta_g=\chi_g(p)p^{l_0-1}\alpha_g^{-1}$, $\alpha_h=a_p(\hh_{z_0})$, $\beta_h=\chi_h(p)p^{m_0-1}\alpha_h^{-1}$. The corresponding definitions for $f$ yield $\alpha_f=\beta_f=1$. Set $c_0=s_0=(l_0+m_0-1)/2$. The following set of assumptions will be in place in all this section.

\begin{ass}
\hfill
\begin{enumerate}[(i)]
    \item $\gcd(N_g,N_h)=1$;
    \item $\hg$ and $\hh$ are residually irreducible and $p$-distinguished;
    \item $h$ is a newform of level $N_h$ (i.e., $h$ has prime-to-$p$ conductor).
\end{enumerate}
\end{ass}



%With the assumptions fixed in the introduction, we consider the modular forms $g \in S_l(N_g,\chi_g)$ and $h \in S_m(N_h,\chi_h)$, together with the imaginary quadratic field $K/\mathbb Q$. We also fix a prime $p$ which splits in $K$. We prove the main theorem of the paper under the assumption that there are no exceptional zeros arising in the denominator of the Perrin-Riou maps of $L_p(g,h)$ and $L_p(h,g)$ (we need both of them). We write $(\alpha_g,\beta_g)$ and $(\alpha_h,\beta_h)$ for the roots of the $p$-th Hecke polynomials of $g$ and $h$ at $p$, respectively.
%and ordered in such a way that $\ord_p(\alpha_g) \leq \ord_p(\beta_g)$ and $\ord_p(\alpha_h)<\ord_p(\beta_h)$. 
%Hence, our main assumption reads as follows.



\subsection{An equality of $p$-adic $L$-functions}

Set $\Lambda_{\hg\hh}=\Lambda_{\hg}\hat{\otimes}\Lambda_{\hh}$ and $\cl{O}_{\hg \hh}=\Lambda_{\hg\hh}[1/p]$. Let $\Lp^g(f,\hg,\hh)\in I_{\hg}^{-1}\cl{O}_{\hg \hh}$ be the two-variable $p$-adic $L$-function obtained from the three-variable $p$-adic $L$-function $\Lp^g(\hf,\hg,\hh)$ introduced in \S\ref{subsec:tripleproduct} by specializing $\hf$ to $f$. Also, we define $L_p(\hg,\hh,(\mathbf{l}+\mathbf{m}-1)/2)\in I_\hg^{-1}\cl{O}_{\hg \hh}$ as the restriction of the three-variable $p$-adic $L$-function $L(\hg,\hh,\mathbf{s})$ introduced in \S\ref{subsec:Hida-Rankin} to the plane $2\mathbf{s}=\mathbf{l}+\mathbf{m}-1$. Similarly, we define the $p$-adic $L$-function $L_p(\hg,\hh\otimes\varepsilon_K,(\mathbf{l}+\mathbf{m}-1)/2)\in I_{\hg}^{-1}\cl{O}_{\hg \hh}$ by replacing $\hh$ by its twist $\hh\otimes \varepsilon_K$.


\begin{propo}\label{prop:factorization}
    We have the following equality of $p$-adic $L$-functions:
    \[
    \Lp^g(f,\hg,\hh)^2=L_p\left(\hg,\hh,\frac{\mathbf{l}+\mathbf{m}-1}{2}\right)L_p\left(\hg,\hh\otimes\varepsilon_K,\frac{\mathbf{l}+\mathbf{m}-1}{2}\right).
    \]
\end{propo}
\begin{proof}
    Since the points $(y,z)\in \tilde{\cl{W}}_{\hg}^{\mathrm{cl}}\times \tilde{\cl{W}}_{\hh}^{\mathrm{cl}}$ with weights $k_y>k_z$ and with $p\nmid \mathrm{cond}(\hg_y)\cdot \mathrm{cond}(\hh_z)$ form a Zariski dense subset of $\cl{U}_\hg\times\cl{U}_\hh$, it suffices to prove the equality after specialization at each such point. Note that for such a point we have that $(x_0,y,z)\in \cl{W}_{\hf \hg \hh}^g$.
    
    Fix such a point $(y,z)\in \tilde{\cl{W}}_\hg^{\mathrm{cl}}\times \tilde{\cl{W}}^{\mathrm{cl}}_\hh$ of weights $(l,m)$ with $l>m$ and let $c=(l+m-1)/2$. Since $V_f\cong L(1)\oplus L(\varepsilon_K)$, we deduce by Artin formalism the equality of complex $L$-values
    \[
    L(f,\hg_y^\circ,\hh_z^\circ,c)=L(\hg_y^\circ,\hh_z^\circ,c)L(\hg_y^\circ,\hh_z^\circ\otimes\varepsilon_K,c).
    \]
    Then, the equality
    \[
    \Lp^g(f,\hg,\hh)(y,z)^2=L_p\left(\hg,\hh,\frac{\mathbf{l}+\mathbf{m}-1}{2}\right)(y,z)\cdot L_p\left(\hg,\hh\otimes\varepsilon_K,\frac{\mathbf{l}+\mathbf{m}-1}{2}\right)(y,z)
    \]
    follows from Theorem~\ref{thm:Hida-3var} and Theorem~\ref{thm:hsieh}.
\end{proof}

From now on, we make the following assumption.

\begin{ass}\label{ass:nonvanishing}
    The $p$-adic $L$-function $\Lp^g(f,\hg,\hh)$ is not identically zero.
\end{ass}

\begin{remark}
    Note that, in light of Proposition~\ref{prop:factorization}, the assumption immediately implies that the $p$-adic $L$-functions $L_p(\hg,\hh,(\mathbf{l}+\mathbf{m}-1)/2)$ and $L_p(\hg,\hh\otimes\varepsilon_K,(\mathbf{l}+\mathbf{m}-1)/2)$ are not identically zero.
\end{remark}

\begin{comment}

We will make the following assumption throughout this section.

\begin{ass}
    The $p$-adic $L$-function $\Lp^g(f,\hg,\hh)$ is not identically zero.
\end{ass}

\begin{remark}
    By the previous proposition, the assumption above immediately implies that both $L_p(\hg,\hh,(\mathbf{l}+\mathbf{m}-1)/2)$ and $L_p(\hg,\hh\otimes\varepsilon_K,\mathbf{l}+\mathbf{m}-1)/2)$ are not identically zero.
\end{remark}

\end{comment}

%\subsection{General notations}

%We assume that $\ord_p(\alpha_g) < l_0-1$, $\ord_p(\alpha_h) < m_0-1$, and also that $\alpha_g \neq \beta_g, \alpha_h \neq \beta_h$. These conditions guarantee that both $g_{\alpha}$ and $h_{\alpha}$ are \'etale points of the Coleman--Mazur eigencurve. This means that for sufficiently small connected affinoid discs $U_{\hg}$ and $U_{\hh}$ in $\mathcal W_L$, centered at $l_0$ and $m_0$ respectively, there exist unique families $\hg$ and $\hh$ of tame levels $N_g$, $N_h$, character $\chi_g$, $\chi_h$ and slopes $\lambda_{\hg} = \ord_p(\alpha_g)$ and $\lambda_{\hh} = \ord_p(\alpha_h)$, specializing to $\hg_{l_0} = g_{\alpha}$ and $\hh_{m_0} = h_{\alpha}$.

%In this note we work in the setting of triple products, considering $(\hg,\hh)$ as above and $\hf$ the family passing through a precise weight-one Eisenstein series, as described in previous parts of the work. In particular, we retain the notations of Section \ref{subsec:wt1}.





\subsection{Weighted Beilinson--Flach class}\label{subsec:weightedBF}

As in previous sections, we denote by $\bb{V}_{\bm{\xi}}$ the big Galois representation associated with a Hida family $\bm{\xi}$. Note that the representations $\bb{V}_{\hh\otimes\varepsilon_K}$ and $\bb{V}_{\hh}\otimes \varepsilon_K$ are isomorphic. Since $p$ splits in $K$, an isomorphism of $\Lambda_{\hh}[G_{\bb{Q}}]$-modules $\bb{V}_{\hh\otimes\varepsilon_K}\cong \bb{V}_{\hh}\otimes \varepsilon_K$ yields an isomorphism of $\Lambda_{\hh}[G_{\bb{Q}_p}]$-modules $\bb{V}_{\hh\otimes\varepsilon_K}\cong \bb{V}_{\hh}$ and hence an isomorphism of Dieudonn\'e modules $\mathbf{D}(\bb{V}_{\hh\otimes\varepsilon_K}^+(1-\mathbf{m}-\chi_{\hh}))\cong \mathbf{D}(\bb{V}_{\hh}^+(1-\mathbf{m}-\chi_{\hh}))$. We choose an isomorphism $\iota:\bb{V}_{\hh\otimes\varepsilon_K}\rightarrow \bb{V}_\hh\otimes\varepsilon_K$ so that the composition
\[
\mathbf{D}(\bb{V}_{\hh\otimes\varepsilon_K}^+(1-\mathbf{m}-\chi_{\hh}))\xrightarrow{\iota} \mathbf{D}(\bb{V}_{\hh}^+(1-\mathbf{m}-\chi_{\hh})) \xrightarrow[]{\langle\, \cdot ,\,\bm{\omega}_\hh\rangle} \cl{O}_{\hh}
\]
agrees with the map
\[
\mathbf{D}(\bb{V}_{\hh\otimes\varepsilon_K}^+(1-\mathbf{m}-\chi_{\hh}))\xrightarrow[]{\langle\, \cdot ,\,\bm{\omega}_{\hh\otimes\varepsilon_K}\rangle} \cl{O}_{\hh}
\]
and from now on we identify $\bb{V}_{\hh\otimes\varepsilon_K}$ with $\bb{V}_\hh\otimes\varepsilon_K$ via this isomorphism.

We put $\bb{V}_{f\hg\hh}^\dagger=V_f\otimes\bb{V}_{\hg}\hat{\otimes} \bb{V}_\hh (2-\mathbf{t}_1)$ and $\bb{V}_{\hg\hh}^\dagger=\bb{V}_{\hg}\hat{\otimes}\bb{V}_\hh(2-\mathbf{t}_1)$. Note that
\[
\bb{V}_{f\hg\hh}^\dagger\cong(1\oplus \varepsilon_K)\otimes \bb{V}_{\hg\hh}^\dagger.
\]
After specializing to the plane $2\mathbf{s}=\mathbf{l}+\mathbf{m}-1$, the class $_c\kappa(\hg,\hh,\mathbf{s})$ introduced in \S\ref{subsec:BF} yields a class $_c\kappa(\hg,\hh)$ in $H^1(\bb{Q},\bb{V}_{\hg\hh}^\dagger)$. Replacing $\hh$ by $\hh\otimes\varepsilon_K$, we also obtain a class $_c\kappa(\hg,\hh\otimes\varepsilon_K)$ in $H^1(\bb{Q},\bb{V}_{\hg\hh}^\dagger\otimes\varepsilon_K)$. We define
\[
\kappa_{\hg,\hh}=\frac{(-1)^{\mathbf{s}}}{c^2-\chi_g^{-1}(c)\chi_h^{-1}(c)c} \cdot _c\kappa(\hg,\hh)\quad \text{and}\quad \kappa_{\hg,\hh\otimes\varepsilon_K}=\frac{(-1)^{\mathbf{s}}}{c^2-\chi_g^{-1}(c)\chi_h^{-1}(c)\varepsilon_K(c)c} \cdot _c\kappa(\hg,\hh\otimes\varepsilon_K).
\]
Note that, on account of Remark~\ref{rk:dependenceofBFonc}, the classes $\kappa_{\hg,\hh}$ and $\kappa_{\hg,\hh\otimes\varepsilon_K}$ do not depend on $c$. Note also that
\[
\kappa_{\hg,\hh}\in H^1_{\bal}(\bb{Q},\bb{V}_{\hg\hh}^\dagger) \quad \text{and}\quad \kappa_{\hg,\hh\otimes \varepsilon_K}\in H^1_{\bal}(\bb{Q},\bb{V}_{\hg\hh}^\dagger\otimes\varepsilon_K),
\]
where the Selmer groups $H^1_{\bal}(\bb{Q},\bb{V}_{\hg\hh}^\dagger)$ and $H^1_{\bal}(\bb{Q},\bb{V}_{\hg\hh}^\dagger\otimes\varepsilon_K)$ are defined as in \S\ref{subsec:BF}, specializing now all objects to the plane $2\mathbf{s}=\mathbf{l}+\mathbf{m}-1$.
Further, recall the basis $\{v_{f,1},v_{f,\varepsilon_K}\}$ of $V_f$ introduced in \S\ref{subsec:wt1}. We now use these elements to define the Beilinson--Flach class that we will use in the comparison. 

%In this section, as a shortcut, we write $c=(l+m-1)/2$. Recall that $V_g$ and $V_h$ are the (dual of Deligne's) $p$-adic Galois representations associated to $g$ and $h$, respectively. As in the introduction, let \[ V_{g,h} := V_g \otimes V_h(1-c). \]

%\begin{defi}
%Let \[ \kappa_{g,h} \in H^1(\mathbb Q, V_{g,h}) \quad \text{ and } \quad \kappa_{g,h \otimes \varepsilon_K} \in H^1(\mathbb Q, V_{g,h} \otimes \varepsilon_K) \] be the global Beilinson--Flach elements associated with the pair $(g,h)$ and its twist by $\varepsilon_K$, respectively.
%\end{defi} 

%Since $p$ splits in $K/\mathbb Q$, the restriction to $G_{\mathbb Q_p}$ of $V_{g,h} \otimes \varepsilon_K$ is equal to that of $V_{g,h}$. We implicitly use the identifications between $V_g \otimes V(h \otimes \varepsilon_K)$ and $V_g \otimes V_h \otimes \varepsilon_K$, as discussed in \cite[\S4.3]{BSV}. Further, recall the Galois representation $V_f$ and the basis $\{v_{f,1},v_{f,\varepsilon_K}\}$ introduced in preceding sections.



\begin{comment}

The following result allows to relate the values of the triple product $p$-adic $L$-function with the Hida--Rankin. This may be as a direct consequence of the definition of both $p$-adic $L$-functions in terms of a suitable Petersson inner product. Before doing that, we introduce the following epsilon factor, following \cite[\S9]{Das}.

\begin{defi}
Let $\epsilon(\hg,\hh):=\epsilon(\hg,\hh,\frac{{\bf l} + {\bf m}+1}{2})$ stand for the epsilon factor $\epsilon = A(\hg,\hh) \cdot B(\hg,\hh)^s$, where $s = \frac{{\bf l} + {\bf m}-1}{2}$, $B$ is the conductor of $g \otimes h$ and $A$ is a suitable algebraic number introduced in \cite[\S9]{Das}. Let $\tilde \epsilon(\hg,\hh) : =\tilde \epsilon(\hg,\hh,\frac{{\bf l} + {\bf m}-1}{2})$ denote the $p$-adic function which agrees with $A(\hg,\hh) \cdot B(\hg,\hh)^s$ on the set of integer numbers. 
%\textcolor{red}{Perd\'on por insistir, pero no creo que esto sea una funci\'on L p-\'adica (no en el sentido integral al menos, quiero decir). Y no creo que necesitemos interpolar el epsilon factor en cualquier caso.} {\color{blue} Si si, de acuerdo, lo dejamos como $p$-adic function o como prefieras!} \textcolor{red}{Es que en realidad creo que no hace falta esta definici\'on. $\Log(BF)$ deber\'ia coincidir directamente con $\Log(\kappa)^2$, sin que aparezca el epsilon factor.} {\color{blue} Es decir, yo entiendo que al comparar puede quedar asi, pero de entrada para definir el weighted BF y mirar como difieren las dos funciones L p-adicas si necesitariamos este factor epsilon, no?} \textcolor{red}{Yo creo que no. Si aplicas el logaritmo a la clase BF tal y como est\'a definida ahora yo creo que te sale $L_p(g,h)L_p(g,h,\varepsilon_K)$, que es directamente $L_p(f,g,h)$, salvo factores y dem\'as. Cuando est\'e escrita esa parte mejor igual vemos m\'as claro qu\'e necesitamos.}
\end{defi}

Note that a priori the functions $A(\hg,\hh)$ and $B(\hg,\hh)$ do depend on the chosen specialization for each of the families. 

%\textcolor{red}{
%Yo sigo sin ver esta definici\'on. Dasgupta solo var\'ia la variable ciclot\'omica ($g$ y $h$ son fijas), por eso puede decir que el epsilon factor es de la forma $\epsilon=A\cdot B^s$, con $A$ y $B$ fijadas. En nuestro caso esto no es as\'i. Cuando dije lo de cambiar la definici\'on de $BF(f,g,h)$ era con la idea de evitar lidiar con este asunto. 
%Con la nueva definici\'on lo que tenemos es que $\log_p(BF(f,g,h))$ coincide con $\log_p(\kappa(f,g,h))^2$ salvo factores algebraicos, con un cuadrado en el segundo t\'ermino, con lo cual no podemos comparar las clases. En cierto modo es como si en este caso el Teorema A de BSV fuese casi inmediato con $BF(g,h)$ jugando el papel de $\zeta_A^{\mathrm{Kato}}$ y $\kappa(f,g,h)$ jugando el papel del punto de Heegner.}

%{\color{blue} Vale, ahora veo lo que dices. Quizas partiendo de $\log_p(BF(f,g,h)) = \log_p(\kappa(f,g,h))^2$ podamos pasar a algo del tipo $BF(f,g,h) = \log_p(\kappa(f,g,h)) \cdot \kappa(f,g,h)$, del estilo de lo que hacian en DR3 (al menos siempre y cuando podamos asegurar que se cumple BK y tal vez con algunas hipotesis extra). Es decir, si tenemos dos clases que viven en el mismo espacio de dimension 1 una sera multiplo de otra y para determinar el coeficiente es suficiente compararlas aplicando el logaritmo (salvo que sean 0; que eso es de hecho lo que haciamos antes ya, pero ahora me estan entrando dudas sobre si no metiamos la pata al tachar el logaritmo, es algo que faltaba por justificar).}


\begin{propo}\label{prop:factL}
With the previous notations, the following relation between the Hida--Rankin $p$-adic $L$-function and the triple product $p$-adic $L$-function holds \[ \Lp^g(f,\hg,\hh)^2 = \epsilon(\hg,\hh) \cdot L_p \left(\hg,\hh, \frac{{\bf l} + {\bf m}+1}{2} \right)^2. \]
\end{propo}

\begin{proof}

We begin by noting that, by the classical Artin formalism, \[ L((1 \oplus \varepsilon_K) \otimes g \otimes h,s) = L(g \otimes h,s) L(g \otimes h \otimes \varepsilon_K, s), \] where $(g,h)$ are specializations of the families $(\hg,\hh)$ and $s$ is any integer.

Then, comparing the classical $L$-values along the region where the weight of $\hf$ is dominant, we have that \[ \Lp^g(f,\hg,\hh)^2 = L_p \left(\hg,\hh, \frac{{\bf l} + {\bf m}-1}{2} \right) \cdot L_p \left(\hg,\hh \otimes \varepsilon_K, \frac{{\bf l} + {\bf m}-1}{2} \right), \] where we have used a standard density argument and that the different Euler and gamma factors involved in both of them.

Finally, applying the functional equation for the Hida--Rankin $p$-adic $L$-function, as discussed e.g. in \cite[Thm. 9.4]{Das}, we get the desired result.
\end{proof}

\end{comment}

%\subsection{Weighted Beilinson--Flach classes}

%The following definition introduces the Beilinson--Flach class that we will use in this article, which is a suitable combination of the class $\kappa_{\hg,\hh}$ introduced in \S\ref{subsec:Hida-Rankin} and the class $\kappa_{\hg,\hh\otimes \varepsilon_K}$ obtained by replacing $\hh$ by its twist $\hh\otimes\varepsilon_K$. 

\begin{defi}
The \emph{weighted Beilinson--Flach class} associated with the pair $(\hg,\hh)$ is the element
\[
\mathrm{BF}(f,\hg,\hh)=v_{f,1}\otimes L_p\left(\hg,\hh\otimes\varepsilon_K,\frac{\mathbf{l}+\mathbf{m}-1}{2}\right)\kappa_{\hg,\hh}+ v_{f,\varepsilon_K} \otimes L_p\left(\hg,\hh,\frac{\mathbf{l}+\mathbf{m}-1}{2}\right)\kappa_{\hg,\hh\otimes\varepsilon_K}
\]
in $H^1(K,\bb{V}_{f\hg\hh}^\dagger)$.   
\end{defi}
 
%Observe that here the situation is different to the setting of Kato elements, since we should not weight each contribution to $\BF$ by a $p$-adic $L$-function. The reason will be clear from the reciprocity law for Beilinson--Flach elements.
Since complex conjugation acts trivially on $\BF(f, \hg, \hh)$, according to the definitions of the vectors $v_{f,1}$ and $v_{f,\varepsilon_K}$ given in \ref{subsec:wt1}, %\textcolor{blue}{(habr\'ia que justificar esto?)} \textcolor{blue}{O: puse una frase, pero igual se puede matizar mas}
this element descends to a class in $H^1(\mathbb Q, \bb{V}_{f\hg\hh}^\dagger)$ that we still denote with the same notation, i.e., we have
\[ 
\mathrm{BF}(f,\hg,\hh) \in H^1(\mathbb Q, \bb{V}_{f\hg \hh}^\dagger).
\]

\subsection{Selmer conditions}\label{subsec:Selmergroups}

Fix a finite set $\Sigma$ of places of $\mathbb Q$ containing $\infty$ and the primes dividing $N_fN_gN_hp$ and let $\mathbb Q^\Sigma$ be the maximal extension of $\mathbb Q$ unramified outside $\Sigma$.

Let $\mathscr{F}^2\bb{V}_{f\hg\hh}^\dagger\subset \bb{V}_{f\hg\hh}^\dagger$ be the $\Lambda_{\hg\hh}$-submodule of rank $4$ defined by
\[
\mathscr{F}^2\mathbb V_{f\hg\hh}^\dag=(\bb{V}_f\hat{\otimes}\bb{V}_{\hg}^+\hat{\otimes}\bb{V}_{\hh}^+ +\bb{V}_f^+\hat{\otimes}\bb{V}_{\hg}\hat{\otimes}\bb{V}_{\hh}^+ +\bb{V}_f^+\hat{\otimes}\bb{V}_{\hg}^+\hat{\otimes}\bb{V}_{\hh})(2-\mathbf{t}_1).
\]
Then, as in \S\ref{subsec:diagonalcycles}, we define the \textit{balanced Selmer group of $\bb{V}_{f\hg\hh}^\dagger$} by
\[
H^1_{\bal}(\bb{Q},\mathbb V_{f \hg \hh}^{\dag}) = \ker \bigg(
H^1(\mathbb Q^\Sigma/\mathbb Q, \mathbb V_{f \hg \hh}^{\dag}) \longrightarrow
\frac{H^1(\mathbb Q_p,\mathbb V_{f \hg \hh}^\dag)}{H^1_{\bal}(\bb{Q}_p,\mathbb V_{f \hg \hh}^\dag)}\bigg),
\]
where
\[
H^1_{\bal}(\bb{Q}_p,\mathbb V_{f \hg \hh}^\dag)=\ker \big(H^1(\mathbb Q_p, \mathbb V_{f \hg \hh}^{\dag}) \longrightarrow H^1(\mathbb Q_p, \mathbb V_{f \hg \hh}^{\dag}/\mathscr{F}^2\mathbb V_{f \hg \hh}^{\dag}) \big).
\]

We also need to introduce other Selmer conditions. Let $\bb{V}_{f\hg\hh}^g\subset \bb{V}_{f\hg\hh}^\dagger$ be the $\Lambda_{\hg\hh}$-submodule of rank $4$ defined by
\[
\bb{V}_{f\hg\hh}^g=V_f\otimes \bb{V}_{\hg}^+\hat{\otimes}\bb{V}_\hh(2-\mathbf{t}_1)
\]
and let $\bb{V}_{f\hg\hh}^{g\cup +}=\bb{V}_{f\hg\hh}^g+\mathscr{F}^2\bb{V}_{f\hg\hh}^\dagger$, which is a $\Lambda_{\hg\hh}$-submodule of $\bb{V}_{f\hg\hh}^\dagger$ of rank $5$.% and $\bb{V}_{f\hg\hh}^{g\cap +}=\bb{V}_{f\hg\hh}^g\cap\mathscr{F}^2\bb{V}_{f\hg\hh}^\dagger$, which are $\Lambda_{\hg\hh}$-submodules of $\bb{V}_{f\hg\hh}^\dagger$ of ranks $5$ and $3$, respectively.

\begin{defi}
For $\mathcal{L}\in\{\mathcal{G},\mathcal{G}\cup +\}$, the Selmer group $H^1_{\mathcal{L}}(\bb{Q},\mathbb V_{f\hg\hh}^{\dag})$ is defined by
\[
H^1_{\mathcal{L}}(\bb{Q},\mathbb V_{f\hg\hh}^{\dag}) = \ker \bigg(
H^1(\mathbb Q^\Sigma/\mathbb Q, \mathbb V_{f \hg\hh}^{\dag}) \longrightarrow
\frac{H^1(\mathbb Q_p,\mathbb V_{f \hg\hh}^\dag)}{H^1_{\mathcal{L}}(\mathbb Q_p,\mathbb V_{f \hg\hh}^\dag)}\bigg),
\]
where
\[
H_{\mathcal L}^1(\mathbb Q_p, \mathbb V_{f\hg\hh}^{\dag}) =
\begin{cases}
\ker \big(H^1(\mathbb Q_p, \mathbb V_{f\hg\hh}^{\dag}) \longrightarrow H^1(\mathbb Q_p, \mathbb V_{f \hg\hh}^{\dag}/\bb{V}_{f\hg\hh}^g) \big) &\textrm{if $\mathcal{L}=\cl{G}$},\\[0.5em]
\ker \big(H^1(\mathbb Q_p, \mathbb V_{f \hg\hh}^{\dag}) \longrightarrow H^1(\mathbb Q_p, \mathbb V_{f\hg\hh}^{\dag}/ \mathbb V_{f\hg\hh}^{g\cup +}) \big) &\textrm{if $\mathcal{L}=\mathcal{G}\cup +$}. %\\[0.5em]
%\ker \big(H^1(\mathbb Q_p, \mathbb V_{f\hg\hh}^{\dag}) \longrightarrow H^1(\mathbb Q_p, \mathbb V_{f \hg\hh}^{\dagger}/ \mathbb V_{f\hg\hh}^{g\cap +}) \big) &\textrm{if $\mathcal{L}=\mathcal{G}\cap +$}.
\end{cases}
\]
\end{defi}

\begin{comment}

\textcolor{blue}{I think it might be better to introduce the specialization $\kappa(f,g,h)$ in the subsection where we use it. I will reorganize this later.}

Using the balanced Selmer condition, we may define a {\it balanced Iwasawa Selmer group} \[ H_{\bal}^1(\mathbb Q, V_{f,\hg,\hh}) \hookrightarrow H^1(\mathbb Q, V_{f,\hg,\hh}) \] and write \[ \rho = \rho_{f,g,h} \colon V_{\hf,\hg,\hh} \rightarrow V_{f,g,h} \] for the composition of the specialization isomorphism in weights $(1,l_0,m_0)$ and the $p$-sta\-bi\-li\-za\-tion isomorphism for $(g_{\alpha},h_{\alpha})$. This induces a specialization map \[ \rho_* \colon H_{\bal}^1(\mathbb Q, V_{\hf,\hg,\hh}) \rightarrow H_{\bal}^1(\mathbb Q, V_{f,g,h}). \]  

\begin{defi}
Write \[ \kappa(f,g,h) \in H_{\bal}^1(\mathbb Q, V_{f,g,h}) \] for the global Selmer class obtained as the image of $\kappa(\hf,\hg,\hh)$ under the previous maps. Similarly, there is a class \[ \kappa(f,\hg,\hh) \in H_{\bal}^1(\mathbb Q, V_{f,\hg,\hh}) \] obtained by considering the specialization just at $\hf$, and allow $\hg$ and $\hh$ to vary in families.
\end{defi}

\end{comment}

%In particular, the class $\rho_*(\kappa(\hf,\hg,\hh))$ agrees with $\kappa(f,g,h)$ up to multiplication by a degree 4 Euler factor which may be explicitly described. \textcolor{blue}{Tal y como est\'a definida, $\kappa(f,g,h)=\rho_\ast(\kappa(\hf,\hg,\hh))$. No hay ning\'un factor de Euler.} Ok!

%\begin{remark}
%The class $\kappa(\hf,\hg,\hh)$ depends on the choice of $p$-stabilizations for $\hg$ and $\hh$, that we fix once and for all. We drop the dependence on the notation, contrary to what happened in \cite{BDV}, but we emphasize that such a dependence still exists. \textcolor{blue}{No s\'e si tiene sentido hablar de $p$-estabilizaciones de $\hg$ y $\hh$. Tiene sentido hablar de $p$-estabilizaciones de $g$ y $h$, y las familias $\hg$ y $\hh$ interpolan una de estas especializaciones. Puesto que estas especializaciones ya fueron fijadas al principio de la secci\'on, quiz\'a este remark es superfluo?} {\color{blue} Vale, saco esto, que lo tenia cuando no considerabamos familias}
%\end{remark}

\begin{comment}
As a final piece of notation, we need to introduce the following period in $L^{\times}$. Moreover, this period depends on the choice of the non-canonical isomorphism $\Gamma$.
\begin{defi}\label{def:period}
Let $\Omega_{f, \gamma} \in L^{\times}$ be the period given by \[ \Omega_{f,\gamma} = 2 \cdot \langle v_f^+, \omega_f \rangle_f. \]   
\end{defi}

\end{comment}


\subsection{Perrin-Riou maps}

In this subsection we introduce the Perrin-Riou maps that we will need for the comparison.

\begin{propo}\label{prop:2variablePR}
There exists a
homomorphism of $\Lambda_{\hg \hh}$-modules
\[
\mathcal L_{\hg \hh}^{-+}: H^1(\mathbb Q_p, \mathbb V_{\hg}^- \hat \otimes \mathbb V_{\hh}^+ (2-\mathbf{t}_1)) \rightarrow  I_{\hg}^{-1}\cl{O}_{\hg \hh},
\]
such that for every point $(y,z) \in \tilde{\cl{W}}_{\hg}^{\mathrm{cl}}\times \tilde{\cl{W}}_{\hh}^{\mathrm{cl}}$ of weights $(l,m)$ with $\alpha_{\hg_y}\beta_{\hh_z}\neq p^{(l+m-1)/2}$ the specialization of $\mathcal L_{\hg \hh}^{-+}$ at $(y,z)$ is the homomorphism
\[
\mathcal L_{{\bf gh}}^{-+}(y,z): H^1\left(\mathbb Q_p, V_{\hg_y}^- \otimes V_{\hh_z}^+\left(\frac{3-l-m}{2}\right)\right) \rightarrow \mathbb C_p
\]
given by
\[
\mathcal L_{{\bf gh}}^{-+}(y,z) = \frac{1-p^{(l+m-3)/2} \alpha_{\hg_y}^{-1} \beta_{\hh_z}^{-1}}{1-p^{(1-l-m)/2} \alpha_{\hg_y} \beta_{\hh_z}} \times \begin{cases} \frac{(-1)^{(m-l-1)/2}}{\left(\frac{m-l-1}{2}\right)!} \times \langle \log_{\BK}(\cdot), \eta_{\hg_y}^\alpha \otimes \omega_{\hh_z} \rangle & \text{ if } l < m+1 \\ \left(\frac{l-m-1}{2}\right)! \times \langle \exp_{\BK}^*(\cdot), \eta_{\hg_y}^\alpha \otimes \omega_{\hh_z} \rangle & \text{ if } l \geq m+1, \end{cases}
\]
where $\log_{\BK}$ is the Bloch--Kato logarithm and $\exp_{\BK}^*$ is the dual exponential map.
\end{propo}
\begin{proof}
    The existence and properties of the map $\cl{L}_{\hg\hh}^{-+}$ can be deduced from those of the map $\cl{L}_{\hg\hh\mathbf{s}}^{-+}$ in Proposition~\ref{perrin1} following the argument in \cite[Prop.~7.3]{BSV}.
\end{proof}

\begin{remark}
    Similarly, we have a $\Lambda_{\hg\hh}$-module homomorphism
    \[
    \mathcal L_{\hg \hh}^{+-}: H^1(\mathbb Q_p, \mathbb V_{\hg}^+ \hat \otimes \mathbb V_{\hh}^- (2-\mathbf{t}_1)) \rightarrow  I_{\hh}^{-1}\cl{O}_{\hg \hh},
    \]
    with analogous interpolation properties.
\end{remark}

\begin{lemma}
    The $\Lambda_{\hg\hh}$-module $H^1(\bb{Q}_p,\bb{V}_{\hg}^{-}\hat{\otimes}\bb{V}_{\hh}^+(2-\mathbf{t}_1))$ is torsion-free.
\end{lemma}
\begin{proof}
    To shorten notation, we write $\bb{V}=\bb{V}_{\hg}^{-}\hat{\otimes}\bb{V}_{\hh}^+(2-\mathbf{t}_1)$. Note that, as a $\Lambda_{\hg\hh}[G_{\bb{Q}_p}]$-module, we have that
    \[
    \bb{V}\cong \Lambda_{\hg\hh}(\underline{\chi}_g \underline{\alpha}_{\hg}^{-1} \underline{\alpha}_{\hh}\Theta),
    \]
    where $\underline{\chi}_g$ is the unramified character of $G_{\bb{Q}_p}$ defined by $\underline{\chi}_g(\Fr_p)=\chi_g(p)$, $\underline{\alpha}_{\hg}$ is the unramified character of $G_{\bb{Q}_p}$ defined by $\underline{\alpha}_{\hg}(\Fr_p)=a_p(\hg)$, $\underline{\alpha}_{\hh}$ is the unramified character of $G_{\bb{Q}_p}$ defined by $\underline{\alpha}_{\hh}(\Fr_p)=a_p(\hh)$ and $\Theta$ is the character of $G_{\bb{Q}_p}$ defined by
    \[
    \sigma\mapsto \omega_{\cyc}^{(l_0-m_0-1)/2}(\sigma) \langle \varepsilon_{\cyc}(\sigma)\rangle^{-1/2} [\langle \varepsilon_{\cyc}(\sigma)\rangle^{1/2}]_{\Lambda_{\hg}} \otimes [\langle \varepsilon_{\cyc}(\sigma)\rangle^{-1/2}]_{\Lambda_{\hh}}.
    \]
    Let $\lambda\in \Lambda_{\hg\hh}$. Then, we have an exact sequence
    \[
    H^0(\bb{Q}_p,\bb{V}/\lambda\bb{V})\rightarrow H^1(\bb{Q}_p,\bb{V})\rightarrow H^1(\bb{Q}_p,\bb{V}),
    \]
    where the second arrow is multiplication by $\lambda$. Thus, in order to show that the $\Lambda_{\hg\hh}$-module $H^1(\bb{Q}_p,\bb{V})$ is torsion-free, it suffices to show that $H^0(\bb{Q}_p,\bb{V}/\lambda\bb{V})=0$ for all non-zero $\lambda\in \Lambda_{\hg\hh}$. For that, it suffices to show that $H^0(\bb{Q}_p,\bb{V}/\frk{Q}\bb{V})=0$ for all height-$1$ prime ideal $\frk{Q}$ of $\Lambda_{\hg\hh}$. Note that, if $(l_0-m_0-1)/2\not\equiv 0 \pmod{p-1}$, then the inertia group $I_{\bb{Q}_p}$ acts non-trivially on $\bb{V}/\cl{I}\bb{V}$ for any proper ideal $\cl{I}$ of $\Lambda_{\hg\hh}$, so the result is immediate. Thus, we assume from now on that $(l_0-m_0-1)/2\equiv 0 \pmod{p-1}$ and therefore that $l_0\equiv m_0+1 \pmod{2(p-1)}$.

    Note that $\Lambda_{\hg\hh}$ is a finite flat extension of the unique factorization domain $\Lambda\hat{\otimes}\Lambda\simeq \bb{Z}_p[[X,Y]]$. Let $\gamma_0$ be a topological generator of $1+p\bb{Z}_p$ (e.g., we can take $\gamma_0=1+p$). If $\frk{q}$ is a height-one prime ideal of $\Lambda\hat{\otimes}\Lambda$ different from the ideal generated by $[\gamma_0]\otimes 1-1\otimes \gamma_0[\gamma_0]$, then the character $\Theta\vert_{I_{\bb{Q}_p}}:I_{\bb{Q}_p}\rightarrow ((\Lambda\hat{\otimes}\Lambda)/\frk{q})^\times$ is non-trivial and therefore $H^0(\bb{Q}_p,\bb{V}/\frk{Q}\bb{V})=0$ for any height-one prime ideal $\frk{Q}$ of $\Lambda_{\hg\hh}$ lying above $\frk{q}$.

    Now let $\frk{q}_1$ be the height-one prime ideal of $\Lambda\hat{\otimes}\Lambda$ generated by $[\gamma_0]\otimes 1-1\otimes \gamma_0[\gamma_0]$ and let $\frk{Q}_1$ be a height-one prime ideal of $\Lambda_{\hg\hh}$ above $\frk{q}_1$. Let $\frk{q}_2$ be the height-two prime ideal of $\Lambda\hat{\otimes}\Lambda$ corresponding to the arithmetic point $(P_{l-2,\bm{1}},P_{l-3,\bm{1}})\in\cl{W}\times\cl{W}$ for some integer $l>3$ such that $l\equiv l_0\pmod{2(p-1)}$. Note that $\frk{q}_1\subset \frk{q}_2$. Since $\Lambda_{\hg\hh}$ is a finite extension of $\Lambda\hat{\otimes}\Lambda$, we can find a prime ideal $\frk{Q}_2$ of $\Lambda_{\hg\hh}$ lying above $\frk{q}_2$ and such that $\frk{Q}_1\subset \frk{Q}_2$. The prime $\frk{Q}_2$ corresponds to a crystalline point $(y,z)\in \tilde{\cl{W}}_{\hg}^{\mathrm{cl}}\times \tilde{\cl{W}}_{\hh}^{\mathrm{cl}}$ of weights $(l,l-1)$. Since $l>3$, both $\hg_y$ and $\hh_z$ are $p$-old. Therefore, by the Ramanujan--Petersson conjecture, we have that the complex absolute value $\vert \alpha_{\hh_z}/\alpha_{\hg_y}\vert$ is $p^{-1/2}$. In particular $\chi_g(p)\alpha_{\hh_z}/\alpha_{\hg_y}\neq 1$. It follows that $H^0(\bb{Q}_p,\bb{V}/\frk{Q}_2\bb{V})=0$ and therefore $H^0(\bb{Q}_p,\bb{V}/\frk{Q}_1\bb{V})=0$.
\end{proof}


\begin{propo}\label{prop:injectivityofPR}
    The map $\cl{L}_{\hg\hh}^{-+}$ is injective.
\end{propo}
\begin{proof}
    By Theorem~\ref{reclaw},
    \[
    \cl{L}_{\hg\hh}^{-+}(\mathrm{pr}^{-+}(\res_p(\kappa_{\hg,\hh})))=\lambda_{N_g}(\hg)^{-1}\cdot L_p\left(\hg,\hh,\frac{\mathbf{l}+\mathbf{m}-1}{2}\right).
    \]
    In particular, it follows by Assumption~\ref{ass:nonvanishing} that the $\Lambda_{\hg\hh}$-module homomorphism $\cl{L}_{\hg\hh}^{-+}$ is non-zero. Since $H^1(\bb{Q}_p,\bb{V}_{\hg}^{-}\hat{\otimes}\bb{V}_{\hh}^+(2-\mathbf{t}_1))$ is a torsion-free $\Lambda_{\hg\hh}$-module of rank $1$, it immediately follows that $\cl{L}_{\hg\hh}^{-+}$ is injective.
\end{proof}

\begin{propo}\label{perrin2}
There exists a homomorphism of $\Lambda_{\hg \hh}$-modules
\[
\mathcal L_{f \hg \hh}^{+-+}: H^1(\mathbb Q_p, V_f^+ \otimes\mathbb V_{\hg}^- \hat \otimes \mathbb V_{\hh}^+ (2-\mathbf{t}_1)) \rightarrow I_\hg^{-1} \cl{O}_{\hg \hh },
\]
such that for every point $(y,z) \in \tilde{\cl{W}}_{\hg}^{\mathrm{cl}}\times \tilde{\cl{W}}_{\hh}^{\mathrm{cl}}$ of weights $(l,m)$ with $\alpha_{\hg_y}\beta_{\hh_z}\neq p^{(l+m-1)/2}$ the specialization of $\mathcal L_{f \hg \hh}^{+-+}$ at $(y,z)$ is the homomorphism
\[
\mathcal L_{{f\hg\hh}}^{+-+}(y,z): H^1\left(\mathbb Q_p, V_{f}^+\otimes V_{\hg_y}^- \otimes V_{\hh_z}^+\left(\frac{3-l-m}{2}\right)\right) \rightarrow \mathbb C_p
\]
given by
\[
\mathcal L_{{f\hg\hh}}^{+-+}(y,z) = \frac{1-p^{(l+m-3)/2} \alpha_{\hg_y}^{-1}\beta_{\hh_z}^{-1}}{1-p^{(1-l-m)/2} \alpha_{\hg_y} \beta_{\hh_z}} \times \begin{cases} \frac{(-1)^{(m-l-1)/2}}{\left(\frac{m-l-1}{2}\right)!} \times \langle \log_{\BK}(\cdot), \omega_{f}\otimes\eta_{\hg_y}^\alpha \otimes \omega_{\hh_z} \rangle & \text{ if } l < m+1 \\ \left(\frac{l-m-1}{2}\right)! \times \langle \exp_{\BK}^*(\cdot), \omega_{f} \otimes\eta_{\hg_y}^\alpha \otimes \omega_{\hh_z} \rangle & \text{ if } l \geq m+1, \end{cases}
\]
where $\log_{\BK}$ denotes the Bloch-Kato logarithm and $\exp_{\BK}^*$ denotes the dual exponential map.
\end{propo}
\begin{proof}
This follows as in \cite[Proposition~7.3]{BSV}, working with $V_f$ instead of $\bb{V}_{\hf}$.
\end{proof}

\begin{remark}
    Similarly, we can define maps $\cl{L}_{f\hg\hh}^{++-}$ and $\cl{L}_{f\hg\hh}^{-++}$.
\end{remark}

\begin{remark}
    The relation between $\cl{L}_{\hg\hh}^{-+}$ and $\cl{L}_{f\hg\hh}^{+-+}$ is as follows: given an element
    \[
    v\otimes z\in V_f^+\otimes H^1(\mathbb Q_p, \mathbb V_{\hg}^- \hat \otimes \mathbb V_{\hh}^+ (2-\mathbf{t}_1))\cong H^1(\mathbb Q_p, V_f^+ \otimes\mathbb V_{\hg}^- \hat \otimes \mathbb V_{\hh}^+ (2-\mathbf{t}_1)),
    \]
    we have that
    \[
    \cl{L}_{f\hg\hh}^{+-+}(v\otimes z)=\langle v,\omega_f\rangle \cdot \cl{L}_{\hg\hh}^{-+}(z).
    \]
\end{remark}



\subsection{The explicit comparison}\label{subsec:explicit_comparison}

We can now present one of the main results of this note, which was already anticipated in the introduction. The result gives a direct comparison between the weighted Beilinson--Flach class introduced in \S\ref{subsec:weightedBF} and a diagonal cycle class, in analogy with the main results of \cite{LR1}. 

\begin{comment}
As we later emphasize, the way that has been used to define the Beilinson--Flach class (as a sum of two different contributions) makes unnecessary to take a projection of the diagonal cycle class. \textcolor{blue}{No s\'e si queda muy claro este comentario en cuanto a que no es necesario tomar una proyecci\'on de la clase diagonal.} \textcolor{red}{Yo diria que se entiende bien, y me parece instructivo en el sentido de decir que al haber definido ese BF como suma ponderada ya tenemos las dos clases en el mismo espacio.} \textcolor{blue}{Lo que no veo claro es la frase `makes unnecessary to take a projection of the diagonal cycle class'. Qu\'e proyecci\'on habr\'ia que tomar si no definimos BF as\'i?} \textcolor{red}{Me estas convenciendo de que igual es ambigua, yo lo pensaba que por ejemplo en lo que habiamos trabajado con David (que eran comparaciones de sistemas de Euler, pero desde un punto de vista distinto) teniamos el ciclo diagonal en algo de dimension 8, desde donde proyectabamos de forma natural a algo de dimension 4, y esa clase proyeccion era la que comparabamos con el BF. Aqui es mas bien al reves, que pegamos los dos BF para tener algo que ya vive en el mismo sitio que el ciclo diagonal.} \textcolor{blue}{Bueno, igual se puede modificar un poco entonces el comentario. O dejarlo as\'i de momento, como veas.}
\end{comment}

%{\color{blue} O: cambio entonces lo de $\log_p$ por $\mathcal L^{+-+} \circ \res_p$ que es consistente con lo anterior}


\begin{comment}
Comento esto puesto que la suposici\'on se sigue ahora de las hip\'otesis del teorema.

\begin{defi}\label{def-log-bk-families}
We write $\mathcal L^{+-+} \circ \res_p$ for the Perrin-Riou regulator defined in $H_{\bal}^1(\mathbb Q, V_{f,\hg,\hh})$ corresponding to localization at $p$ followed by the composition with the natural projection \[ H^1(\mathbb Q_p, V_{f,\hg,\hh}) \longrightarrow H^1(\mathbb Q_p, V_f^+ \otimes V_{\hg}^- \otimes V_{\hh}^+), \] the Bloch--Kato logarithm and the pairing with the canonical differentials.
\end{defi}

We introduce the following assumption, which may be thought as a rather mild condition.

\begin{ass}\label{ass-inj-fam}
The Perrin-Riou logarithm $\mathcal L^{+-+} \circ \res_p$ of Definition~\ref{def-log-bk-families} is injective. {\color{blue}{En el caso en familias entiendo que una cosa asi deberia ser mas facil, pero tampoco veo claro del todo como hacerlo, asi que en principio lo dejo como assumption.}}
\end{ass}

\end{comment}

Recall that the class $\kappa_{\hg,\hh}$ belongs to the balanced Selmer group $H^1_{\bal}(\bb{Q},\bb{V}_{\hg\hh}^\dagger)$. Therefore, the image of $\res_p(\kappa_{\hg,\hh})$ in $H^1(\bb{Q}_p,\bb{V}_{\hg}^-\hat{\otimes}\bb{V}_{\hh}(2-\mathbf{t}_1))$ lands in $H^1(\bb{Q}_p,\bb{V}_{\hg}^-\hat{\otimes}\bb{V}_{\hh}^+(2-\mathbf{t}_1))$ and the image of $\res_p(\kappa_{\hg,\hh})$ in $H^1(\bb{Q}_p,\bb{V}_{\hg}\hat{\otimes}\bb{V}_{\hh}^-(2-\mathbf{t}_1))$ lands in $H^1(\bb{Q}_p,\bb{V}_{\hg}^+\hat{\otimes}\bb{V}_{\hh}^-(2-\mathbf{t}_1))$. We denote these images by $\res_p(\kappa_{\hg,\hh})^{-+}$ and $\res_p(\kappa_{\hg,\hh})^{+-}$, respectively. Similarly we define $\res_p(\kappa_{\hg,\hh\otimes\varepsilon_K})^{-+}$ and $\res_p(\kappa_{\hg,\hh\otimes\varepsilon_K})^{+-}$.



We now establish the key proposition towards the comparison that we will prove later.

\begin{proposition}\label{prop:localconditionforBF}
    The image of $\res_p(\BF(f, \hg, \hh))$ in $H^1(\bb{Q}_p, V_f \otimes \bb{V}_{\hg}^- \hat{\otimes} \bb{V}_{\hh}(2-\mathbf{t}_1))$ lands in $H^1(\bb{Q}_p,V_f^+\otimes \bb{V}_\hg^- \hat{\otimes} \bb{V}_\hh^+(2-\mathbf{t}_1))$ and is given by
    \[
    v_f^+\otimes \left(L_p\left(\hg,\hh\otimes\varepsilon_K,\frac{\mathbf{l}+\mathbf{m}-1}{2}\right)\res_p(\kappa_{\hg,\hh})^{-+}+L_p\left(\hg,\hh,\frac{\mathbf{l}+\mathbf{m}-1}{2}\right)\res_p(\kappa_{\hg,\hh\otimes\varepsilon_K})^{-+}\right)
    \]
\end{proposition}
\begin{proof}
    Since $\kappa_{\hg,\hh}$ belongs to $H^1_{\bal}(\bb{Q},\bb{V}_{\hg\hh}^\dagger)$ and $\kappa_{\hg,\hh\otimes\varepsilon_K}$ belongs to $H^1_{\bal}(\bb{Q},\bb{V}_{\hg\hh}^\dagger\otimes\varepsilon_K)$, the image of $\res_p(\BF(f, \hg, \hh))$ in $H^1(\bb{Q}_p, V_f \otimes \bb{V}_{\hg}^- \hat{\otimes} \bb{V}_{\hh}(2-\mathbf{t_1}))$ lies in $H^1(\bb{Q}_p, V_f \otimes \bb{V}_{\hg}^- \hat{\otimes} \bb{V}_{\hh}^+(2-\mathbf{t_1}))$. Therefore, to conclude the proof it suffices to show that the class
    \[
    v_f^-\otimes \left(L_p\left(\hg,\hh\otimes\varepsilon_K,\frac{\mathbf{l}+\mathbf{m}-1}{2}\right)\res_p(\kappa_{\hg,\hh})^{-+}-L_p\left(\hg,\hh,\frac{\mathbf{l}+\mathbf{m}-1}{2}\right)\res_p(\kappa_{\hg,\hh\otimes\varepsilon_K})^{-+}\right)
    \]
    in $H^1(\bb{Q}_p, V_f^- \otimes \bb{V}_{\hg}^- \hat{\otimes} \bb{V}_{\hh}^+(2-\mathbf{t_1}))$ is zero. Using Theorem~\ref{reclaw}, we have that
    \begin{align*}
        \mathcal{L}_{\hg \hh}^{-+} &\left(L_p\left(\hg,\hh\otimes\varepsilon_K,\frac{\mathbf{l}+\mathbf{m}-1}{2}\right)\res_p(\kappa_{\hg,\hh})^{-+}-L_p\left(\hg,\hh,\frac{\mathbf{l}+\mathbf{m}-1}{2}\right)\res_p(\kappa_{\hg,\hh\otimes\varepsilon_K})^{-+}\right) \\
        &= \lambda_{N_g}(\hg)^{-1}\cdot L_p\left(\hg,\hh\otimes\varepsilon_K,\frac{\mathbf{l}+\mathbf{m}-1}{2}\right)L_p\left(\hg,\hh,\frac{\mathbf{l}+\mathbf{m}-1}{2}\right) \\ &-\lambda_{N_g}(\hg)^{-1}\cdot L_p\left(\hg,\hh,\frac{\mathbf{l}+\mathbf{m}-1}{2}\right)L_p\left(\hg,\hh\otimes\varepsilon_K,\frac{\mathbf{l}+\mathbf{m}-1}{2}\right)=0.
    \end{align*}
    Since the map $\cl{L}_{\hg \hh}^{-+}$ is injective by Proposition~\ref{prop:injectivityofPR}, the result follows.
\end{proof}

\begin{remark}\label{rk:localconditionforBF}
    As a consequence of the previous proposition, we have that
    \[
    \BF(f,\hg,\hh)\in H^1_{\cl{G}\cup+}(\bb{Q},\bb{V}_{f\hg\hh}^\dagger).
    \]
\end{remark}


Let $\kappa(f,\hg,\hh)\in H^1_{\bal}(\bb{Q},\bb{V}_{f\hg\hh}^\dagger)$ be the class obtained from the class $\kappa(\hf,\hg,\hh)\in H^1_{\bal}(\bb{Q},\bb{V}_{\hf\hg\hh}^\dagger)$ introduced in Theorem~\ref{thm:reclawdiagonalcycles} by specializing $\hf$ to $f$. We denote by $\res_p(\kappa(f,\hg,\hh))^{+-+}$ the image of the class $\res_p(\kappa(f,\hg,\hh))$ in $H^1(\bb{Q}_p,V_f^+\otimes\bb{V}_\hg^-\hat{\otimes}\bb{V}_\hh^+(2-\mathbf{t}_1))$. Similarly, we denote by $\res_p(\BF(f,\hg,\hh))^{+-+}$ the image of $\res_p(\BF(f,\hg,\hh))$ in $H^1(\bb{Q}_p,V_f^+\otimes\bb{V}_\hg^-\hat{\otimes}\bb{V}_\hh^+(2-\mathbf{t}_1))$.

We also introduce the following definition.

\begin{defi}\label{def:period}
Let $\Omega_{f, \gamma} \in L^{\times}$ be the $p$-adic period given by
\[
\Omega_{f,\gamma} = 2 \cdot \langle v_f^+, \omega_f \rangle_f.
\] 
\end{defi}

\begin{remark}
The definition of $\Omega_{f,\gamma}$ depends on the isomorphism $\gamma$ fixed in \S\ref{subsec:wt1}, hence the notation.    
\end{remark}

%To shorten notation, we put $\BF(f,\hg,\hh)_p=\res_p(\BF(f,\hg,\hh)$ and we denote by $\BF(f,\hg,\hh)_p^{+-+}$ the image of $\BF(f,\hg,\hh)_p$ in $H^1(\bb{Q}_p,V_f^+\otimes \bb{V}_\hg^- \hat{\otimes} \bb{V}_\hh^+(2-\mathbf{t}_1))$.

\begin{comment}
\begin{proposition}
The image of $\res_p(\BF(f, \hg, \hh))$ under the map \[ H^1(\mathbb Q_p, V_{f,\hg,\hh}) \longrightarrow H^1(\mathbb Q_p, V_f^- \otimes V_{\hg} \otimes V_{\hh}(1-{\bf c})) \] induced by the projection $V_f \otimes V_{\hg} \otimes V_{\hh} \rightarrow V_f^- \otimes V_{\hg} \otimes V_{\hh}$ is equal to \[ v_f^-(L_p(\hg,\hh\otimes\varepsilon_K,(\mathbf{l}+\mathbf{m}-1)/2) \cdot \res_p(\kappa_{\hg,\hh}) - L_p(\hg,\hh,(\mathbf{l}+\mathbf{m}-1)/2) \cdot \res_p(\kappa_{\hg,\hh \otimes\varepsilon_K})), \] and lies in the subspace corresponding to $V_f^- \otimes V_{\hg}^+ \otimes V_{\hh}^+$.  \textcolor{blue}{No creo que queramos esto, no? Lo que queremos al final es una clase en $V_f^+\otimes\bb{V}_\hg^-\otimes\bb{V}_\hh^+$.}
\end{proposition}

\begin{proof}
The first assertion directly follows from the definition. In order to verify the second claim, we show that the other three components are zero.

\begin{itemize}
\item[(a)] The projection to the quotient $V_{\hg}^- \otimes V_{\hh}^-$ vanishes because of the local properties of Beilinson--Flach classes recalled in Proposition \ref{prop:local-property}.

\item[(b)] Next, we consider the projection to $V_{\hg}^- \otimes V_{\hh}^+$. This also vanishes, but now this is because its image under the isomorphism given by the Perrin-Riou map is zero, using our definition of the class as certain weighted sum of the two Beilinson--Flach classes. Then, we may conclude since the Perrin-Riou map is an isomorphism thanks to Assumption \ref{non-exc-zeros}.

\item[(c)] The projection to $V_{\hg}^+ \otimes V_{\hh}^-$ is zero for the same reason, since the epsilon factors involved in the functional equation for $L_p(\hg,\hh)$ and $L_p(\hh,\hg)$ are the same. %{\color{blue}{O: Podemos poner algo mas de explicacion, diciendo que la ecuacion funcional es la misma porque tanto la A como la B vienen de la ecuacion funcional compleja donde no importa el orden}}. 
In particular, \[ \begin{aligned} L_p(\hg,\hh\otimes\varepsilon_K,(\mathbf{l}+\mathbf{m}-1)/2) \cdot L_p(\hh,\hg,(\mathbf{l}+\mathbf{m}-1)/2) & - \\ L_p(\hg,\hh,(\mathbf{l}+\mathbf{m}-1)/2) \cdot L_p(\hh \otimes \varepsilon_K,\hg,(\mathbf{l}+\mathbf{m}-1)/2) & = 0. \end{aligned} \]

{\color{blue} O: Este argumento me esta liando mas de lo que pensaba e igual se puede perfilar o escribir mejor, ahora no esta muy detallado, pero deberia funcionar bien.} We now justify the previous equality. Observe that the complex functional equation of \cite[Thm. 9.3]{Das} holds regardless of the constraints among the weights, since the complex $L$-function is symmetric in both modular forms. Then, in Thm. 9.4, we may also consider the case where the weight inequality is reversed, taking instead the comparison with the complex object $D(g,f,\cdot, \cdot)$; this leads us to a pair of equations like (71) and (72), now for $D(g,f,\cdot,\cdot)$ and $D_p(f,g,\cdot,\cdot)$. Then, we apply equation (73), which allows us to come back from $D(g,f,\cdot,\cdot)$ to a (complex) symmetric expression in both variables, and the conclusion is the same, getting an identical relation between the $p$-adic $L$-function and the complex one, via the epsilon factor. This means that we get again the same epsilon factors when the weights are not in the order prescribed in \cite[Thm. 9.4]{Das} and the result is still valid.
% Alternativamente, intuyo que podemos jugar con la forma del factor epsilon, trabajando con las funciones p-adicas A y B y viendo si son analiticas y demas.

Note that, again, the Perrin-Riou map is an isomorphism thanks to Assumption \ref{non-exc-zeros}.
\end{itemize} 
\end{proof}

\end{comment}

\begin{theorem}\label{thm:main}
    The equality
    \[
    \Omega_{f,\gamma} \cdot \cl{L}_{f\hg\hh}^{+-+}(\res_p(\kappa(f,\hg,\hh))^{+-+})^2 =\lambda_{N_g}(\hg) \cdot \mathcal{L}_{f\hg\hh}^{+-+}(\res_p(\BF(f, \hg, \hh))^{+-+})
    \]
    holds.
\end{theorem}
\begin{proof}
    %We first note that, for any element
    %\[
    %v\otimes z \in V_f^+ \otimes H^1(\bb{Q}_p,\bb{V}_\hg^-\hat{\otimes} \bb{V}_\hh^+(2-\mathbf{t}_1))\cong H^1(\bb{Q}_p,V_f^+\otimes \bb{V}_{\hg}^-\hat{\otimes}\bb{V}_{\hh}^+(2-\mathbf{t}_1)),
    %\]
    %we have that
    %\[
    %\mathcal{L}_{f\hg\hh}^{+-+}(v\otimes z)=\langle v,\omega_f\rangle \cdot\mathcal{L}_{\hg\hh}^{-+}(z).
    %\]
    By Proposition~\ref{prop:localconditionforBF}, we have that
    \begin{align*}
        \mathcal{L}^{+-+}_{f\hg\hh}(\res_p(\BF(f,\hg,\hh))^{+-+})&=\langle v_f^+,\omega_f\rangle \cdot L_p\left(\hg,\hh\otimes\varepsilon_K,\frac{\mathbf{l}+\mathbf{m}-1}{2}\right)\cl{L}_{\hg\hh}^{-+}\left(\res_p(\kappa_{\hg,\hh})^{-+}\right) \\
        &+\langle v_f^+,\omega_f\rangle \cdot L_p\left(\hg,\hh,\frac{\mathbf{l}+\mathbf{m}-1}{2}\right)\cl{L}_{\hg\hh}^{-+}\left(\res_p(\kappa_{\hg,\hh\otimes\varepsilon_K})^{-+}\right).
    \end{align*}
    Using Theorem~\ref{reclaw}, it follows that
    \[
    \mathcal{L}^{+-+}_{f\hg\hh}(\res_p(\BF(f,\hg,\hh))^{+-+})=\frac{\Omega_{f,\gamma}}{\lambda_{N_g}(\hg)}\cdot L_p\left(\hg,\hh,\frac{\mathbf{l}+\mathbf{m}-1}{2}\right)L_p\left(\hg,\hh\otimes\varepsilon_K,\frac{\mathbf{l}+\mathbf{m}-1}{2}\right).
    \]
    Hence, by Proposition~\ref{prop:factorization},
    \[
    \lambda_{N_g}(\hg)\cdot \mathcal{L}^{+-+}_{f\hg\hh}(\res_p(\BF(f,\hg,\hh))^{+-+})=\Omega_{f,\gamma}\cdot \Lp^g(f,\hg,\hh)^2,
    \]
    which is equal to
    \[
    \Omega_{f,\gamma} \cdot \cl{L}^{+-+}_{f\hg\hh}(\res_p(\kappa(f,\hg,\hh))^{+-+})^2
    \]
    by Theorem~\ref{thm:reclawdiagonalcycles}.
\end{proof}

%\begin{remark}
%Concerning the Rankin--Selberg periods, the corresponding result of \cite{BDV} involves the logarithm of a unit, which plays the role of a period. In our setting, however, the reciprocity laws for Beilinson--Flach classes and for diagonal cycles are formulated directly in terms of Rankin--Selberg periods, so the unit appearing in \cite[Thm. 5.2]{BDV} is no longer present. Thus, in our case there is no such ambiguity in the choice of periods, and we naturally obtain the Petersson norm on both sides. \textcolor{blue}{Revisar si este remark deber\'ia ir despu\'es del corolario anterior o en otro sitio. Tambi\'en estar\'ia bien citar el resultado preciso de BDV al que nos estamos refiriendo.} \textcolor{red}{Tampoco se si aqui pega mucho, si ves que es innecesario tambien lo puedes sacar.} \textcolor{blue}{Como veas}
%\end{remark}

\begin{corollary}\label{corollary:comparison}
    Assume that $H^1_{\cl{G}\cup +}(\bb{Q}, \bb{V}_{f\hg\hh}^\dagger)$ is a torsion-free $\Lambda_{\hg\hh}$-module of rank $1$. Then $\BF(f,\hg,\hh)$ belongs to $H^1_{\bal}(\bb{Q}, \bb{V}_{f\hg\hh}^\dagger)$ and
    \[
    \lambda_{N_g}(\hg)\cdot\BF(f,\hg,\hh)=\Omega_{f,\gamma}\cdot \Lp^g(f,\hg,\hh) \cdot \kappa(f,\hg,\hh).
    \]

    \begin{comment}
    \textcolor{blue}{Ideally one would like to assume that the balanced Selmer group (rather than the $\cl{G}\cup +$ Selmer group) is torsion-free of rank $1$. If we only vary one of the modular forms in a CM family, that would suffice using the arguments in \cite[\S7.4]{ACR}. But in the present setting it is not so clear to me.} \textcolor{red}{Mi idea para esto seria ver que la clase esta en el balanced Selmer group, para lo que digamos seria la parte $(-,+,-)$ tambien es cero, haciendo un razonamiento como el de la proposicion 5.4 pero para la $h$, lo ves razonable o no? (Es decir, algo asi pasaria por relacionar los epsilon factor para las funciones $L$ $p$-adicas de $(g,h)$ y $(h,g)$) En cualquier caso, si te parece que es mucho rollo podemos dejarlo asi (ya te digo, yo tenia ese sketch en mente pero igual escribirlo bien no es tan trivial). Lo que venia a decir es que cuando hacemos el mismo juego para la $h$ nos quedan cuatro funciones $L$ $p$-ádicas: para $(g,h)$ y $(g,h\otimes \epsilon_K)$ por un lado y para $(h,g)$ y $(h \otimes \epsilon_K,g)$ por otro. Y yo lo que decia es que podemos pasar de $(g,h)$ a $(g,h \otimes \epsilon_K)$ via un cierto epsilon factor y de $(h,g)$ a $(h \otimes \epsilon_K)$ via el mismo factor epsilon, adaptando un poco los argumentos sobre la ecuacion funcional que hay en Dasgupta. Y ademas es que esto al fin y al cabo no deja de ser equivalente a lo que estamos suponiendo, que este cacharro de dimension 5 sea libre y de rango 1 viene a decir que esta clase tiene que estar en el subespacio de dimension 4, porque el ciclo diagonal ya sabemos que lo esta. Te dejo aqui comentado lo que esta comentado en otro lado y que era mi idea de prueba.}
    \textcolor{blue}{Entiendo que ver que la parte $(-,+,-)$ es cero es equivalente a ver que $L_p(\hg,\hh)L_p(\hh\otimes\varepsilon_K,\hg)=L_p(\hg,\hh\otimes\varepsilon_K))L_p(\hh,\hg)$, que ahora mismo est\'a puesto como parte del corolario siguiente. Y esto viene a ser lo mismo que decir que los factores epsilon $p$-\'adicos coinciden, es decir que $\epsilon_p(\hg,\hh)=\epsilon_p(\hh,\hg)$. Pero esta parte no la veo obvia. B\'asicamente quieres decir que tienes una funci\'on meromorfa que vive en $\Frac(\cl{O}_\hg\hat{\otimes}\cl{O}_{\hh})$ que interpola el epsilon factor complejo $\epsilon(g,h)$ en las dos regiones. Pero esto no me parece que sea inmediato.}
    \textcolor{red}{Pero dos cosas diferentes, que igual me estoy liando: (1) olvidandonos de los factores epsilon, con esa igualdad del corolario siguiente sobre funciones $L$ $p$-adicas llegariamos a lo que queremos, porque nos quedaria la resta de las dos cosas, no? (2) sobre cada factor epsilon, yo entiendo que lo podemos ver como un cociente de dos funciones $L$ $p$-adicas, es decir $L_1 = \epsilon L_2$, por lo que $\epsilon$ seria una funcion meromorfa, no? En cualquier caso, por lo que dices me parece que con (1) llegaria, no?} \textcolor{blue}{Pero la demostraci\'on del corolario siguiente utiliza este corolario. Y respecto a los factores epsilon, tenemos una funci\'on meromorfa que interpola el factor epsilon en una regi\'on, y una funci\'on meromorfa que lo interpola en la otra, simplemente escribi\'endolos como cocientes de funciones L p-\'adicas. Lo que igual no es obvio es que estas dos funciones meromorfas coinciden, que es lo que nos dice el corolario siguiente.}
    \textcolor{red}{Ah, vale, ok, es decir, que habria que demostrarlo independientemente. Pero por eso digo, mi idea es que son funciones meromorfas, y mi plan era ver que coincidian digamos sobre un conjunto denso, como los enteros. Usando Dasgupta en principio lo que tienes es que en una sabes el factor para $k \geq l$ y en la otra para $l \geq k$, con lo que la interseccion seria solo una linea. Pero es lo que escribia debajo, que yo creo que ese argumento funciona tambien cuando $l \geq k$ (el del Thm. 9.4 de Dasgupta), porque podemos hacer el juego con $D(g,f,s)$ y llegar a lo mismo.} \textcolor{blue}{Pero yo creo que las regiones son $k>l$ y $k<l$, entonces no hay intersecci\'on. Yo no veo claro que el argumento usando Dasgupta funcione.} \textcolor{red}{Es decir, pero la idea es que tu puedes tener el $D_p(f,g,...)$ cuando el peso de $f$ es menor (porque eso esta definido en familias y ahi nada afecta), y ahi considerar para comparar el $D(g,f,...)$. Entonces en Dasgupta deriva las ecuaciones (71) y (72) diciendo, bueno, pues basicamente estamos cambiando $f$ y $g$ por unos multiplos y estas cosas son lineales en $f$ y $g$, asi que si lo cambiamos por unos ciertos multiplos entonces el factor de escala es el mismo, de ahi lo de $R(sigma)$ y $\tilde R(sigma)$. Y tu luego tienes una ecuacion como la 73 que ya te devuelve a $\Lambda(f \otimes g,)$ porque eso es simetrico, y yo entiendo que funciona igual (es decir, no hay ninguna comparacion basada en que estes en la region de interpolacion, sino que el unico problema es que para definir $D(f,g,...)$ necesitas que el peso de $f$ sea mayor, pero tu puedes coger $D(g,f,...)$ y no pasaria nada} \textcolor{blue}{Pero es que Dasgupta no est\'a variando $f$ y $g$, est\'a variando la variable ciclot\'omica. $R$ es una funci\'on meromorfa en $s$. No sabemos c\'omo var\'ia con los pesos.} \textcolor{red}{Ya ya, eso por supuesto, pero lo que quiero decir es que tu podrias hacer ese juego para $L_p(g,h)$ y para $L_p(h,g)$, yendote a la funcion compleja. Y lo que verias es que la $R$ asociada a ambas es la misma. Entonces tu tienes dos funciones de tres variables cada una $R_1$ y $R_2$. Y lo que estamos diciendo es que para cada $s$ fijo $R_1(k,l,s)=R_2(k,l,s)$, por lo que para cada $s$, $R_1$ tendria que coincidir con $R_2$, no?} \textcolor{blue}{Pero $f$ y $g$ no son intercambiables, no? O sea, la definici\'on de $D(f,g)$ requiere $k\geq l$, yo creo. Entonces, digamos que igual si $k=l$, entonces puedes hacer el juego de intercambiar $f$ y $g$ y ver que $R_1(k,k,s)=R_2(k,k,s)$} \textcolor{red}{Pero mira, dicho de otra manera, lo que quiero decir es que si supones que $k \geq l$, puedes considerar entonces $D(f,g,...)$, pero tienes dos objetos $p$-adicos, $D_p(f,g,...)$ y $D_p(g,f,...)$, y en ambos se traduce de la misma manera que cambies $f$ y $g$ por un multiplo, por lo que esas ecuaciones funcionales van a ser iguales y van a replicar lo que pasa para el objeto correspondiente a $D(f,g,...)$, donde aqui claro, lo ordenas poniendo primero el del peso mayor y luego el del menor para que eso exista, pero en esencia no importa mucho porque este objeto luego se te relaciona con uno que es simetrico. Entonces lo que digo es que al fin y al cabo $D_p(f,g,...)$ y $D_p(g,f,...)$ replican ambos por linealidad las transformaciones del unico objeto complejo, y por eso el factor de Euler seria el mismo (no se si me explico bien).} \textcolor{blue}{Pero una cosa, c\'omo sabes que la ecuaci\'on funcional para $D_p(g,f)$ va a replicar lo que ocurre con $D(f,g)$. Yo entiendo que para $D_p(f,g)$ lo sabes porque tienes la f\'ormula de interpolaci\'on de la proposici\'on 5.4.1 de LLZ. Digamos, para demostrar (71) en Dasgupta, yo tengo que relacionar $D(W_N(f),W_N(g),y,s)$ con $D(f^\ast,g^\ast,1/N,s)$, y luego s\'e que $D_p(W_N(f),W_N(g),y,s)$ y $D_p(f^\ast,g^\ast,1/N,s)$ se relacionan de la misma manera porque para $k\geq l$ vienen dados por la f\'ormula en la proposici\'on 5.4.1 de LLZ. Pero no veo por qu\'e la relaci\'on entre $D_p(W_N(g),W_N(g),y,s)$ y $D_p(g^\ast,f^\ast,1/N,s)$ habr\'ia de ser la misma que la relaci\'on entre $D(W_N(f),W_N(g),y,s)$ y $D(f^\ast,g^\ast,1/N,s)$.} {\color{red} Ah, vale, este me parece el punto importante: tu dices que las ecuaciones (71) y (72) vienen de que la $p$-adica interpola la compleja y que sale de ahi entonces el multiplo. Es que yo lo estaba pensando de forma diferente (olvidandonos de la interpolacion): tu tienes que ambas cosas, la $D$ y la $D_p$, son lineales en el sentido de que si cambias $f$ (o $g$) por un multiplo entonces la expresion correspondiente a $D$ o a $D_p$ cambia por el mismo multiplo. Entonces yo lo que entiendo es que dices tienes un factor de escala en la ecuacion funcional compleja que es el mismo que en la $p$-adica, porque el multiplo o factor de escala es el mismo, y por eso decia, que va a ser el mismo en $D_p(f,g)$ que en $D_p(g,f)$. Es que de hecho Dasgupta no menciona la propiedad de interpolacion ni se complica con que tengas puntos en el rango de interpolacion ni nada, para mi que eso viene de la linealidad, entonces te permite hacer el mismo razonamiento que Dasgupta hace para $D_p(f,g)$ para $D_p(g,f)$.} \textcolor{blue}{Pero es que yo no creo que venga solo de la linealidad. Tenemos que $W_N(f)=f^\ast(nz)$ y $W_N(g)=g^\ast(mz)$, donde $n=N/N_f$ y $m=N/N_g$. Para llegar a (72) a partir de (70) tienes que relacionar $D(f^\ast(nz),g^\ast(mz),y,s)$ para todo $y\in \frac{1}{N}\bb{Z}/\bb{Z}$ con $D(f^\ast,g^\ast,1/N,s)$, y entiendo que aqu\'i hace uso del hecho de que esencialmente esto son productos de Petersson. No veo que baste apelar a la linealidad para relacionar estas dos expresiones.} \textcolor{red}{Le dare una vuelta, pero se me hace raro todo, es decir, tu dices que el argumento que vale para $D_p(f,g,...)$ no vale para $D_p(g,f,...)$ porque se esta usando la interpolacion, no? En cualquier caso, lo podemos dejar asi por ahora y pensar si vemos alguna manera viable de mejorar el resultado, pero se me hace raro no ser capaces de probar algo asi (que moralmente es verdad porque es equivalente a la suposicion que estamos haciendo). Y sobre todo que Dasgupta no apela en absoluto a la interpolacion, asi que da como la sensacion que tuviese otro argumento en mente, no? (yo entendia lo de la linealidad, pero es cierto que no esta claro). Por ejemplo, tu crees entonces que lo de que diga que sirve tb para $k=l$ esta mal?} \textcolor{blue}{No creo que est\'e mal lo de que sirva para $k=l$, porque la proposici\'on 5.4.1 de LLZ tambi\'en incluye el caso $k=l$. Pero en todo caso hace falta algo m\'as que la linealidad para derivar (72). Y yo creo que para ver que el factor $R$ que sale en (71) coincide con el que sale en (72) para $s$ entero se requiere usar alguna propiedad sobre $D_p(f,g)$, aunque no est\'e dicho expl\'icitamente. Para m\'i un motivo para ser un poco esc\'eptico es que una identidad de la forma $L_p(\hg,\hh)L_p(\hh^\ast,\hg^\ast)=L(\hg^\ast,\hh^\ast)L(\hh,\hg)$, que es lo que se tendr\'ia si el argumento funcionase, parece lo suficientemente interesante como para que est\'e enunciada en alg\'un lado.}

    {\color{red} Y otra cosa: si te fijas la propiedad de interpolacion sirve solo para los $s$ tq $l \leq s \leq k-1$, es decir, que si fijas $k$ y $l$ solo serviria para un numero finito (porque aqui no puedes variar caracter) y sin embargo dicen que lo tienen siempre... lo que digo es que esto me hace pensar que la demo no va via densidad.} \textcolor{blue}{Para $D_p(f,g)$ puedes usar la Proposici\'on 5.4.1 de LLZ y para $D(f,g)$ puedes usar la f\'ormula antes del Teorema 4.2.3. Ninguna de las dos impone ninguna restricci\'on sobre $s$. Y yo creo que de alguna manera usando estas f\'ormulas ves que $R(s)$ y $\tilde{R}(s)$ coinciden cuando $s$ es entero. No necesitas estar dentro del rango de interpolaci\'on, pero s\'i necesitas $k\geq l$.}
    \end{comment}

\begin{comment}
\item[(c)] The projection to $V_{\hg}^+ \otimes V_{\hh}^-$ is zero for the same reason, since the epsilon factors involved in the functional equation for $L_p(\hg,\hh)$ and $L_p(\hh,\hg)$ are the same. %{\color{blue}{O: Podemos poner algo mas de explicacion, diciendo que la ecuacion funcional es la misma porque tanto la A como la B vienen de la ecuacion funcional compleja donde no importa el orden}}. 
In particular, \[ \begin{aligned} L_p(\hg,\hh\otimes\varepsilon_K,(\mathbf{l}+\mathbf{m}-1)/2) \cdot L_p(\hh,\hg,(\mathbf{l}+\mathbf{m}-1)/2) & - \\ L_p(\hg,\hh,(\mathbf{l}+\mathbf{m}-1)/2) \cdot L_p(\hh \otimes \varepsilon_K,\hg,(\mathbf{l}+\mathbf{m}-1)/2) & = 0. \end{aligned} \]

{\color{blue} O: Este argumento me esta liando mas de lo que pensaba e igual se puede perfilar o escribir mejor, ahora no esta muy detallado, pero deberia funcionar bien.} We now justify the previous equality. Observe that the complex functional equation of \cite[Thm. 9.3]{Das} holds regardless of the constraints among the weights, since the complex $L$-function is symmetric in both modular forms. Then, in Thm. 9.4, we may also consider the case where the weight inequality is reversed, taking instead the comparison with the complex object $D(g,f,\cdot, \cdot)$; this leads us to a pair of equations like (71) and (72), now for $D(g,f,\cdot,\cdot)$ and $D_p(f,g,\cdot,\cdot)$. Then, we apply equation (73), which allows us to come back from $D(g,f,\cdot,\cdot)$ to a (complex) symmetric expression in both variables, and the conclusion is the same, getting an identical relation between the $p$-adic $L$-function and the complex one, via the epsilon factor. This means that we get again the same epsilon factors when the weights are not in the order prescribed in \cite[Thm. 9.4]{Das} and the result is still valid.
% Alternativamente, intuyo que podemos jugar con la forma del factor epsilon, trabajando con las funciones p-adicas A y B y viendo si son analiticas y demas.
\end{comment}
\end{corollary}
\begin{proof}
    Let $\mathrm{pr}^{+-+}:H^1_{\cl{G}\cup +}(\bb{Q}_p,\bb{V}_{f\hg\hh}^\dagger)\rightarrow H^1(\bb{Q}_p,V_f^+\otimes\bb{V}_\hg^-\hat{\otimes}\bb{V}_\hh^+(2-\mathbf{t}_1))$ be the map induced by the natural projection $\bb{V}_{f\hg\hh}^{g\cup +}\rightarrow V_f^+\otimes\bb{V}_\hg^-\hat{\otimes}\bb{V}_\hh^+(2-\mathbf{t}_1)$. Since the element $\Lp^g(f,\hg,\hh)\in I_{\hg}^{-1}\cl{O}_{\hg\hh}$ is non-zero by Assumption~\ref{ass:nonvanishing}, it follows from Theorem~\ref{thm:reclawdiagonalcycles} that the map
    \[
    \cl{L}_{f\hg\hh}^{+-+}\circ\mathrm{pr}^{+-+}\circ\res_p: H^1_{\cl{G}\cup +}(\bb{Q},\bb{V}_{f\hg\hh}^\dagger) \longrightarrow I_{\hg}^{-1} \cl{O}_{\hg\hh}
    \]
    is non-zero and therefore injective by our assumptions on $H^1_{\cl{G}\cup +}(\bb{Q}, \bb{V}_{f\hg\hh}^\dagger)$. Since both $\kappa(f,\hg,\hh)$ and $\BF(f,\hg,\hh)$ belong to $H^1_{\cl{G}\cup +}(\bb{Q},\bb{V}_{f\hg\hh}^\dagger)$, the result now follows immediately from the previous theorem. 
\end{proof}


\begin{remark}\label{rk:torsionfreeness}
    Regarding the assumptions in the previous corollary, it is expected by sign considerations that the $\Lambda_{\hg\hh}$-module $H^1_{\bal}(\bb{Q}, \bb{V}_{f\hg\hh}^\dagger)$ has rank $1$ and that the $\Lambda_{\hg\hh}$-module $H^1_{\cl{G}}(\bb{Q}, \bb{V}_{f\hg\hh}^\dagger)$ has rank zero, which would easily imply that the $\Lambda_{\hg\hh}$-module $H^1_{\cl{G}\cup +}(\bb{Q},\bb{V}_{f\hg\hh}^\dagger)$ has rank $1$. Moreover, we can ensure that $H^1_{\cl{G}\cup +}(\bb{Q}, \bb{V}_{f\hg\hh}^\dagger)$ is torsion-free by imposing the condition that $H^0(\bb{Q},\overline{\rho}^\dagger)=0$, where $\overline{\rho}^\dagger$ is the residual representation attached to $\bb{V}_{f\hg \hh}^\dagger$.
\end{remark}


\begin{remark}
It may be instructive to discuss the analogies and differences between a comparison of this kind and that carried out in \cite{LR1}, where we also used a Coleman family passing through a critical Eisenstein series. Here, the main idea is to construct cohomology classes over $K$, that we can then descend and compare with a class over $\mathbb Q$. However, in \cite{LR1}, the classes are all defined over $\mathbb Q$, so the comparison is between a suitable projection of one of the classes and the other.
Moreover, weight-one modular forms behave in an ostensibly different way, since both $p$-stabilizations are ordinary. 
\end{remark}



\begin{corollary}
    Assume that $H^1_{\cl{G}\cup +}(\bb{Q}, \bb{V}_{f\hg\hh}^\dagger)$ is a torsion-free $\Lambda_{\hg\hh}$-module of rank $1$. Then
    \begin{align*}
     \Lp^g(f,\hg,\hh)\Lp^h(f,\hg,\hh)&=\frac{\lambda_{N_g}(\hg)}{\lambda_{N_h}(\hh)}\cdot L_p\left(\hg,\hh,\frac{\mathbf{l}+\mathbf{m}-1}{2}\right)L_p\left(\hh\otimes\varepsilon_K,\hg,\frac{\mathbf{l}+\mathbf{m}-1}{2}\right) \\ &=\frac{\lambda_{N_g}(\hg)}{\lambda_{N_h}(\hh)}\cdot L_p\left(\hg,\hh\otimes\varepsilon_K,\frac{\mathbf{l}+\mathbf{m}-1}{2}\right)L_p\left(\hh,\hg,\frac{\mathbf{l}+\mathbf{m}-1}{2}\right).
    \end{align*}
\end{corollary}
\begin{proof}
    By the previous corollary, we know that $\BF(f,\hg,\hh)\in H^1_{\bal}(\bb{Q},\bb{V}_{f\hg\hh}^\dagger)$. Therefore, the image of $\BF(f,\hg,\hh)$ in $H^1(\bb{Q}_p,V_f^-\otimes\bb{V}_{\hg}\hat{\otimes}\bb{V}_{\hh}^{-}(2-\mathbf{t}_1))$ is zero. This image is given by
    \[
    v_f^-\otimes \left(L_p\left(\hg,\hh\otimes\varepsilon_K,\frac{\mathbf{l}+\mathbf{m}-1}{2}\right)\res_p(\kappa_{\hg,\hh})^{+-}-L_p\left(\hg,\hh,\frac{\mathbf{l}+\mathbf{m}-1}{2}\right)\res_p(\kappa_{\hg,\hh\otimes\varepsilon_K})^{+-}\right),
    \]
    so it follows that
    \[
    L_p\left(\hg,\hh\otimes\varepsilon_K,\frac{\mathbf{l}+\mathbf{m}-1}{2}\right)\res_p(\kappa_{\hg,\hh})^{+-}-L_p\left(\hg,\hh,\frac{\mathbf{l}+\mathbf{m}-1}{2}\right)\res_p(\kappa_{\hg,\hh\otimes\varepsilon_K})^{+-}=0.
    \]
    Applying the Perrin-Riou map $\mathcal{L}_{\hg\hh}^{+-}$, we deduce by Theorem~\ref{reclaw} that
    \begin{align}
    &L_p\left(\hg,\hh,\frac{\mathbf{l}+\mathbf{m}-1}{2}\right)L_p\left(\hh\otimes\varepsilon_K,\hg,\frac{\mathbf{l}+\mathbf{m}-1}{2}\right) \nonumber \\
    &=L_p\left(\hg,\hh\otimes\varepsilon_K,\frac{\mathbf{l}+\mathbf{m}-1}{2}\right)L_p\left(\hh,\hg,\frac{\mathbf{l}+\mathbf{m}-1}{2}\right). \label{eq:equalityofHRpadicLfunctions}
    \end{align}
    Now note that, for any element
    \[
    v\otimes z \in V_f^+ \otimes H^1(\bb{Q}_p,\bb{V}_\hg^+\hat{\otimes} \bb{V}_\hh^-(2-\mathbf{t}_1))\cong H^1(\bb{Q}_p,V_f^+\otimes \bb{V}_{\hg}^+\hat{\otimes}\bb{V}_{\hh}^-(2-\mathbf{t}_1)),
    \]
    we have that
    \[
    \mathcal{L}_{f\hg\hh}^{++-}(v\otimes z)=\langle v,\omega_f\rangle \cdot\mathcal{L}_{\hg\hh}^{+-}(z).
    \]
    Let $\mathrm{pr}^{++-}:H^1_{\bal}(\bb{Q}_p,\bb{V}_{f\hg\hh}^\dagger)\rightarrow H^1(\bb{Q}_p,V_f^+\otimes\bb{V}_\hg^+\hat{\otimes}\bb{V}_\hh^-(2-\mathbf{t}_1))$ be the map induced by the natural projection $\mathscr{F}^2\bb{V}_{f\hg\hh}^\dagger\rightarrow V_f^+\otimes\bb{V}_\hg^+\hat{\otimes}\bb{V}_\hh^-(2-\mathbf{t}_1)$. Using Theorem~\ref{reclaw} and equation~\ref{eq:equalityofHRpadicLfunctions}, the image of $\BF(f,\hg,\hh)$ by $\cl{L}_{f\hg\hh}^{++-}\circ \mathrm{pr}^{++-}\circ \res_p$ is equal to
    \[
    \Omega_{f,\gamma}\cdot\lambda_{N_h}(\hh)^{-1}\cdot L_p\left(\hg,\hh,\frac{\mathbf{l}+\mathbf{m}-1}{2}\right)L_p\left(\hh\otimes\varepsilon_K,\hg,\frac{\mathbf{l}+\mathbf{m}-1}{2}\right).
    \]
    Also, by Theorem~\ref{thm:reclawdiagonalcycles}, the image of $\kappa(f,\hg,\hh)$ by $\cl{L}_{f\hg\hh}^{++-}\circ \mathrm{pr}^{++-}\circ \res_p$ is equal to $\Lp^h(f,\hg,\hh)$. Hence, the result now follows from the previous corollary.
    
\end{proof}

\begin{remark}
    Note that the $p$-adic $L$-functions involved in the statement have disjoint interpolation ranges, so the result does not follow by a direct comparison of complex $L$-values.
\end{remark}

\begin{comment}
\begin{theorem}\label{thm:bf-vs-cycles-fam}
The equality \[ \Omega_{f,\gamma} \cdot \log_p(\kappa(f,\hg,\hh))^2 = \log_p(\BF(f, \hg, \hh)) \] holds. In particular, if Assumption~\ref{ass-inj-fam} is true, then \[ \Omega_{f,\gamma} \cdot \log_p(\kappa(f,\hg,\hh)) \cdot \kappa(f,\hg,\hh) = \BF(f, \hg, \hh) \]  holds in the balanced Selmer group $H_{\bal}^1(\mathbb Q, V_{f,\hg,\hh})$.  \textcolor{blue}{No hemos visto que $\BF(f, \hg, \hh)$ pertenezca al balanced Selmer group.}

{\color{blue}{Tal y como hemos normalizado creo que no hay dependencia $\mathcal C(l,m)$ como estaba antes en la derecha, pero revisa eso, Raul, que podria ser (no afectaria nada).}}
\end{theorem}

\begin{proof}

We begin by noting that under the current assumptions the balanced Selmer group is one-dimensional, so it suffices to compare the images of both (two-variable) classes under the corresponding Perrin-Riou map. Let $\chi$ denote either $\varepsilon_K$ or the trivial character $1$. Then, the reciprocity law for the Beilinson--Flach class asserts that \[ \langle \log_{\hg \hh}^{-+}(\res_p(\BF^{-+}(f, \hg, \hh)), \eta_{\hg} \otimes \omega_{\hh} \rangle = L_p(\hg, \hh \otimes \chi, 1+s). \] Recall that we are working under the assumption that the $c$-factor does not vanish (and in particular we have inverted it). As usual, the superscript $-+$ in the local Beilinson--Flach element stands for the projection to the corresponding local piece. Note that this morphism makes sense because of the previous proposition. 

From the previous results, we have that the class belongs to the kernel of the composition of the Perrin-Riou logarithm and the pairing with the canonical differentials: \[ \langle \log_{\hg \hh}^{-+}, \eta_{\hg} \otimes \omega_{\hh} \rangle \colon H^1(\mathbb Q, V_{\hg,\hh}) \longrightarrow I_{\hg}^{-1} \cl{O}_{\hf \hg \hh}. \] The composition factors through the morphism induced in cohomology by the projection $D(V_{\hg}) \otimes D(V_{\hh}) \rightarrow D(V_{\hg}^-) \otimes D(V_{\hh}^+)$, and the resulting map is injective under the assumption that there are no exceptional zeros. We will use now the injectivity of that map to compare both classes.

Note that the morphism $\mathcal L_{f \hg \hh}^{+-+} \, : \, H^1(\mathbb Q_p, D(V_{\hg})^- \otimes D(V_{\hh})^+) \otimes_L V_f^+ \longrightarrow I_{\hg}^{-1} \cl{O}_{\hf \hg \hh}$ may be written as \[ \mathcal L_{fgh}^{+-+}(v \otimes z) = \langle v,\omega_f \rangle_f \cdot \log_{gh}^{-+}(z). \] Then, \[ \mathcal L_{f \hg \hh}^{+-+}(\res_p(\BF(f, \hg, \hh))) = \Omega_{f,\gamma} \cdot L_p(\hg,\hh,1+s), \] and the explicit reciprocity law for diagonal cycles gives \[ \log_{f \hg \hh}^{+-+} \circ \res_p(\kappa(f,\hg,\hh)) = \Lp^g(f,\hg,\hh). \]

As we have seen before in Prop. \ref{prop:factL}, \[ \Lp^g(f,\hg,\hh) = \epsilon(\hg,\hh) \cdot L_p(\hg,\hh). \]
Hence, we have that the difference between $\Omega_{f,\gamma} \cdot \kappa(f,\hg,\hh)$ and $\BF(f, \hg, \hh)$ is killed by the linear form $\log_{f\hg \hh}^{+-+} \circ \res_p$, so the difference defines an element  in the balanced Selmer group. Then, the explicit reciprocity law of Kings--Loeffler--Zerbes \cite{KLZ} implies that the Bloch--Kato Selmer group vanishes, thus concluding the proof.
\end{proof}
\end{comment}


\begin{comment}
\begin{remark}
Recall that in the definition of the Beilinson--Flach classes we had inverted the $c$-factor, and this is the reason of why it does not appear in the following statement.
\end{remark}
\end{comment}

\subsection{Formulae for specializations}\label{subsec:nonexceptional}


We can now specialize the Hida families $(\hg,\hh)$ and obtain similar results for the fixed modular forms $(g,h)$. We write $(g_\alpha,h_\alpha)=(\hg_{y_0},\hh_{z_0})$.

Let $\cl{L}_{g_\alpha h_\alpha}^{-+}:H^1(\bb{Q}_p,V_g^-\otimes V_h^+(1-c_0))\rightarrow \bb{C}_p$ and $\cl{L}_{fg_\alpha h_\alpha}^{+-+}:H^1(\bb{Q}_p,V_f^+\otimes V_g^-\otimes V_h^+(1-c_0))\rightarrow \bb{C}_p$ be the maps obtained from the maps $\cl{L}_{\hg\hh}^{+-}$ and $\cl{L}_{f\hg\hh}^{+-+}$ introduced above by specializing $(\hg,\hh)$ to $(g_\alpha,h_\alpha)$. Note that, for any element
\[
v\otimes z \in V_f^+ \otimes H^1(\bb{Q}_p,V_g^-\otimes V_h^+(1-c_0))\cong H^1(\bb{Q}_p,V_f^+\otimes V_g^-\otimes V_h^+(1-c_0)),
\]
we have that
\[
\mathcal{L}_{fg_\alpha h_\alpha}^{+-+}(v\otimes z)=\langle v,\omega_f\rangle \cdot\mathcal{L}_{gh}^{-+}(z).
\]
Also note that, if $\alpha_g \beta_h\neq p^{l_0+m_0-1}$, then we have
\[
\cl{L}_{g_\alpha h_\alpha}^{-+}= \frac{1-p^{(l_0+m_0-3)/2} \alpha_{g}^{-1} \beta_{h}^{-1}}{1-p^{(-l_0-m_0+1)/2} \alpha_{g} \beta_{h}} \times \begin{cases} \frac{(-1)^{(m_0-l_0-1)/2}}{\left(\frac{m_0-l_0-1}{2}\right)!} \times \langle \log_{\BK}(\cdot), \eta_{g_\alpha}^\alpha \otimes \omega_{h_\alpha} \rangle & \text{ if } l_0 \leq m_0 \\ \left(\frac{l_0-m_0-1}{2}\right)! \times \langle \exp_{\BK}^*(\cdot), \eta_{g_\alpha}^\alpha \otimes \omega_{h_\alpha} \rangle & \text{ if } l_0 > m_0. \end{cases}
\]

To shorten notation, we write $V_{fgh}^\dagger=V_f\otimes V_g\otimes V_h(1-c_0)$ and we define Selmer conditions and Selmer groups for this representation analogous to the ones introduced in \S\ref{subsec:Selmergroups} for the representation $\bb{V}_{f\hg\hh}^\dagger$. We denote by $\kappa(f,g_\alpha,h_\alpha)\in H^1_{\bal}(\bb{Q},V_{fgh}^\dagger)$ the specialization of the class $\kappa(f,\hg,\hh)$ and we denote by $\BF(f,g_\alpha,h_\alpha)\in H^1_{\cl{G}\cup +}(\bb{Q},V_{fgh}^\dagger)$ the specialization of the class $\BF(f,\hg,\hh)$.

\begin{proposition}\label{prop:main-nonexceptional}
    The equality
    \[
    \Omega_{f,\gamma} \cdot \cl{L}_{fg_\alpha h_\alpha}^{+-+}(\res_p(\kappa(f,g_\alpha,h_\alpha))^{+-+})^2 =\lambda_{N_g}(g) \cdot \mathcal{L}_{f g_\alpha h_\alpha}^{+-+}(\res_p(\BF(f, g_\alpha, h_\alpha))^{+-+})
    \]
    holds.
\end{proposition}
\begin{proof}
    This is an immediate consequence of Theorem~\ref{thm:main}.
\end{proof}


\begin{corollary}\label{cor:non-exceptionalcomparison}
    Assume that $H^1_{\bal}(\bb{Q}, V_{fgh}^\dagger)$ is $1$-dimensional and that $\Lp^g(f,g_\alpha,h_\alpha)\neq 0$. If $l_0=2$ and $m_0=1$, further assume that $\alpha_g\beta_h\neq 1$. Then $\BF(f,g_\alpha,h_\alpha)$ belongs to $H^1_{\bal}(\bb{Q}, V_{fgh}^\dagger)$ and
    \[
    \lambda_{N_g}(g)\cdot\BF(f,g_\alpha,h_\alpha)=\Omega_{f,\gamma}\cdot \Lp^g(f,g_\alpha,h_\alpha) \cdot \kappa(f,g_\alpha,h_\alpha).
    \]
\end{corollary}
\begin{proof}
    Arguing as in the proof of \cite[Thm~9.5]{ACR}, we have $H^1_{\cl{G}\cup+}(\bb{Q},V_{fgh}^\dagger)=H^1_{\bal}(\bb{Q},V_{fgh}^\dagger)$. Now the proof follows as in Corollary~\ref{corollary:comparison}.
\end{proof}

\begin{remark}
We emphasize that the class $\kappa(f,g_\alpha,h_\alpha)$ can be interpreted as an anticyclotomic cohomology class via the identifications discussed in \S\ref{subsec:ac-cycles}.
\end{remark}

\begin{comment}


\textcolor{blue}{Estoy pensando que igual no hace falta realmente hacer ninguna suposici\'on aqu\'i sobre ceros excepcionales. Simplemente en ese caso los resultados que obtenemos ser\'an triviales porque nos quedar\'a b\'asicamente $0=0$, pero entiendo que siguen siendo ciertos.}

\begin{ass}\label{non-exc-zeros}
With the notations introduced above, it holds that
\[
\alpha_g \beta_h \neq p^{\frac{l+m-1}{2}}.
\]
%\quad \text{ and } \quad \alpha_h \beta_g \neq p^{\frac{l+m-1}{2}}. \] {\color{blue} O: aqui creo que estaba mal porque ponia $l+m+1$, cuando el centro de la ecuacion funcional y lo que sale en el contexto del triple product es $l+m+1-2=l+m-1$.}
\end{ass}

This assumption implies that the denominator that appears in the interpolation property characterizing the maps $\cl{L}_{\hg\hh}^{-+}$ and $\cl{L}_{f\hg\hh}^{+-+}$ does not vanish at the point $(y_0,z_0)$ corresponding to $(g,h)$.
\end{comment}

%\begin{defi}\label{def-log-bk}
%We write $\mathcal L^{+-+} \circ \res_p$ for the Perrin-Riou regulator defined in $H_{\bal}^1(\mathbb Q, V_{f,g,h})$ corresponding to localization at $p$ followed by the composition with the natural projection \[ H^1(\mathbb Q_p, V_{f,g,h}) \longrightarrow H^1(\mathbb Q_p, V_f^+ \otimes V_g^- \otimes V_h^+), \] the Bloch--Kato logarithm and the pairing with the canonical differentials.
%\end{defi}

\begin{comment}
Esto se sigue ahora de las hip\'otesis.

\begin{ass}\label{ass-inj}
The Perrin-Riou logarithm $\mathcal L^{+-+} \circ \res_p$ is injective on $H_{\bal}^1(\mathbb Q, V_{f,g,h})$.
\end{ass}

\begin{remark}
This assumption may be reformulated in terms of a Bloch--Kato conjecture about the Bloch--Kato Selmer group $H_{\bal}^1(\mathbb Q, V_{f,g,h})$ being one-dimensional, which is expected to hold quite generally.
\end{remark}
\end{comment}

%In this section, we impose that there is not an exceptional zero at $(x_0,y_0,z_0)$ for the $p$-adic $L$-function $\Lp^g(\hf,\hg,\hh)$ arising from the Euler factor
%\[
%1-\frac{\alpha_g\beta_h}{p^{c_0-1}}
%\]
%in the interpolation property of the Perrin--Riou map $\mathcal{L}_{\hf\hg\hh}^{+-+}$ introduced in Proposition~\ref{perrin}. This amounts to the following assumption.



%{\color{blue} No se hasta que punto puede ser viable demostrar esto, aunque sea en algun caso. Es decir, $V_f$ deberia ser muy bien comportada, pero no lo veo del todo claro. Del mismo modo que en la conclusion de BDV (al final de la seccion 4.3 justifican que si los logaritmos son iguales las clases tambien, aqui tenemos que hacer lo mismo, aunque sea con alguna hipotesis extra (tipo, que una cierta funcion L tiene algun valor no nulo o asi).}

\begin{comment}

\begin{corollary}\label{thm:bf-vs-cycles}
The equality \[ \Omega_{f,\gamma} \cdot \log_p(\kappa(f,g,h))^2 = \log_p(\BF(f, g, h)) \] holds. In particular, if Assumption~\ref{ass-inj} is true, then \[ \Omega_{f,\gamma} \cdot \log_p(\kappa(f,g,h)) \cdot \kappa(f,g,h) = \BF(f, g, h) \]  holds in the balanced Selmer group $H_{\bal}^1(\mathbb Q, V_{f,g,h})$.
\end{corollary}

\end{comment}

%\begin{proof}
%This follows from Theorem \ref{thm:bf-vs-cycles-fam}.


\begin{comment}
We begin by noting that under the current assumptions the balanced Selmer group is one-dimensional, so it suffices to compare the images of both classes under the corresponding Perrin-Riou map.

Let $\chi$ denote either $\varepsilon_K$ or the trivial character $1$. Then, the reciprocity law for the Beilinson--Flach class asserts that \[ \langle \log_{gh}^{-+}(\res_p(\BF^{-+}(f, g, h)), \eta_g \otimes \omega_h \rangle = L_p(g_{\alpha}, h \otimes \chi, 1+s), \] where we have placed ourselves in the assumptions that the $c$-factor does not vanish and in particular we have inverted it. Further, the superscript $-+$ in the Beilinson--Flach elements stands for the projection to the corresponding local piece.

Moreover, note that the previous morphism makes sense because the image of $\res_p(\BF(f, g, h))$ under the map \[ H^1(\mathbb Q_p, V_{f,g,h}) \longrightarrow H^1(\mathbb Q_p, V_f^- \otimes V_g \otimes V_h(1-c)) \] induced by the projection $V_f \otimes V_g \otimes V_h \rightarrow V_f^- \otimes V_g \otimes V_h$ is equal to \[ v_f^-(\res_p(\kappa_{g,h}) - \epsilon(g,h) \cdot \res_p(\kappa_{g,h \otimes\varepsilon_K})), \] which lies in the subspace corresponding to $V_g^+ \otimes V_h^+$. In order to verify this claim, we make the following observations:

\begin{itemize}
\item[-] The projection to the quotient $V_g^- \otimes V_h^-$ vanishes because of the local properties of Beilinson--Flach classes; see e.g. \cite[Prop. 8.1.7]{KLZ}.
\item[-] Once we know that, we may consider the projection to $V_g^- \otimes V_h^+$. This also vanishes, but now because its image under the isomorphism given by the Perrin-Riou map is zero, using our definition of the class as certain weighted sum of the two Beilinson--Flach classes. Then, we may conclude since the Perrin-Riou map is an isomorphism thanks to Assumption \ref{non-exc-zeros}.
\item[-] The projection to $V_g^+ \otimes V_h^-$ is zero for the same reason. Again, the Perrin-Riou map is an isomorphism thanks to Assumption \ref{non-exc-zeros}.
\end{itemize}

From the previous results, we have that the class belongs to the kernel of the composition of the Perrin-Riou logarithm and the pairing with the canonical differentials: \[ \langle \log_{gh}^{-+}, \eta_g \otimes \omega_h \rangle \colon H^1(\mathbb Q, V_{g,h}) \longrightarrow \mathbb Q_p. \] The composition factors through the morphism induced in cohomology by the projection $D(V_g) \otimes D(V_h) \rightarrow D(V_g^-) \otimes D(V_h^+)$, and the resulting map is injective under the assumption that there are no exceptional zeros. We will use now the injectivity of that map to compare both classes.

Note that the morphism $\mathcal L_{fgh}^{+-+} \, : \, H^1(\mathbb Q_p, D(V_g)^- \otimes D(V_h)^+) \otimes_L V_f^+ \longrightarrow \mathbb Q_p$ may be written as \[ \mathcal L_{fgh}^{+-+}(v \otimes z) = \langle v,\omega_f \rangle_f \cdot \log_{gh}^{-+}(z). \] Then, \[ \mathcal L_{fgh}^{+-+}(\res_p(\BF(f, g, h))) = \Omega_{f,\gamma} \cdot L_p(g_{\alpha},h,1+s). \]

The explicit reciprocity law for diagonal cycles gives \[ \log_{fgh}^{+-+} \circ \res_p(\kappa(f,g,h)) = \Lp^g(f,g_{\alpha},h). \]

As we have seen before in Prop. \ref{prop:factL}, \[ \Lp^g(f,g_{\alpha},h) = \epsilon(g,h) \cdot L_p(g_{\alpha},h). \]
Hence, we have that the difference between $\Omega_{f,\gamma} \cdot \kappa(f,g,h)$ and $\BF(f, g, h)$ is killed by the linear form $\log_{fgh}^{+-+} \circ \res_p$, so the difference defines an element  in the balanced Selmer group. Then, the explicit reciprocity law of Kings--Loeffler--Zerbes \cite{KLZ} implies that the Bloch--Kato Selmer group vanishes, thus concluding the proof.
\end{comment}
%\end{proof}




\begin{comment}
\subsection{Proof of the main theorem: anticyclotomic diagonal cycles}

\textcolor{blue}{En esta secci\'on no estamos diciendo realmente nada nuevo, no? Es b\'asicamente reenunciar lo de la secci\'on anterior.} {\color{red} O: si si, cierto, es que inicialmente en mi cabeza la idea era comparar nuestras clases con los BF usando los ciclos diagonales normales como puente, pero vamos, que al final nuestras clases y los ciclos diagonales ya son de forma muy obvia lo mismo. Igual llega hacer un remark diciendo que la clase de ciclos diagonales se puede interpretar como el $\kappa_{\psi,g,h}$ de ACR.} \textcolor{blue}{S\'i, yo quiz\'a har\'ia simplemente un remark. Comentamos esta secci\'on, entonces?}

We establish now an explicit comparison between the anticyclotomic diagonal classes of \cite{ACR} and the diagonal cycles that we used in previous sections. Recall in particular the anticyclotomic class $\kappa_{\psi,g,h}$.

\begin{theorem}\label{thm:acr-vs-cycles}
Assume that the complex $L$-series $L(g,h,s)$ vanishes at the central critical point $s=(l+m-1)/2$. Then, the class $\kappa(f,g,h)$ belongs to the Selmer group $\Sel(\mathbb Q, V(f,g,h))$ and the equality \[ \mathcal L^{+-+} \circ \res_p(\kappa(f,g,h)) = \langle \log_{fgh}^{+-+}(\res_{\mathfrak p}(\kappa^{+-+}_{\psi,g,h})), \omega_{\psi} \otimes \eta_g \otimes \omega_h \rangle \] holds in $L$ up to multiplication by an explicit non-zero constant in the number field $K(g,h)$. \textcolor{blue}{Si estamos en la regi\'on no balanceada la anulaci\'on del valor L nos dice que la igualdad anterior es $0=0$. Y si estamos en la regi\'on no balanceada yo creo que la anulaci\'on del valor L no nos dice nada.}
\end{theorem}

\begin{proof}
This follows by the construction of the classes in \cite{ACR}, which are obtained by applying the Shapiro isomorphism to the diagonal cycle classes of \cite{BSV}. Then, these elements satisfy an explicit reciprocity law relating them with the triple product $p$-adic $L$-function (see e.g. \S7 of \emph{loc.\,cit.}).
\end{proof}

\begin{remark}
In \cite{ACR}, the result is affected by multiplication by a factor of $\log_p(u_{\mathfrak p})$, where $u_{\mathfrak p}$ in $\mathcal O_K[1/p]^{\times}$ is a generator of $\mathfrak p^{h_K}$, with $h_K$ the class number of $K$. This is a consequence of the use of different periods for the $p$-adic $L$-functions used there, which does not occur in our case since in all the cases we use the Petersson norm.
\end{remark}


The main theorem of the introduction is a consequence of Theorem \ref{thm:bf-vs-cycles} and Theorem \ref{thm:acr-vs-cycles}.

\begin{theorem}
With the previous notations, the following equality holds:
\[ \mathcal L^{-+} \circ \res_p(\kappa_{g,h}) = \mathcal L^{+-+} \circ \res_p(\kappa_{\psi,g,h}). \] %where the $\log_p$ in the left is the map $\langle \log_{gh}^{-+}(\cdot), \eta_g \otimes \omega_h \rangle$ and the $\log_p$ in the right is the map $\langle \log_{fgh}^{+-+}(\cdot), \omega_{\psi} \otimes \eta_g \otimes \omega_h \rangle$
\end{theorem}

\begin{proof}
According to Theorem \ref{thm:bf-vs-cycles} and Theorem \ref{thm:acr-vs-cycles}, both sides agree with the logarithm of the diagonal cycle, so the result follows.    
\end{proof}

\begin{remark}
Following the development of \cite{LR1} and \cite{PR25}, it is natural to wonder if one consider instead a single Hida family passing through a weight one Eisenstein series, and look for a connection with the system of circular units. Unfortunately, Iwasawa cohomology vanishes over one half of the weight space, and the na\"ive approach would just yield an equality of the form $0=0$. We expect to investigate further this issue in forthcoming work.
\end{remark}

\end{comment}

\subsection{A factorization formula for the big logarithm of a Beilinson--Flach class}

As a consequence of the cyclotomic results of \cite{BDV}, B\"uy\"ukboduk, Casazza, and Sakamoto \cite{BC}, \cite{BS} obtained an expression of the Hida--Rankin $p$-adic $L$-function in terms of the Ochiai big logarithm of the Kato class. Here, we discuss an analogue of \cite[Prop. 8.9]{BC}, showing that the image of $\BF(f,\hg,\hh)$ under a suitable Perrin-Riou map factors as the product of two triple product $p$-adic $L$-functions. 
%\textcolor{blue}{No s\'e si acabo de ver la analog\'ia con BC. En nuestro caso lo que tenemos es una factorizaci\'on del logaritmo de una cierta clase como producto de funciones $L$ $p$-\'adicas.} %In particular, we express the triple product $p$-adic $L$-function as a product of a Rankin--Selberg and a logarithm of a Beilinson--Kato class using our previous result. However, the setting here is slightly different: the self-duality condition of the triple product makes that the two $p$-adic $L$-functions appearing in the formula are the same, and hence one does not get the extra BDP-type $p$-adic $L$-function.

\begin{comment}
\textcolor{blue}{No s\'e si tiene sentido introducir aqu\'i $\bm{\omega}_{\hg}$ y $\bm{\omega}_{\hh}$ cuando ya han aparecido varias veces a lo largo del art\'iculo.}

We begin by fixing some notations, following \cite[\S3]{BC}. For that purpose, we consider the $\Lambda$-adic module \[ M = \mathbb V_{\hg}^+(\varepsilon_{\cyc} \varepsilon_{\hg}^{-1}) \otimes \mathbb V_{\hh}^+(\varepsilon_{\cyc} \varepsilon_{\hh}^{-1}) \simeq \Lambda_{\hg}(\alpha_{\hg}^{-1}) \otimes \Lambda_{\hh}(\alpha_{\hh}^{-1}). \] To canonically trivialize the module $\mathbb D_{\cris}(M)$, we need Ohta's work \cite{ohta00}, in the same way as it is applied in \cite[Prop. 10.1.1]{KLZ}. Then, there exists a canonical isomorphism \[ \omega_{\hg} \otimes \omega_{\hh} \colon \mathbb D_{\cris}(M) \longrightarrow \Lambda_{\hg} \otimes \Lambda_{\hh} \] coming from the tensor product of the two isomorphisms \[ \omega_{\hg} \colon \mathbb D_{\cris}(\mathbb V_{\hg}^+) \longrightarrow \Lambda_{\hg}, \quad \omega_{\hh} \colon \mathbb D_{\cris}(\mathbb V_{\hh}^+) \longrightarrow \Lambda_{\hh} \]
\end{comment}

% (Si lo ponemos en esta secci\'on ya se entiende) Along this section, we keep same the assumptions and notations from \S\ref{subsec:explicit_comparison}. In particular, once we know that the Beilinson--Flach class lies in $H^1_{\bal} (\bb{Q},\bb{V}_{f\hg\hh}^\dagger)$, we can project to $\mathbb V_{\hf}^- \otimes \mathbb V_{\hg}^+ \otimes \mathbb V_{\hh}^+$ and consider the pairing with $\eta_f \otimes \omega_{\hg} \otimes \omega_{\hh}$


Let $\cl{L}_{f\hg\hh}^{-++}:H^1(\bb{Q}_p,V_f^-\otimes \bb{V}_{\hg}^+\hat{\otimes}\bb{V}_{\hh}^+(2-\mathbf{t}_1))\rightarrow \cl{O}_{\hg\hh}$ be the Perrin-Riou map obtained from $\cl{L}_{\hf\hg\hh}$ by specializing $\hf$ to $f$. It is given by the composition of a big logarithm map $\log_{f \hg \hh}^{-++}$ with the map obtained by pairing with $\eta_{f_\alpha}^\alpha\otimes \bm{\omega}_{\hg}\otimes \bm{\omega}_{\hh}$. Note that the pairing is taken with respect to the differential $\eta_{f_\alpha}^\alpha$ corresponding to the $p$-stabilized form $f_\alpha$, as opposed to $\eta_f^\alpha$.

\begin{defi}\label{def:big-log}
We denote by $\Log_{\bm{\omega}_{\hg} \otimes \bm{\omega}_{\hh}}$ the $\Lambda_{\hg\hh}$-module homomorphism
\[
\Log_{\bm{\omega}_{\hg} \otimes \bm{\omega}_{\hh}}=\cl{L}^{-++}_{f\hg\hh}\circ\mathrm{pr}^{-++}\circ\res_p : H^1_{\bal}(\bb{Q},\bb{V}_{f\hg\hh}^\dagger)\longrightarrow  \cl{O}_{\hg\hh}.
\]
%application \[ \Log_{\omega_{\hg} \otimes \omega_{\hh}} \colon H^1(\mathbb Q_p) \mathbb V_f^- \otimes \mathbb V_{\hg}^+ \otimes \mathbb V_{\hh}^+) \longrightarrow \Lambda_{\hg \hh} := \Lambda_{\hg} \otimes \Lambda_{\hh} \] given by the pairing with the canonical differentials $\eta_f \otimes \omega_{\hg} \otimes \omega_{\hh}$.

%\textcolor{blue}{Tal y como est\'a definido no tiene mucho sentido. La localizaci\'on en $p$ de un elemento arbitrario de $H^1(\bb{Q},\bb{V}_{\hg\hh})$ vive en $H^1(\bb{Q}_p,\bb{V}_{\hg\hh})$. Para poder definir el logaritmo hay que garantizar que vive en $H^1(\bb{Q}_p,\bb{V}_\hg^+\otimes \bb{V}_{\hh}^+)$.} \textcolor{red}{Si si, toda la razon, pero ahora entonces me estoy liando. En BS es cierto tambien que lo definen tomando tambien la parte +, pero eso en general no va a ser cierto, que nos caiga en la $(+,+)$, asi que solo estaria definido si las otras dos proyecciones, la $(+,-)$ y la $(-,+)$ se anulan, no?} \textcolor{blue}{La cuesti\'on es que una vez demuestras que $\BF$ vive en $H^1_{\bal}(\bb{Q},\bb{V}_{f\hg\hh}^\dagger)$, puedes tomar la proyecci\'on $(-,+,+)$ y aplicar $\log_{\eta_f \otimes\bm{\omega}_{\hg}\otimes \bm{\omega}_{\hh}}$, que viene a ser el homomorfismo $\cl{L}_{f\hg\hh}^f$ de la secci\'on 4.2. O sea, quieres coger el corolario 5.10 y aplicar $\cl{L}_{f\hg\hh}^f$ a ambos lados. En realidad lo que estamos haciendo en esta secci\'on no deja de ser un corolario casi inmediato del Corolario 5.10.} \textcolor{red}{Si si, totalmente, pues igual se puede poner alli como Corolario despues de 5.10, no? Luego me quede yo pensandolo y es cierto que es bastante tonto, pero es cierto que puede quedar gracioso (pero igual no haciendo una seccion propia).} \textcolor{blue}{S\'i, quiz\'a quede mejor ponerlo despu\'es del Corolario 5.10. Me parece bien.}
\end{defi}

\begin{remark}\label{rk:other-proj}
In terms of the Beilinson--Flach classes, we can think that the pairing with $\eta_{\hg} \otimes \omega_{\hh}$ (resp. $\omega_{\hg} \otimes \eta_{\hh}$) allows us to recover the $p$-adic $L$-function $L_p(\hg,\hh)$ (resp. $L_p(\hh,\hg)$), while here we are considering instead the pairing with $\omega_{\hg} \otimes \omega_{\hh}$.
\end{remark}

Let $\Lp^f(f,\hg,\hh)\in \cl{O}_{\hg\hh}$ be the two-variable $p$-adic $L$-function obtained from the three-variable $f$-dominant $p$-adic $L$-function $\Lp^f(\hf,\hg,\hh)$ introduced in \S\ref{subsec:tripleproduct} by specializing $\hf$ to $f$. In this subsection, we relate the $p$-adic $L$-function $\Lp^f(f,\hg,\hh)$ to the weighted Beilinson--Flach class $\BF(f,\hg,\hh)$. % (Parece repetitivo poner esto aqu\'i. Ya lo hemos recordado al principio de la secci\'on) As a piece of notation, recall that $\hf$ is the Hida family passing through the irregular weight one Eisenstein series $f=\Eis_1(\varepsilon_K)$ introduced in \S\ref{subsec:wt1}.
Note that, while $\Lp^g(f, \hg, \hh)^2$ and $\Lp^h(f, \hg, \hh)^2$ may be factored as the product of two Hida--Rankin $p$-adic $L$-functions by $p$-adic Artin formalism, %\textcolor{blue}{No veo muy bien c\'omo factorizas esas funciones como productos de BDPs} \textcolor{red}{Perdon, queria decir Hida--Rankin, no BDP (es decir, que si no es la dominante factorizas como producto de Hida--Rankins},
the case where the dominant $p$-adic $L$-function is Eisenstein is subtler, as already shown in \cite{BC}. %We write $\Lp^f(f,\hg,\hh)$ for the two-variable $p$-adic $L$-function corresponding to the specialization of $\Lp^f(\hf,\hg,\hh)$ in the first argument.

\begin{proposition}
Assume that $H^1_{\cl{G}\cup +}(\bb{Q}, \bb{V}_{f\hg\hh}^\dagger)$ is a torsion-free $\Lambda_{\hg\hh}$-module of rank $1$. Then we have the following factorization of the image of the Beilinson--Flach class under the Perrin-Riou map $\Log_{\bm{\omega}_{\hg} \otimes \bm{\omega}_{\hh}}$: %\textcolor{blue}{No creo que tenga sentido llamarle factorizaci\'on de la triple product $p$-adic $L$-function. En todo caso lo que estamos factorizando es el otro lado.} 
\[ \Omega_{f,\gamma} \cdot \Lp^g(f,\hg,\hh) \cdot \Lp^f(f, \hg, \hh) =  \lambda_{N_g}(\hg) \cdot \Log_{\omega_{\hg} \otimes \omega_{\hh}}(\BF(f, \hg, \hh)). \] 
\end{proposition}

\begin{proof}
By Corollary~\ref{corollary:comparison}, we have
\[
\lambda_{N_g}(\hg)\cdot\BF(f,\hg,\hh)=\Omega_{f,\gamma}\cdot \Lp^g(f,\hg,\hh) \cdot \kappa(f,\hg,\hh).
\]
The result then follows immediately by taking $\Log_{\bm{\omega}_{\hg} \otimes \bm{\omega}_{\hh}}$ and using Theorem~\ref{thm:reclawdiagonalcycles}. 
%We may consider, as in \cite{BC}, the big-logarithm map of Definition \ref{def:big-log} that corresponds to the projection to the $(-,+,+)$-component, thus yielding the desired equality: in the left hand side, this is just the $f$-dominant $p$-adic $L$-function, while in the right hand side it corresponds with the big Beilinson--Kato map introduced earlier on.
\end{proof}

\begin{comment}
\begin{remark}
In terms of Remark \ref{rk:other-proj}, we can think about this result as an expression for the pairing with $\omega_{\hg} \otimes \omega_{\hh}$, which is encoded not in a Hida--Rankin $p$-adic $L$-function, but in a triple product one. We hope to expand this idea in forthcoming work in the setting of the critical Eisenstein degeneration.
\end{remark}
\end{comment}

\section{The $p$-exceptional case}\label{sec:expectional}

We keep the notations and assumptions in the previous section, and consider now the situation where there is a trivial zero for the $p$-adic $L$-function $\Lp^g(f,\hg,\hh)$, as well as for the $p$-adic $L$-functions $L_p(\hg,\hh)$ and $L_p(\hg,\hh\otimes\varepsilon_K)$, at the point $(y_0,z_0)$ corresponding to the modular forms $(g,h)$ arising from the vanishing of an Euler factor at $p$. In this situation, the statements in \S\ref{subsec:nonexceptional} become trivial, so we introduce \textit{improved} classes and \textit{improved} $p$-adic $L$-functions to remove the vanishing Euler factor. %We note that this case is indeed rather relevant since it arises naturally in certain cases of great arithmetic interest, as we will remark later.  \textcolor{blue}{Hacer comentario m\'as adelante con ejemplo de caso donde puede ser relevante. Seguramente algo del estilo: $h$ forma modular de peso $1$, $g$ forma asociada a curva el\'iptica con reducci\'on multiplicativa en $p$.)} 

Hence, along this section we will work under the following assumption.

\begin{ass}
With the previous notations, it holds that \[ \alpha_g \beta_h = p^{\frac{l_0+m_0-3}{2}}. \] 
\end{ass}

Since $g$ and $h$ are $p$-ordinary and $p\nmid \mathrm{cond}(h)$, it follows from this assumption together with the Ramanujan--Petersson conjecture that $(l_0,m_0)=(2,1)$ and that $g$ has conductor $N_gp$. Note that this situation includes in particular the case where $g$ is the modular form attached to an elliptic curve over $\bb{Q}$ with multiplicative reduction at $p$, which is of great arithmetic interest.


%\textcolor{red}{En realidad creo que podemos decir m\'as que eso. Viendo la discusi\'on  en la intro de BSV parece que en el caso ordinario la \'unica opci\'on es $l_0=2$ y $m_0=1$. Usando Ramanujan sale f\'acil que no puede ser $m_0\geq 3$. No veo tan claro c\'omo descartan $l_0=3$, $m_0=2$.}
%\textcolor{blue}{Yo igual diria entonces que por Ramanujan se tiene que $m_0 \geq 3$ y no me complicaria con el otro caso.} \textcolor{red}{No, lo que se tiene es que $m_0=1$ o $m_0=2$. Y viendo BSV parece que la \'unica opci\'on es $m_0=1$. Entonces no s\'e si ser\'ia conveniente decir eso, que en esta secci\'on $l_0=2$ y $m_0=1$.}

%{\color{blue} Ya he entendido lo que pasa (creo): hay una doble anulacion en la funcion $L$, y uno de los factores que se anula viene de la clase de cohomologia y otro del Perrin-Riou map.}

We consider the following codimension-$1$ subvariety of $\cl{U}_{\hg}\times \cl{U}_{\hh}$:
\[
\cl{Y}=\lbrace (y,z)\in \cl{U}_{\hg}\times\cl{U}_{\hh}\;\vert\; \kappa_{\hg}(y)=\kappa_{\hh}(z)+1\rbrace.
\]
Let $\cl{C}$ be an irreducible component of $\cl{Y}$ containing the point $(y_0,z_0)$. We denote by $\cl{O}_{\cl{C}}$ the corresponding ring of global functions. Note that an arithmetic point $(y,z)\in \cl{C}\subset \cl{U}_\hg\times\cl{U}_\hh$ will have weights $(l,l-1)$ for some integer $l$. We denote by $\cl{C}^{\mathrm{cl}}$ the intersection of $\cl{C}$ with $\tilde{\cl{W}}_{\hg}^{\mathrm{cl}}\times \tilde{\cl{W}}_{\hh}^{\mathrm{cl}}$.
%Define the line
%\begin{equation}
%\mathcal C := \mathcal S_{l,l-1,l-1} := \{ (y,z,\sigma) \in \mathcal W_{\hg} \times \mathcal W_{\hh} \times \mathcal W: \, \mathrm{w}(z) = \mathrm{w}(y)-1, \, \mathrm{w}(z)=\sigma \},
%\end{equation}
%where $\mathrm{w}$ is either the weight map $\mathrm{w} \colon \mathcal W_{\hg} \rightarrow \mathcal W$ or $\mathrm{w} \colon \mathcal W_{\hh} \rightarrow \mathcal W$. 
%Note that this is a sub-variety of $ \mathcal W_{{\bf ghs}}$ all whose crystalline arithmetic points have weights $(l,l-1,l-1)$ for some $l \geq 1$.
%\textcolor{red}{Yo quiz\'a definir\'ia $\cl{C}$ como subvariedad de $\cl{W}_{\hg}\times\cl{W}_{\hh}$. De lo contrario luego no tiene mucho sentido restringir $\cl{O}_{\hf\hg\hh}$ a $\cl{C}$ (o al menos habr\'ia que explicar lo que eso significa).} \textcolor{blue}{Ok, me parece bien.}


Note that, for any point $(y,z)\in \cl{C}^{\mathrm{cl}}$ with $p\nmid \mathrm{cond}(\hg_y)\cdot\mathrm{cond}(\hh_z)$, the Euler factor $\cl{E}(x_0,y,z)$ appearing in the interpolation formula for $\Lp^g(f,\hg,\hh)$ and the Euler factor $\cl{E}(y,z,s(y,z))$ appearing in the interpolation formula for $L_p(\hg,\hh,(\mathbf{l}+\mathbf{m}-1)/2)$ are given by
\[
\cl{E}(x_0,y,z)=\cl{E}(y,z,s(y,z))=\left(1-\chi_g(p)\frac{\alpha_{\hh_z}}{\alpha_{\hg_y}}\right)^2 \left(1-\chi_g(p)\frac{\beta_{\hh_z}}{\alpha_{\hg_y}}\right)^2.
\]
For varying values of $(y,z)\in \cl{C}^{\mathrm{cl}}$, the factors $1-\chi_g(p)\alpha_{\hh_z}\alpha_{\hg_y}^{-1}$ are interpolated by the $p$-adic $L$-function
\[
1-\chi_g(p)\frac{a_p(\hh)}{a_p(\hg)} \in \cl{O}_{\cl{C}},
\]
which vanishes at the point $(y_0,z_0)$.

%The restriction to $\mathcal C$ of the {\em second multiplier} in the Euler-like factor appearing in the interpolation formula for Hida--Rankin's $p$-adic $L$-function is \[ 1-\frac{p^{s-1}}{\alpha_{g_y^{\circ}} \beta_{h_z^{\circ}}} = 1- \frac{\chi_h^{-1}(p) \alpha_{h_z^{\circ}}}{\alpha_{g_y^{\circ}}}. \] This expression interpolates to an Iwasawa function in $\Lambda_{\hg}$, which by abuse of notation we continue to denote with the same symbol. Similarly, the {\em third multiplier} is \[ 1 - \frac{\beta_{g_y^{\circ}} \alpha_{h_z^{\circ}}}{p^s} = 1 - \frac{\chi_g(p) \alpha_{h_z^{\circ}}}{\alpha_{g_y^{\circ}}} p^{l-s-1} = 1 - \frac{\chi_g(p) \alpha_{h_z^{\circ}}}{\alpha_{g_y^{\circ}}}. \] Hence, the extra condition that $2s=l+m-1$, the vanishing of the second multiplier is equivalent to the vanishing of the third multiplier. Observe that $\chi_h^{-1}(p) = \chi_g(p)$ since $\varepsilon_K(p)=1$.

\begin{remark}
The appearance of the double factor
\[
\left(1-\chi_g(p)\frac{a_p(\hh)}{a_p(\hg)}\right)^2
\]
in the interpolation formula for $\Lp^g(f,\hg,\hh)$ (resp. $L_p(\hg,\hh,(\mathbf{l}+\mathbf{m}-1)/2)$) at points in $\cl{C}^{\mathrm{cl}}$ comes from two different sources:
\begin{itemize}
    \item the numerator in the interpolation formula for the Perrin-Riou map $\cl{L}_{f\hg\hh}^{+-+}$ (resp. $\cl{L}_{\hg\hh}^{-+})$;
    \item the comparison between the corresponding $p$-stabilized and non-$p$-stabilized cohomology classes.
\end{itemize}
%Now, we make the following distinction:
%\begin{itemize}
%\item[-] The second multiplier corresponds to the numerator of the Perrin-Riou map. In this setting, we will therefore work with the so-called improved Perrin-Riou map.
%\item[-] The third multiplier corresponds to a factor which is also present in the cohomology class (but not in the Perrin-Riou map). Working assuming that this is zero would be the analogue of \cite{BDV}.
%\end{itemize}
\end{remark}

%One naturally expects $1-\frac{\alpha_{\chi_h(p) h_z^{\circ}}}{\alpha_{g_y^{\circ}}}$ should divide the two-variable $p$-adic $L$-function $L_p(\hg,\hh)(y,z,s)$, where $(y,z)$ vary in $\mathcal W_{\hg} \times \mathcal W_{\hh}$ and $s=\mathrm{w}(z)$ is determined by the weight of $z$. Hida's work \cite{hida2} allows us to conclude the following when $\hh \neq \hg^*$.

\subsection{Improved $p$-adic $L$-functions}

In this subsection we introduce the improved $p$-adic $L$-functions that we will use.

%We denote by $L_p(\hg,\hh)_{\vert \cl{C}}\in I_{\hg}^{-1}\cl{O}_{\cl{C}}$ the restriction of the $p$-adic $L$-function $L_p(\hg,\hh)$ to $\cl{C}$. Similarly, we define $L_p(\hg,\hh\otimes\varepsilon_K)_{\vert\cl{C}}\in I_{\hg}^{-1}\cl{O}_{\cl{C}}$ and $L_p(\hh,\hg)_{\vert\cl{C}}, L_p(\hh\otimes\varepsilon_K, \hg)_{\vert\cl{C}}\in I_{\hh}^{-1}\cl{O}_{\cl{C}}$.

\begin{theorem}\label{thm:improvedHRpadicLfunction}
    There exists a $p$-adic $L$-function
    \[
    \widehat{L}_p(\hg,\hh)\in I_{\hg}^{-1}\cl{O}_{\cl{C}}
    \]
    such that
    \[
    L_p(\hg,\hh,(\mathbf{l}+\mathbf{m}-1)/2)\vert_{\cl{C}}=\left(1-\chi_g(p)\frac{a_p(\hh)}{a_p(\hg)}\right)^2 \widehat{L}_p(\hg,\hh).
    \]
    % (quito esto porque no estoy del todo seguro de c\'omo justificar (ii), aunque creo que deber\'ia ser verdad) satisfying the following properties:
    %\begin{enumerate}[(i)]
        %\item for every $(y,z)\in\cl{C}^{\mathrm{cl}}$ of weights $(l,l-1)$ with $p\nmid \mathrm{cond}(\hg_y)\cdot \mathrm{cond}(\hh_z)$ we have
        %\[
        %\widehat{L}_p(\hg,\hh)(y,z) = \frac{\mathcal{E}^\ast(y,z)}{\mathcal{E}_0(y)\mathcal{E}_1(y)} \frac{\Gamma(l-1)}{\pi^{l}(-i)2^{2l-1} \langle \hg_y^\circ,\hg_y^\circ \rangle}  \times L(\hg_y^\circ,\hh_z^\circ,l-1),
        %\]
        %where
        %\begin{eqnarray}
        %\nonumber \cl{E}_0(y) &:=& 1-\chi_g^{-1}(p) \beta_{\hg_y}^2p^{1-l}, \\
        %\nonumber \mathcal E_1(y) &:=& 1-\chi_g(p)\alpha_{\hg_y}^{-2}p^{l-2}, \\
        %\nonumber \mathcal E^\ast(y,z) &:=& \Big(1-\chi_g(p)\frac{\beta_{\hh_z}}{\alpha_{\hg_y}} \Big)^2;
        %\end{eqnarray}
     %   \item $L_p(\hg,\hh,(\mathbf{l}+\mathbf{m}-1)/2)=\left(1-\chi_g(p)\frac{a_p(\hh)}{a_p(\hg)}\right)^2 \widehat{L}_p(\hg,\hh)$;
      %  \item $\widehat{L}_p(\hg,\hh)(y_0,z_0)=0$ if and only if $L(g,h,1)=0$.
    %\end{enumerate}
\end{theorem}

%\begin{remark}
 %   Since the set of points $(y,z)\in\cl{C}^{\mathrm{cl}}$ with $p\nmid \mathrm{cond}(\hg_y)\cdot \mathrm{cond}(\hh_z)$ is Zariski dense in $\cl{C}$, it follows immediately from the interpolation formula that
  %  \[
   % L_p(\hg,\hh)=\left(1-\chi_g(p)\frac{a_p(\hh)}{a_p(\hg)}\right)^2 \widehat{L}_p(\hg,\hh). 
    %\]
%\end{remark}

%The next result is related with \cite{hida2}, where he deletes one of the Euler factors passing to a codimension one subvariety. Now, we delete two of them, and we have to pass to a codimension two subvariety.

%\begin{theorem}\label{improved}
%Let $\hat L_p(\hg,\hh)$ be the unique element in the fraction field of $\Lambda_{\hg} \hat \otimes \Lambda_{\hh}$ such that \[ L_p(\hg,\hh)(y,z,\mathrm{w}(z)) = \Big(1-\frac{\chi_h(p) \alpha_{h_z^{\circ}}}{\alpha_{g_y^{\circ}}} \Big)^2 \cdot \hat L_p(\hg, \hh)(y,z). \] {\color{blue} This function is analytic in $\mathcal C$}.
%\end{theorem}

\begin{proof}
Let $N=\lcm(N_g,N_h)$. Let $E_1=E^{(1)}_{1/N}(\tau)$ be the weight-$1$ holomorphic Eisenstein series introduced in \cite[\S5.1]{LLZ0}. Then, as in the proof of \cite[Lemma~9.8]{BSV}, one shows that the $p$-adic $L$-function
%Following the construction in \cite{hida2} (\emph{cf}. also the proof of \cite[Lemma~9.8]{BSV}), for a suitable holomorphic weight-$1$ Eisenstein series $E_1$, the $p$-adic $L$-function
\[
\widehat{L}_p(\hg,\hh)=\frac{\langle W_{N}(\hg), \mathrm{e}_{\ord}(E_1\cdot \hh)\rangle_N}{\langle \hg,\hg\rangle_N} \in I_{\hg}^{-1} \cl{O}_{\cl{C}}
\]
satisfies the required properties.

%{\color{blue} Esto es algo que hacen BSV en algun articulo, de decir que se pueden quitar dos factores de Euler restringiendo cada vez a una subvariedad de dimension menor. Obviamente se puede hacer, pero hay que decir por que no hay polos. No se si lo ves claro, sino igual se puede citar el BSV que toca.} \textcolor{red}{Una cosa, si sabemos que tenemos una clase mejorada y un Perrin-Riou map mejorado, no podemos definir $\hat{L}_p$ como la imagen de la clase mejorada por el Perrin-Riou map mejorado? Haci\'endolo as\'i no ser\'ia inmediato el resultado? Y no hace falta restringirse a codimensi\'on dos.} \textcolor{blue}{Si, pero luego igual estaria bien poder relacionar esa funcion $\hat L_p$ con la otra, no? que igual no es dificil, o como lo ves?} \textcolor{red}{Si lo \'unico que queremos decir es lo que est\'a enunciado en el teorema puedes simplemente tomar $\hat{L}_p=\hat{\cl{L}}(\hat{\kappa})$ y tienes que $L_p=\cl{L}(\kappa)=\cl{E} \cdot\hat{\cl{L}}(\cl{E}\cdot\hat{\kappa})=\cl{E}^2\cdot \hat{L}_p$. Supongo que la cuesti\'on es que queremos decir algo m\'as sobre $\hat{L}_p$.} \textcolor{blue}{Igual podemos hacerlo en principio asi y luego ver si es posible decir algo mas.}
\end{proof}

%\begin{remark}
%The point where the vanishing occurs satisfies $m \leq s < l$, since we have the conditions $2s=l+m-1$ and $l=m+1$. Hence, we do know that the $p$-adic $L$-functions is zero, since it is the product of the complex $L$-value by a vanishing Euler factor.
%\end{remark}


%{\color{blue} Introducir de forma similar la improved para el producto triple (ajustar bien segun como pongamos la notacion).}

%Following \cite[\S8.4]{Hs}, there is also an improved triple product $p$-adic $L$-function.


\begin{theorem}\label{thm:improvedtripleproductpadicLfunction}
    There exists a $p$-adic $L$-function
    \[
    \widehat{\Lp}^g(f,\hg,\hh)\in I_{\hg}^{-1}\cl{O}_{\cl{C}}
    \]
    such that
    \[
    \Lp^g(f,\hg,\hh)\vert_{\cl{C}}=\left(1-\chi_g(p)\frac{a_p(\hh)}{a_p(\hg)}\right)^2 \widehat{\Lp}^g(f,\hg,\hh).
    \]
    %satisfying the following properties:
    %\begin{enumerate}[(i)]
        %\item for every $(y,z)\in\cl{C}^{\mathrm{cl}}$ of weights $(l,l-1)$ with $p\nmid \mathrm{cond}(\hg_y)\cdot \mathrm{cond}(\hh_z)$ we have
        %\[
        %\widehat{\Lp}(f,\hg,\hh)^2(y,z) = \frac{\mathcal{E}^\ast(y,z)^2}{\mathcal{E}_0(y)^2\mathcal{E}_1(y)^2} \frac{\Gamma(l-1)^2}{(2\pi i)^{2l-4} \langle \hg_y^\circ,\hg_y^\circ \rangle}  \times L(f,\hg_y^\circ,\hh_z^\circ,l-1),
        %\]
        %where
        %\begin{eqnarray}
        %\nonumber \cl{E}_0(y) &:=& 1-\chi_g^{-1}(p) \beta_{\hg_y}^2p^{1-l}, \\
        %\nonumber \mathcal E_1(y) &:=& 1-\chi_g(p)\alpha_{\hg_y}^{-2}p^{l-2}, \\
        %\nonumber \mathcal E^\ast(y,z) &:=& \Big(1-\chi_g(p)\frac{\beta_{\hh_z}}{\alpha_{\hg_y}} \Big)^2;
        %\end{eqnarray}
     %   \item $\Lp^g(f,\hg,\hh)\vert_{\cl{C}}=\left(1-\chi_g(p)\frac{a_p(\hh)}{a_p(\hg)}\right)^2 \widehat{\Lp}(f,\hg,\hh)$;
     %   \item $\widehat{\Lp}(f,\hg,\hh)(y_0,z_0)=0$ if and only if $L(f,g,h,1)=0$.
    %\end{enumerate}
\end{theorem}
\begin{proof}
    This is proved as in \cite[Lemma~9.8]{BSV}, taking $(\hf,\hg,\hh)$ in \emph{loc.\,cit.} to be our $(\hg,\hh,\hf)$.%, and proceeding as in \cite{Hs} to find a suitable choice of test vectors.
\end{proof}

%\begin{remark}
 %   Since the set of points $(y,z)\in\cl{C}^{\mathrm{cl}}$ with $p\nmid \mathrm{cond}(\hg_y)\cdot \mathrm{cond}(\hh_z)$ is Zariski dense in $\cl{C}$, it follows immediately from the interpolation formula that
  %  \[
   % \widehat{\Lp}^g(f,\hg,\hh)=\left(1-\chi_g(p)\frac{a_p(\hh)}{a_p(\hg)}\right)^2 \widehat{L}_p(\hg,\hh). 
    %\]
%\end{remark}

%\begin{theorem}[Hsieh]
%There is an analytic $p$-adic $L$-function $\widehat{\Lp}^g(\hf,\hg,\hh)$ defined along the line $\mathcal C$ and satisfying that \[ \Lp^g(\hf,\hg,\hh) =  \left( 1 - \frac{\chi_g(p) \alpha_{f_x^{\circ}} \alpha_{h_z^{\circ}}}{\alpha_{g_y^{\circ}}} \right)^2 \widehat{\Lp}^g(\hf,\hg,\hh) \]
%\end{theorem}

%\begin{proof}
%{\color{blue} Comentar algo.}
%\end{proof}

\subsection{Improved Perrin-Riou maps}\label{subsec:improvedPerrinRiou}

%{\color{blue} Esta parte es la que mas hay que repasar, porque ahora lo que queremos usar es el improved Perrin-Riou map, que no tiene el factor de Euler que se anula.}

In this subsection we introduce improved Perrin-Riou maps.  Given a $\Lambda_{\hg\hh}$-module $M$, we denote by $M\vert_{\cl{C}}$ the $\cl{O}_{\cl{C}}$-module $M\otimes_{\Lambda_{\hg\hh}}\cl{O}_{\cl{C}}$.

\begin{propo}\label{perrin-improved1}
There exists a homomorphism of $\mathcal{O}_{\cl{C}}$-modules
\[
\widehat{\mathcal L}_{\hg \hh}^{-+}: H^1(\mathbb{Q}_p,\mathbb{V}_{\hg}^- \hat{\otimes} \mathbb{V}_{\hh}^+ (2-\mathbf{t}_1)|_{\mathcal C}) \rightarrow I_\hg^{-1} \cl{O}_{\cl{C}},
\]
such that for every point $(y,z) \in \mathcal{C}^{\mathrm{cl}}$ of weights $(l,l-1)$ the specialization of $\widehat{\mathcal L}_{\hg \hh}^{-+}$ at $(y,z)$ is the homomorphism
\[
\widehat{\mathcal L}_{{\hg\hh}}^{-+}(y,z): H^1(\mathbb Q_p, V_{\hg_y}^- \otimes V_{\hh_z}^+(2-l)) \rightarrow \mathbb C_p
\]
given by
\[
\widehat{\mathcal L}_{{\hg \hh}}^{-+}(y,z) = \frac{1}{1-p^{1-l} \alpha_{\hg_y} \beta_{\hh_z}}  \times \langle \exp_{\BK}^*(\cdot), \omega_{f} \otimes\eta_{\hg_y}^\alpha \otimes \omega_{\hh_z} \rangle.
\]
\end{propo}
\begin{proof}
This follows from \cite[Lemma~9.4]{BSV}.
\end{proof}


\begin{propo}\label{perrin-improved2}
There exists a homomorphism of $\mathcal{O}_{\cl{C}}$-modules
\[
\widehat{\mathcal L}_{f \hg \hh}^{+-+}: H^1(\mathbb{Q}_p, V_f^+\otimes\mathbb{V}_{\hg}^- \hat{\otimes} \mathbb{V}_{\hh}^+ (2-\mathbf{t}_1)|_{\mathcal C}) \rightarrow I_\hg^{-1} \cl{O}_{\cl{C}},
\]
such that for every point $(y,z) \in \mathcal{C}^{\mathrm{cl}}$ of weights $(l,l-1)$ the specialization of $\widehat{\mathcal L}_{f\hg \hh}^{+-+}$ at $(y,z)$ is the homomorphism
\[
\widehat{\mathcal L}_{{f\hg\hh}}^{+-+}(y,z): H^1(\mathbb Q_p, V_f^+\otimes V_{\hg_y}^- \otimes V_{\hh_z}^+(2-l)) \rightarrow \mathbb C_p
\]
given by
\[
\widehat{\mathcal L}_{{f\hg \hh}}^{+-+}(y,z) = \frac{1}{1-p^{1-l} \alpha_{\hg_y} \beta_{\hh_z}}  \times \langle \exp_{\BK}^*(\cdot), \omega_{f} \otimes\eta_{\hg_y}^\alpha \otimes \omega_{\hh_z} \rangle.
\]
\end{propo}
\begin{proof}
This follows from \cite[Lemma~9.4]{BSV}.
\end{proof}

\begin{remark}
As before, given an element
\[
v\otimes z\in V_f^+\otimes H^1(\mathbb{Q}_p,\mathbb{V}_{\hg}^- \hat{\otimes} \mathbb{V}_{\hh}^+ (2-\mathbf{t}_1)|_{\mathcal C}) \cong H^1(\mathbb{Q}_p, V_f^+\otimes\mathbb{V}_{\hg}^- \hat{\otimes} \mathbb{V}_{\hh}^+ (2-\mathbf{t}_1)|_{\mathcal C}),
\]
we have that
\[
\widehat{\cl{L}}_{f\hg\hh}^{+-+}(v\otimes z)=\langle v,\omega_f\rangle\cdot \widehat{\cl{L}}_{\hg\hh}^{+-}(z).
\]
\end{remark}

\subsection{Improved cohomology classes}

We now introduce the improved version of the cohomology classes that have been used in this work. %We consider the $\cl{O}_{\cl{C}}[G_{\bb{Q}}]$-modules $\bb{V}_{\hg\hh\vert\cl{C}}^\dagger=\bb{V}_{\hg\hh}\otimes_{\Lambda_{\hg\hh}}\cl{O}_{\cl{C}}$ and $\bb{V}_{f\hg\hh\vert\cl{C}}^\dagger=\bb{V}_{f\hg\hh}\otimes_{\Lambda_{\hg\hh}}\cl{O}_{\cl{C}}$.

We can define Selmer conditions for the $\cl{O}_{\cl{C}}[G_{\bb{Q}}]$-modules $\bb{V}_{\hg\hh}^\dagger\vert_{\cl{C}}$ and $\bb{V}_{f\hg\hh}^\dagger\vert_{\cl{C}}$ analogous to the ones introduced for $\bb{V}_{\hg\hh}^\dagger$ and $\bb{V}_{f\hg\hh}^\dagger$ in the previous section and consider the corresponding Selmer groups. In particular, we can consider the Selmer groups $H^1_{\bal}(\bb{Q},\bb{V}_{\hg\hh}^\dagger\vert_{\cl{C}})$, $H^1_{\bal}(\bb{Q},\bb{V}_{f\hg\hh}^\dagger\vert_{\cl{C}})$ and $H^1_{\cl{G}\cup +}(\bb{Q},\bb{V}_{f\hg\hh}^\dagger\vert_{\cl{C}})$.

While in the case of diagonal cycles there is a construction of an improved cohomology class, this is not the case for Beilinson--Flach elements. As it occurs with the $p$-adic $L$-function, the factor $\left( 1 - \chi_g(p)\frac{a_p(\hh)}{a_p(\hg)} \right)$ appears in the interpolation property when considering its variation in families (see e.g. \cite[\S8]{KLZ}). Hence, the following conjecture can be seen as a standard expectation in the theory of exceptional zeros, and we expect to come back to it in forthcoming work.

\begin{conj}\label{conj:improvedbeilinsonclass}
There exists a cohomology class $\widehat{\kappa}_{\hg,\hh}\in H^1_{\bal}(\bb{Q},\bb{V}_{\hg\hh}^\dagger\vert_{\cl{C}})$ such that
\[
\kappa_{\hg,\hh}\vert_{\cl{C}} = \left( 1 - \chi_g(p)\frac{a_p(\hh)}{a_p(\hg)} \right) \widehat{\kappa}_{\hg,\hh}.
\]
\end{conj}

\begin{comment}
{\color{blue} Del estilo de BSV para los ciclos, pero comentar algo.} \textcolor{red}{No veo muy claro c\'omo justificar esto. Entiendo que habr\'ia que meterse en la construcci\'on de KLZ y ver si bajo la condici\'on $\mathbf{s}=\mathbf{l}-1$ se puede sacar un factor de Euler de la clase, pero la construcci\'on es un poco farragosa y no parece obvio. Entiendo que moralmente queremos decir que podemos construir una clase que interpola los $\Eis^{[f_0,g,j]}$ en vez de los $\Eis^{[f,g,j]}$.} 

\textcolor{blue}{Yo estaba dandole vueltas e inicialmente tambien pensaba que era mas claro (porque claro, en BSV consideran un improved pero sacando ese factor). Me extrana que en ningun sitio este enunciado un teorema asi, pero tampoco lo encuentro, pero es justo lo que dices, que moralmente tiene que ser una clase que interpola los $\Eis^{[f_0,g,j]}$, solo que no se como de inmediato es justificarlo.}
\end{comment}

\begin{remark}
    Similarly, we also expect to have an analogous improved cohomology class $\widehat{\kappa}_{\hg,\hh\otimes\varepsilon_K} \in H^1_{\bal}(\bb{Q},\bb{V}_{\hg\hh}^\dagger\vert_{\cl{C}}\otimes\varepsilon_K)$.
\end{remark}

For the rest of this section, we work under the following assumption.

\begin{ass}
    Conjecture~\ref{conj:improvedbeilinsonclass} holds.
\end{ass}

Let $\mathrm{pr}^{-+}:H^1_{\bal}(\bb{Q}_p,\bb{V}_{\hg\hh}^\dagger\vert_{\cl{C}})\rightarrow H^1(\bb{Q}_p,\bb{V}_{\hg}^-\hat{\otimes}\bb{V}_\hh^+(2-\mathbf{t}_1)\vert_{\cl{C}})$ be the map induced by the natural projection $\mathscr{F}^2\bb{V}_{\hg\hh}^\dagger\vert_{\cl{C}}\rightarrow \bb{V}_\hg^-\hat{\otimes}\bb{V}_\hh^+(2-\mathbf{t}_1)\vert_{\cl{C}}$ and let
\[
\widehat{\mathcal{L}}_{\hg\hh}^{\hg}:H^1_{\bal}(\bb{Q},\bb{V}_{\hg\hh}^\dagger\vert_{\cl{C}})\longrightarrow I_{\hg}^{-1}\cl{O}_{\cl{C}}
\]
be the map defined by $\widehat{\mathcal{L}}_{\hg\hh}^\hg=\widehat{\mathcal{L}}_{\hg\hh}^{-+}\circ \mathrm{pr}^{-+}\circ \res_p$.

\begin{propo}\label{prop:improvedreclaw1}
    The class $\widehat{\kappa}_{\hg,\hh}$ satisfies
    \[
    \widehat{\cl{L}}_{\hg\hh}^{\hg}(\widehat{\kappa}_{\hg,\hh})=\widehat{L}_p(\hg,\hh).
    \]
\end{propo}
\begin{proof}
   By Proposition~\ref{perrin-improved1}, Theorem~\ref{reclaw} and Theorem~\ref{thm:improvedHRpadicLfunction}, we have
    \begin{align*}
        \left( 1 - \chi_g(p)\frac{a_p(\hh)}{a_p(\hg)} \right)^2  \widehat{\cl{L}}_{\hg\hh}^\hg (\widehat{\kappa}_{\hg,\hh})&=\cl{L}_{\hg\hh}^{\hg}(\kappa_{\hg,\hh})\vert_{\cl{C}}= L_p(\hg,\hh)\vert_{\cl{C}} \\
        &=\left(1-\chi_g(p)\frac{a_p(\hh)}{a_p(\hg)}\right)^2 \widehat{L}_p(\hg,\hh).
    \end{align*}
    By the Ramanujan--Petersson conjecture, the factor $1-\chi_g(p)\frac{a_p(\hh)}{a_p(\hg)}$ does not vanish at any point $(y,z)\in\cl{C}^{\mathrm{cl}}$ with $p\nmid \mathrm{cond}(\hg_y)\cdot \mathrm{cond}(\hh_z)$, so in particular it is non-zero in $\cl{O}_{\cl{C}}$. The result now follows from the above equality. 
\end{proof}

\begin{remark}
    Similarly $\widehat{\cl{L}}_{\hg\hh}^{\hg}(\widehat{\kappa}_{\hg,\hh\otimes\varepsilon_K})=\widehat{L}_p(\hg,\hh\otimes\varepsilon_K)$.
\end{remark}

%Similarly, arguing as in \cite{BSV}, we can get an improved diagonal cycle class $\kappa_h(\hf,\hg,\hh)$, which agrees with the former up to multiplication by the corresponding Euler factor. We also write $\kappa_h(f,g,h)$ for the corresponding specialization.

\begin{propo}\label{prop:improvedtripleproductclass}
There exists a cohomology class $\widehat{\kappa}(f,\hg,\hh)\in H^1_{\bal}(\bb{Q},\bb{V}_{f\hg\hh}^\dagger\vert_{\cl{C}})$ such that
\[
\kappa(f,\hg,\hh)\vert_{\cl{C}} = \left( 1 - \chi_g(p)\frac{a_p(\hh)}{a_p(\hg)} \right) \widehat{\kappa}(f,\hg,\hh).
\]
\end{propo}
\begin{proof}
This is proved in \cite[\S9.3]{BSV}.
\end{proof}

Let $\mathrm{pr}^{+-+}:H^1_{\bal}(\bb{Q}_p,\bb{V}_{f\hg\hh}^\dagger\vert_{\cl{C}})\rightarrow H^1(\bb{Q}_p,V_f^+\otimes\bb{V}_{\hg}^-\hat{\otimes}\bb{V}_\hh^+(2-\mathbf{t}_1)\vert_{\cl{C}})$ be the map induced by the natural projection $\mathscr{F}^2\bb{V}_{f\hg\hh}^\dagger\vert_{\cl{C}}\rightarrow V_f^+\otimes\bb{V}_\hg^-\hat{\otimes}\bb{V}_\hh^+(2-\mathbf{t}_1)\vert_{\cl{C}}$ and let
\[
\widehat{\mathcal{L}}_{f\hg\hh}^{\hg}:H^1_{\bal}(\bb{Q},\bb{V}_{f\hg\hh}^\dagger\vert_{\cl{C}})\longrightarrow I_{\hg}^{-1}\cl{O}_{\cl{C}}
\]
be the map defined by $\widehat{\mathcal{L}}_{f\hg\hh}^\hg=\widehat{\mathcal{L}}_{f\hg\hh}^{+-+}\circ \mathrm{pr}^{+-+}\circ \res_p$.

\begin{propo}
    The class $\widehat{\kappa}(f,\hg,\hh)$ satisfies
    \[
    \widehat{\cl{L}}_{f\hg\hh}^{\hg}(\widehat{\kappa}(f,\hg,\hh))=\widehat{\Lp}^g(f,\hg,\hh).
    \]
\end{propo}
\begin{proof}
    By Proposition~\ref{perrin-improved2}, Proposition~\ref{prop:improvedtripleproductclass}, Theorem~\ref{thm:reclawdiagonalcycles} and Theorem~\ref{thm:improvedtripleproductpadicLfunction}, we have
    \begin{align*}
        \left( 1 - \chi_g(p)\frac{a_p(\hh)}{a_p(\hg)} \right)^2  \widehat{\cl{L}}_{f\hg\hh}^\hg (\widehat{\kappa}(f,\hg,\hh))&=\cl{L}_{f\hg\hh}^{\hg}(\kappa(f,\hg,\hh))\vert_{\cl{C}}= \Lp^g(f,\hg,\hh)\vert_{\cl{C}} \\
        &=\left(1-\chi_g(p)\frac{a_p(\hh)}{a_p(\hg)}\right)^2 \widehat{\Lp}(f,\hg,\hh).
    \end{align*}
    By the Ramanujan--Petersson conjecture, the factor $1-\chi_g(p)\frac{a_p(\hh)}{a_p(\hg)}$ does not vanish at any point $(y,z)\in\cl{C}^{\mathrm{cl}}$ with $p\nmid \mathrm{cond}(\hg_y)\cdot \mathrm{cond}(\hh_z)$, so in particular it is non-zero in $\cl{O}_{\cl{C}}$. The result now follows from the above equality.
\end{proof}



\subsection{The main theorem}

In this subsection, we adapt the main result of \S\ref{subsec:nonexceptional} to the current exceptional case.

\begin{propo}\label{prop:improvedfactorization}
    We have the following equality of $p$-adic $L$-functions:
    \[
    \widehat{\Lp}^g(f,\hg,\hh)^2=\widehat{L}_p(\hg,\hh)\widehat{L}_p(\hg,\hh\otimes\varepsilon_K).
    \]
\end{propo}
\begin{proof}
    By Proposition~\ref{prop:factorization}, Theorem~\ref{thm:improvedHRpadicLfunction} and Theorem~\ref{thm:improvedtripleproductpadicLfunction}, we have the equality
    \[
    \left(1-\chi_g(p)\frac{a_p(\hh)}{a_p(\hg)}\right)^4 \widehat{\Lp}^g(f,\hg,\hh)^2=\left(1-\chi_g(p)\frac{a_p(\hh)}{a_p(\hg)}\right)^4\widehat{L}_p(\hg,\hh)\widehat{L}_p(\hg,\hh\otimes\varepsilon_K)
    \]
    in $I_{\hg}^{-2}\cl{O}_{\cl{C}}$. By the Ramanujan--Petersson conjecture, the factor $1-\chi_g(p)\frac{a_p(\hh)}{a_p(\hg)}$ does not vanish at any point $(y,z)\in\cl{C}^{\mathrm{cl}}$ with $p\nmid \mathrm{cond}(\hg_y)\cdot \mathrm{cond}(\hh_z)$, so in particular it is non-zero in $\cl{O}_{\cl{C}}$. The result now follows from the above equality.
\end{proof}

We now introduce the improved Beilinson--Flach class that we will use. For the remaining of the section, we assume that Conjecture~\ref{conj:improvedbeilinsonclass} is true.

\begin{defi}
The \emph{improved weighted Beilinson--Flach class} associated with the pair $(\hg,\hh)$ is the element
\[
\widehat{\mathrm{BF}}(f,\hg,\hh)=v_{f,1}\otimes \widehat{L}_p(\hg,\hh\otimes\varepsilon_K)\widehat{\kappa}_{\hg,\hh}+v_{f,\varepsilon_K}\otimes\widehat{L}_p(\hg,\hh)\widehat{\kappa}_{\hg,\hh\otimes \varepsilon_K}
\]
in $H^1(\bb{Q},\bb{V}_{f\hg\hh}^\dagger\vert_{\cl{C}})$.   
\end{defi}

%Quito el remark porque creo que no es correcto. No creo que en este caso el improved Perrin-Riou map sea inyectivo (tampoco estoy del todo seguro de que lo sea en el otro caso).
%\begin{remark}
 %   It follows as in Proposition~\ref{prop:localconditionforBF}, using in this case the improved Perrin-Riou map introduced in Proposition~\ref{perrin-improved1} and the reciprocity law given by Proposition~\ref{prop:improvedreclaw1}, that $\widehat{\mathrm{BF}}(f,\hg,\hh)$ lies in $H^1_{\cl{G}\cup +}(\bb{Q},\bb{V}_{f\hg\hh}^\dagger\vert_{\cl{C}})$.
%\end{remark}

\begin{comment}
Arguing as in \cite{BSV}, we can get an improved diagonal cycle class $\kappa_h(\hf,\hg,\hh)$, which agrees with the former up to multiplication by the corresponding Euler factor. We also write $\kappa_h(f,g,h)$ for the corresponding specialization.

\begin{propo}
Under the current assumptions, there exists a cohomology class $\kappa_h(\hf,\hg,\hh)$ defined over the codimension one subvariety corresponding to the points with $l-1=m$, and such that \[ \kappa(\hf,\hg,\hh) = \left( 1 - \frac{\beta_{\hf} \beta_{\hg} \alpha_{\hh}}{p^c} \right) \kappa_h(\hf,\hg,\hh), \] where $c = \frac{k+l+m-2}{2}$.
\end{propo}

\begin{proof}
This is proved in \cite[\S9.3]{BSV}.
\end{proof}
\end{comment}

%We can now adapt our previous ideas to prove the main result in the exceptional setting. In this case, instead of considering the (parallel) variation in both $\hg$ and $\hh$, we have to restrict to the variation of one single variable.

Let $\res_p(\widehat{\kappa}(f,\hg,\hh))^{+-+}$ and $\res_p(\widehat{\BF}(f,\hg,\hh))^{+-+}$ be the images of $\res_p(\widehat{\kappa}(f,\hg,\hh))$ and $\res_p(\widehat{\BF}(f,\hg,\hh))$, respectively, in $H^1(\bb{Q}_p,V_f^+\otimes\bb{V}_\hg^-\hat{\otimes}\bb{V}_\hh^+(2-\mathbf{t}_1)\vert_{\cl{C}})$.

\begin{theorem}\label{thm:bf-vs-cycles-fam}
The equality
\[
\Omega_{f,\gamma} \cdot \widehat{\cl{L}}_{f\hg\hh}^{+-+}(\res_p(\widehat{\kappa}(f,\hg,\hh))^{+-+})^2 =\lambda_{N_g}(\hg) \cdot \widehat{\mathcal{L}}_{f\hg\hh}^{+-+}(\res_p(\widehat{\BF}(f, \hg, \hh))^{+-+})
\]
holds. %In particular, if $H^1_{\cl{G}\cup +}(\bb{Q},\bb{V}_{f\hg\hh}^\dagger\vert_{\cl{C}})$ is a torsion-free $\cl{O}_{\cl{C}}$-module of rank $1$ and the $p$-adic $L$-function $\widehat{\Lp}^g(f,\hg,\hh)$ is not identically zero, then $\widehat{\mathrm{BF}}(f, \hg, \hh)$ belongs to $H^1_{\bal}(\bb{Q},\bb{V}_{f\hg\hh}^\dagger\vert_{\cl{C}})$ and
%\[
%\Omega_{f,\gamma} \cdot \widehat{\Lp}^g(f,\hg,\hh) \cdot \widehat{\kappa}(f,\hg,\hh) = \lambda_{N_g}(\hg)\cdot\widehat{\mathrm{BF}}(f, \hg, \hh).
%\]
\end{theorem}
\begin{proof}
The result follows as in Theorem~\ref{thm:main}, working in this case with the improved Perrin-Riou maps.
\end{proof}


\begin{proposition}
    Assume that $H^1(\bb{Q},\bb{V}_{f\hg\hh}^\dagger\vert_{\cl{C}})$ is a torsion-free $\cl{O}_{\cl{C}}$-module and $H^1_{\cl{G}\cup +}(\bb{Q},\bb{V}_{f\hg\hh}^\dagger)$ is a torsion-free $\Lambda_{\hg\hh}$-module of rank $1$. Then $\widehat{\mathrm{BF}}(f, \hg, \hh)$ belongs to $H^1_{\bal}(\bb{Q},\bb{V}_{f\hg\hh}^\dagger\vert_{\cl{C}})$ and
    \[
    \Omega_{f,\gamma} \cdot \widehat{\Lp}^g(f,\hg,\hh) \cdot \widehat{\kappa}(f,\hg,\hh) = \lambda_{N_g}(\hg)\cdot\widehat{\mathrm{BF}}(f, \hg, \hh).
    \]
\end{proposition}
\begin{proof}
    By Corollary~\ref{corollary:comparison}, we have
    \[
    \left(1-\chi_g(p)\frac{a_p(\hh)}{a_p(\hg)}\right)^3\cdot\Omega_{f,\gamma} \cdot \widehat{\Lp}^g(f,\hg,\hh) \cdot \widehat{\kappa}(f,\hg,\hh) = \left(1-\chi_g(p)\frac{a_p(\hh)}{a_p(\hg)}\right)^3\cdot \lambda_{N_g}(\hg)\cdot\widehat{\mathrm{BF}}(f, \hg, \hh).
    \]
    Since $H^1(\bb{Q},\bb{V}_{f\hg\hh}^\dagger\vert_{\cl{C}})$ is a torsion-free $\cl{O}_{\cl{C}}$-module and $1-\chi_g(p)\frac{a_p(\hh)}{a_p(\hg)}$ is non-zero in $\cl{O}_{\cl{C}}$, it follows that
    \[
    \Omega_{f,\gamma} \cdot \widehat{\Lp}^g(f,\hg,\hh) \cdot \widehat{\kappa}(f,\hg,\hh) = \lambda_{N_g}(\hg)\cdot\widehat{\mathrm{BF}}(f, \hg, \hh).
    \]
    In particular, since $\widehat{\kappa}(f,\hg,\hh)$ belongs to $H^1_{\bal}(\bb{Q},\bb{V}_{f\hg\hh}^\dagger\vert_{\cl{C}})$, the same is true for $\widehat{\BF}(f,\hg,\hh)$.
\end{proof}

\begin{remark}
    As in Remark~\ref{rk:torsionfreeness}, we can ensure that both $H^1_{\cl{G}\cup +}(\bb{Q}, \bb{V}_{f\hg\hh}^\dagger)$ and $H^1(\bb{Q},\bb{V}_{f\hg\hh}^\dagger\vert_{\cl{C}})$ are torsion-free by imposing the condition that $H^0(\bb{Q},\overline{\rho}^\dagger)=0$, where $\overline{\rho}^\dagger$ is the residual representation attached to $\bb{V}_{f\hg \hh}^\dagger$.
\end{remark}


Let $\widehat{\kappa}(f,g_\alpha,h_\alpha)\in H^1_{\bal}(\bb{Q},V_{fgh}^\dagger)$ and $\widehat{\BF}(f,g_\alpha,h_\alpha)\in H^1_{\cl{G}\cup +}(\bb{Q},V_{fgh}^\dagger)$ be the specializations of the classes $\widehat{\kappa}(f,\hg,\hh)\in H^1_{\bal}(\bb{Q},\bb{V}_{f\hg\hh}^\dagger\vert_{\cl{C}})$ and $\widehat{\BF}(f,\hg,\hh)\in H^1(\bb{Q},\bb{V}_{f\hg\hh}^\dagger\vert_{\cl{C}})$, respectively, at the point $(y_0,z_0)$.

\begin{corollary}
    Assume that $H^1(\bb{Q},\bb{V}_{f\hg\hh}^\dagger\vert_{\cl{C}})$ is a torsion-free $\cl{O}_{\cl{C}}$-module and $H^1_{\cl{G}\cup +}(\bb{Q},\bb{V}_{f\hg\hh}^\dagger)$ is a torsion-free $\Lambda_{\hg\hh}$-module of rank $1$. Then $\widehat{\BF}(f,g_\alpha,h_\alpha)$ belongs to $H^1_{\bal}(\bb{Q}, V_{fgh}^\dagger)$ and
    \[
    \lambda_{N_g}(g)\cdot\widehat{\BF}(f,g_\alpha,h_\alpha)=\Omega_{f,\gamma}\cdot \widehat{\Lp}^g(f,g_\alpha,h_\alpha) \cdot \widehat{\kappa}(f,g_\alpha,h_\alpha).
    \]
\end{corollary}
\begin{proof}
    This is an immediate consequence of the previous proposition.
\end{proof}

%\begin{remark}
 %   Note that, in this exceptional case, if $H^1_{\bal}(\bb{Q}, V_{fgh}^\dagger)$ is $1$-dimensional, then $H^1_{\cl{G}\cup+}(\bb{Q}, V_{fgh}^\dagger)$ is $2$-dimensional. Therefore, the proof of Corollary~\ref{cor:non-exceptionalcomparison} cannot be adapted to this case, which leads us to impose different assumptions in the Corollary above.
%\end{remark}

\begin{comment}
The next theorem is stated for a fixed choice of the pair $(g,h)$ for the sake of simplification.

\begin{theorem}
Assume that $L(g,h,s)$ vanishes at $s=(l+m-1)/2$. Then, $\kappa_h(f,g,h)$ and $\kappa(g,h)$ belong to the Selmer group $\Sel(\mathbb Q, V(f,g,h))$ and $\Sel(\mathbb Q, V(g,h))$ respectively and \[ \log_{fgh}^{+-+}(\res_p(\kappa^{+-+}_h(f,g,h))) = \langle \log_{gh}^{-+}(\res_p^{-+}(\kappa(g,h))), \eta_g \otimes \omega_h \rangle.  \]
\end{theorem}

\begin{proof}
As before, we have the (weighted) one-variable element \[ \BF(f, g, h)\in H^1(\mathbb Q, \mathbb V(g, h) \otimes \chi), \] for any choice of character $\chi = 1, \varepsilon_K$, corresponding to the specialization of the three modular forms.

Proceeding as in \cite[Thm. 5.1]{BDV}, the module $H^1(\mathbb Q, \mathbb V(g, h))$ is torsion-free, and the result follows by applying the previous results to each good classical specialization of $h$.
\end{proof}
\end{comment}



\bibliographystyle{amsalpha}
\bibliography{AOR-refs}


\end{document} 
